% Lesson 42 — Vocabulary: Words and expressions related to training activities/courses on democracy — Anglais Tle A — Cours signé OnBuch+
% Programme officiel MINESEC. Compiler avec : tectonic EN42-vocabulary-words-and-expressions-related-to-training-ac.tex
\newcommand{\DOCMATIERE}{Anglais}
\newcommand{\DOCNIVEAU}{Tle A}
\newcommand{\DOCMODULE}{Chapter 4 : Citizenship and Human Rights}
\newcommand{\DOCLECON}{EN42}
\newcommand{\DOCTITRE}{Lesson 42 — Vocabulary: Words and expressions related to training activities/courses on democracy}
% OnBuch+ — préambule « cours signés » ENRICHI (compilé avec tectonic / XeLaTeX)
% Système visuel « Pop » d'OnBuch+ : fond crème (jamais blanc), bordures encre
% épaisses, ombres « dures » décalées, titres Archivo Black en capitales.
% Version MATHÉMATIQUES : graphes (pgfplots/tikz), boîtes pédagogiques
% (définition, propriété, méthode, attention, le savais-tu, exercice résolu,
% exercice, TP, bilan), page de garde et signature OnBuch+ ancrée en bas.
%
% Macros attendues dans main.tex AVANT \input{preamble} :
%   \DOCMATIERE  \DOCNIVEAU  \DOCMODULE  \DOCLECON (numéro)  \DOCTITRE
\documentclass[11pt]{article}

\usepackage{fontspec}
\usepackage[english,french]{babel} % français = langue principale ; l'anglais via \eng{..} / textebox
\usepackage[a4paper,margin=2cm,top=3.4cm,bottom=2.8cm,headheight=1.4cm,footskip=1.3cm]{geometry}
\usepackage{amsmath,amssymb,amsfonts}
\usepackage{mathtools} % \xmapsto, \coloneqq... souvent utilisés par les modèles
\usepackage{cancel} % simplifications barrées (bilans de forces)
% \abs{z} (module, valeur absolue) et \norm{u} (norme), utilisés par les modèles
\providecommand{\abs}{}\let\abs\relax\DeclarePairedDelimiter{\abs}{\lvert}{\rvert}
\providecommand{\norm}{}\let\norm\relax\DeclarePairedDelimiter{\norm}{\lVert}{\rVert}
\usepackage[table]{xcolor} % [table] : fournit aussi \rowcolors (alternance de lignes)
\usepackage{pagecolor}
\usepackage{tikz}
\usetikzlibrary{arrows.meta,positioning,calc,shapes.geometric,shapes.misc,shapes.arrows,decorations.pathmorphing,decorations.pathreplacing,decorations.markings,patterns,shadows,mindmap,trees,fit,backgrounds,matrix,intersections,angles,quotes}
\usepackage{pgfplots}
\usepackage[version=4]{mhchem} % équations nucléaires \ce{^{235}_{92}U}
\usepackage[european,siunitx]{circuitikz} % schémas électriques
\pgfplotsset{compat=1.18}
\usepgfplotslibrary{fillbetween,groupplots} % aires entre courbes (calcul intégral)
\usepackage{enumitem}
\usepackage{fancyhdr}
% [most] charge toutes les bibliothèques tcolorbox (listings, minted, raster,
% documentation...) et gonfle inutilement la mémoire TeX sur un document long
% (~300 boîtes) : on ne charge que ce qui est réellement utilisé.
\usepackage[skins,breakable]{tcolorbox}
\usepackage{graphicx}
\usepackage{colortbl}
\usepackage{booktabs}
\usepackage{tabularx}
\usepackage{diagbox} % case d'en-tête diagonale des tableaux à double entrée
% Colonnes extensibles centrée / gauche / droite (C, L, R), souvent utilisées par les modèles
\newcolumntype{C}{>{\centering\arraybackslash}X}
\newcolumntype{L}{>{\raggedright\arraybackslash}X}
\newcolumntype{R}{>{\raggedleft\arraybackslash}X}
\usepackage{array}
\usepackage{multicol}
\usepackage{multirow}
\usepackage{siunitx}
% parse-numbers=false : accepte aussi \SI{2,5 \times 10^{-3}}{...}
\sisetup{locale=FR,output-decimal-marker={,},group-minimum-digits=4,parse-numbers=false}
\DeclareSIUnit\ohm{\ensuremath{\symup{\Omega}}} % U+2126 absent de la police
\DeclareSIUnit\division{div} % réglages d'oscilloscope (ms/div, V/div)
\DeclareSIUnit\tr{tr} % tours (vitesse de rotation, ex. \SI{1500}{\tr\per\minute})
\DeclareSIUnit\molar{M} % concentration molaire (ex. \SI{3}{\molar}, \SI{1}{\micro\molar})
\DeclareSIUnit\an{an} % année (le modèle confond parfois avec \year, un compteur TeX) — ex. \SI{5}{\kilo\meter\per\an}
\DeclareSIUnit\FCFA{FCFA} % franc CFA (ex. \SI{500}{\FCFA\per\kilo\gram})
\DeclareSIUnit\degreeBrix{\textdegree Brix} % teneur en sucre d'un sirop/jus (réfractométrie)
\DeclareSIUnit\month{mois} % le modèle utilise parfois \month, un compteur TeX, comme unité de durée
\usepackage{qrcode}
% \degree : commande fréquemment utilisée par les modèles pour un angle ou une
% température en dehors de \SI{}{} (non fournie par les paquets chargés).
\providecommand{\degree}{^{\circ}}
% \qed (fin de démonstration), utilisé par les modèles sans amsthm
\providecommand{\qed}{\hfill$\square$}
% \barre : double barre des valeurs interdites dans un tableau de variations (tkz-tab)
\providecommand{\barre}{\ensuremath{\|}}
% environnement proof (amsthm non chargé) utilisé par les modèles
\ifdefined\proof\else\newenvironment{proof}[1][Démonstration]{\par\noindent\textit{#1.}\ }{\qed\par}\fi
\usepackage{lastpage}
\usepackage{hyperref}
\hypersetup{hidelinks}

% ============================================================
% Palette « Pop » OnBuch+ — mêmes hex que la classe P (plus_ui.dart)
% ============================================================
\definecolor{poporange}{HTML}{F59321}
\definecolor{popdark}{HTML}{E07A0C}
\definecolor{popcream}{HTML}{FFF6E8}
\definecolor{popink}{HTML}{1C1714}
\definecolor{poppurple}{HTML}{7A5AE0}
\definecolor{popgreen}{HTML}{2E9E6A}
\definecolor{poppink}{HTML}{E0457B}
\definecolor{popblue}{HTML}{2F7DE1}
\definecolor{popgold}{HTML}{C98A12}
\definecolor{popmuted}{HTML}{8A7F76}
\definecolor{popline}{HTML}{EADFD2}
\definecolor{popgray}{HTML}{8A7F76}
\colorlet{popgrey}{popgray}
% Alias tolérants (le modèle nomme parfois les couleurs différemment)
\colorlet{popwhite}{white}
\colorlet{popblack}{popink}
\colorlet{popred}{poppink}
\colorlet{popyellow}{popgold}
% Teintes claires (fonds de tableaux/encadrés) — le modèle nomme parfois
% les couleurs avec un suffixe L (ex. popblueL) sans les avoir définies.
\colorlet{poporangeL}{poporange!20}
\colorlet{popdarkL}{popdark!20}
\colorlet{poppurpleL}{poppurple!20}
\colorlet{popgreenL}{popgreen!20}
\colorlet{poppinkL}{poppink!20}
\colorlet{popblueL}{popblue!20}
\colorlet{popgoldL}{popgold!20}
\colorlet{popredL}{popred!20}
\colorlet{popgrayL}{popgray!20}
% Teintes claires pour fonds de tableaux / zones
\colorlet{popblueL}{popblue!10!white}
\colorlet{poporangeL}{poporange!14!white}
\colorlet{popgreenL}{popgreen!12!white}
\colorlet{poppurpleL}{poppurple!12!white}
\colorlet{poppinkL}{poppink!12!white}
\colorlet{popgoldL}{popgold!14!white}

\pagecolor{popcream}

% ============================================================
% Typographie
% ============================================================
\defaultfontfeatures{Ligatures=TeX}
\setmainfont{PlusJakartaSans-Medium.ttf}[Path=../../../fonts/,BoldFont=PlusJakartaSans-Bold.ttf]
% Symboles, API et emojis absents de la police principale : repli sur DejaVu Sans (caractères actifs)
\newfontface\symfont{DejaVuSans.ttf}[Path=../../../fonts/]
\makeatletter
\newcommand\symactive[1]{\catcode#1=13 \begingroup\lccode`\~=#1 \lowercase{\endgroup\def~}{{\symfont\char#1}}}
\newcommand\symblank[1]{\catcode#1=13 \begingroup\lccode`\~=#1 \lowercase{\endgroup\def~}{}}
\newcommand\symhyph[1]{\catcode#1=13 \begingroup\lccode`\~=#1 \lowercase{\endgroup\def~}{\mbox{-}}}
\newcount\symc
\newcommand\symrange[2]{\symc=#1 \@whilenum\symc<#2 \do{\expandafter\symactive\expandafter{\the\symc}\advance\symc by 1 }}
\newcommand\symblankrange[2]{\symc=#1 \@whilenum\symc<#2 \do{\expandafter\symblank\expandafter{\the\symc}\advance\symc by 1 }}
\makeatother
\symrange{"0250}{"02FF}\symactive{"2713}\symactive{"2714}\symactive{"2717}\symactive{"2718}\symactive{"2610}\symactive{"2611}\symactive{"2612}
\symhyph{"2011}\symhyph{"2E1A}\symblankrange{"1F300}{"1FAFF}\symblank{"FE0F}
% Maths en sans-serif (Fira Math) avec chiffres et lettres droites de Plus
% Jakarta, pour que les formules se fondent dans le texte.
\usepackage[math-style=ISO,bold-style=ISO]{unicode-math}
\setmathfont{FiraMath-Regular.otf}[Path=../../../fonts/]
\setmathfont{PlusJakartaSans-Medium.ttf}[Path=../../../fonts/,range={up/num,up/{Latin,latin},\mathup}]
\usepackage{icomma} % virgule décimale sans espace en mode math
% Caractères mathématiques Unicode absents des polices de texte (Plus Jakarta,
% Archivo Black) : rendus par leur équivalent mathématique, en texte comme en
% formule — sinon ils s'affichent en carré vide (ex. « Arithmétique dans ℤ »).
\usepackage{newunicodechar}
\newunicodechar{ℝ}{\ensuremath{\mathbb{R}}}
\newunicodechar{ℤ}{\ensuremath{\mathbb{Z}}}
\newunicodechar{ℂ}{\ensuremath{\mathbb{C}}}
\newunicodechar{ℕ}{\ensuremath{\mathbb{N}}}
\newunicodechar{ℚ}{\ensuremath{\mathbb{Q}}}
\newunicodechar{𝔻}{\ensuremath{\mathbb{D}}}
\newunicodechar{α}{\ensuremath{\alpha}}
\newunicodechar{β}{\ensuremath{\beta}}
\newunicodechar{γ}{\ensuremath{\gamma}}
\newunicodechar{δ}{\ensuremath{\delta}}
\newunicodechar{ε}{\ensuremath{\varepsilon}}
\newunicodechar{θ}{\ensuremath{\theta}}
\newunicodechar{λ}{\ensuremath{\lambda}}
\newunicodechar{μ}{\ensuremath{\mu}}
\newunicodechar{σ}{\ensuremath{\sigma}}
\newunicodechar{τ}{\ensuremath{\tau}}
\newunicodechar{φ}{\ensuremath{\varphi}}
\newunicodechar{ψ}{\ensuremath{\psi}}
\newunicodechar{ω}{\ensuremath{\omega}}
\newunicodechar{Σ}{\ensuremath{\Sigma}}
% majuscules produites par \MakeUppercase dans les titres (α-aminés -> Α)
\newunicodechar{Α}{\ensuremath{\alpha}}
\newunicodechar{Β}{\ensuremath{\beta}}
\newunicodechar{Γ}{\ensuremath{\Gamma}}
\newunicodechar{Δ}{\ensuremath{\Delta}}
\newunicodechar{Ε}{\ensuremath{\varepsilon}}
\newunicodechar{Θ}{\ensuremath{\Theta}}
\newunicodechar{Λ}{\ensuremath{\Lambda}}
\newunicodechar{Μ}{\ensuremath{\mu}}
\newunicodechar{Τ}{\ensuremath{\tau}}
\newunicodechar{Φ}{\ensuremath{\Phi}}
\newunicodechar{Ψ}{\ensuremath{\Psi}}
\newunicodechar{Ω}{\ensuremath{\Omega}}
\newunicodechar{↦}{\ensuremath{\mapsto}}
\newunicodechar{∈}{\ensuremath{\in}}
\newunicodechar{∉}{\ensuremath{\notin}}
\newunicodechar{∧}{\ensuremath{\wedge}}
\newunicodechar{⇔}{\ensuremath{\Leftrightarrow}}
\newunicodechar{⟺}{\ensuremath{\Longleftrightarrow}}
\newunicodechar{⇒}{\ensuremath{\Rightarrow}}
\newunicodechar{⟹}{\ensuremath{\Longrightarrow}}
\newunicodechar{∘}{\ensuremath{\circ}}
\newunicodechar{∪}{\ensuremath{\cup}}
\newunicodechar{∩}{\ensuremath{\cap}}
\newunicodechar{≡}{\ensuremath{\equiv}}
\newunicodechar{⊂}{\ensuremath{\subset}}
\newunicodechar{⊥}{\ensuremath{\perp}}
\newunicodechar{∃}{\ensuremath{\exists}}
\newunicodechar{∀}{\ensuremath{\forall}}
\newunicodechar{∛}{\ensuremath{\sqrt[3]{\,}}}
\newunicodechar{∜}{\ensuremath{\sqrt[4]{\,}}}
\newunicodechar{⃗}{\ensuremath{{}^{\to}}}
\newunicodechar{ⁿ}{\ifmmode^{n}\else\textsuperscript{\ensuremath{n}}\fi}
\newunicodechar{ˣ}{\ifmmode^{x}\else\textsuperscript{\ensuremath{x}}\fi}
\newunicodechar{ᵇ}{\ifmmode^{b}\else\textsuperscript{\ensuremath{b}}\fi}
\newunicodechar{ʳ}{\ifmmode^{r}\else\textsuperscript{\ensuremath{r}}\fi}
\newunicodechar{ᵘ}{\ifmmode^{u}\else\textsuperscript{\ensuremath{u}}\fi}
\newunicodechar{ʸ}{\ifmmode^{y}\else\textsuperscript{\ensuremath{y}}\fi}
\newunicodechar{ᵃ}{\ifmmode^{a}\else\textsuperscript{\ensuremath{a}}\fi}
\newunicodechar{ᵉ}{\ifmmode^{e}\else\textsuperscript{\ensuremath{e}}\fi}
\newunicodechar{ᵏ}{\ifmmode^{k}\else\textsuperscript{\ensuremath{k}}\fi}
\newunicodechar{ᵖ}{\ifmmode^{p}\else\textsuperscript{\ensuremath{p}}\fi}
\newunicodechar{ᴺ}{\ifmmode^{N}\else\textsuperscript{\ensuremath{N}}\fi}
\newunicodechar{ᶜ}{\ifmmode^{c}\else\textsuperscript{\ensuremath{c}}\fi}
\newunicodechar{ʰ}{\ifmmode^{h}\else\textsuperscript{\ensuremath{h}}\fi}
\newunicodechar{ᵠ}{\ifmmode^{q}\else\textsuperscript{\ensuremath{q}}\fi}
\newunicodechar{ₙ}{\ifmmode_{n}\else\textsubscript{\ensuremath{n}}\fi}
\newunicodechar{ₐ}{\ifmmode_{a}\else\textsubscript{\ensuremath{a}}\fi}
\newunicodechar{ᵢ}{\ifmmode_{i}\else\textsubscript{\ensuremath{i}}\fi}
\newunicodechar{ᵣ}{\ifmmode_{r}\else\textsubscript{\ensuremath{r}}\fi}
\newunicodechar{ₖ}{\ifmmode_{k}\else\textsubscript{\ensuremath{k}}\fi}
\newunicodechar{ᵦ}{\ifmmode_{\beta}\else\textsubscript{\ensuremath{\beta}}\fi}
\newunicodechar{ₓ}{\ifmmode_{x}\else\textsubscript{\ensuremath{x}}\fi}
\newfontfamily\popdisplay{ArchivoBlack-Regular.ttf}[Path=../../../fonts/]
\newfontfamily\poplogo{SpaceGrotesk-Bold.ttf}[Path=../../../fonts/]
\newfontfamily\popsemi{PlusJakartaSans-SemiBold.ttf}[Path=../../../fonts/]
\newfontfamily\popxbold{PlusJakartaSans-ExtraBold.ttf}[Path=../../../fonts/]
\newfontfamily\popxboldls{PlusJakartaSans-ExtraBold.ttf}[Path=../../../fonts/,LetterSpace=3.0]

\setlist[enumerate]{leftmargin=1.7em,itemsep=5pt,topsep=4pt}
\setlist[itemize]{leftmargin=1.7em,itemsep=4pt,topsep=4pt,label={\color{poporange}\textbullet}}
\setlength{\parindent}{0pt}
\setlength{\parskip}{5pt}
\renewcommand{\arraystretch}{1.25}

% pgfplots : style Pop par défaut
\pgfplotsset{
  every axis/.append style={axis line style={popink,line width=1pt},tick style={popink},
    grid style={popline,line width=0.5pt},label style={font=\small},tick label style={font=\footnotesize}},
  popcurve/.style={poporange,line width=1.8pt,no markers,smooth},
}
% Même style disponible dans un \draw TikZ simple (hors pgfplots)
\tikzset{popcurve/.style={poporange,line width=1.8pt}}
\tikzset{popgrid/.style={popline,very thin}}
\colorlet{popcurve}{poporange} % le nom du style est parfois utilisé comme couleur

% ============================================================
% Pill / squiggle
% ============================================================
\newcommand{\poppill}[3]{%
  \tikz[baseline=-0.55ex]\node[draw=popink,line width=1.2pt,rounded corners=2.6pt,%
    fill=#2,inner xsep=6.5pt,inner ysep=3pt]{%
    \popxboldls\fontsize{7}{7}\selectfont\color{#3}\MakeUppercase{#1}};%
}
\newcommand{\popsquiggle}[1]{%
  \tikz[baseline=-0.4ex]{
    \draw[#1,line width=2.6pt,line cap=round]
      plot [smooth] coordinates {(0,0.09) (0.34,0.01) (0.68,0.21) (1.02,0.01) (1.36,0.10)};
  }%
}
% Mot-clé surligné (marqueur orange sous le texte)
\newcommand{\cle}[1]{{\popxbold\color{popdark}#1}}

% ============================================================
% En-tête / pied
% ============================================================
\pagestyle{fancy}
\fancyhf{}
\renewcommand{\headrulewidth}{0pt}
\renewcommand{\footrulewidth}{0pt}
\lhead{\raisebox{-0.15ex}{\poplogo\fontsize{13}{13}\selectfont\color{popink}OnBuch\color{poporange}+}}
\rhead{\poppill{\DOCMATIERE\ \textbullet\ \DOCNIVEAU\ \textbullet\ Leçon \DOCLECON}{white}{popink}}
\lfoot{\popxbold\fontsize{7.6}{7.6}\selectfont\color{popmuted}\MakeUppercase{Cours signé OnBuch+}\ \textbullet\ {\popsemi\DOCTITRE}}
\rfoot{\tikz[baseline=-0.6ex]\node[rounded corners=8.5pt,fill=popink,minimum height=17pt,minimum width=17pt,inner xsep=5pt,inner sep=0]{\popxbold\fontsize{8}{8}\selectfont\color{popcream}\thepage\,/\,\pageref*{LastPage}};}

% ============================================================
% Titres
% ============================================================
\newcommand{\chaptitre}[2]{%
  \begin{tcolorbox}[enhanced,colback=poporange,colframe=popink,boxrule=2.6pt,arc=11pt,
      drop shadow=popink,left=15pt,right=15pt,top=12pt,bottom=11pt,before skip=3pt,after skip=18pt]
    \poppill{#2}{popink}{popcream}\par\vspace{9pt}
    {\raggedright\hyphenpenalty=10000\exhyphenpenalty=10000\popdisplay\fontsize{22}{24}\selectfont\color{popcream}\MakeUppercase{#1}\par}
  \end{tcolorbox}
}
% Section numérotée : pastille orange avec le numéro + titre Archivo
\newcounter{popsec}
\newcommand{\coursec}[1]{%
  \stepcounter{popsec}\setcounter{popsub}{0}%
  \par\vspace{16pt}\noindent
  \begin{minipage}{\linewidth}
  \tikz[baseline=-0.7ex]\node[draw=popink,line width=1.6pt,fill=poporange,rounded corners=4pt,minimum size=22pt,inner sep=0]{\popdisplay\fontsize{12}{12}\selectfont\color{popcream}\thepopsec};%
  \hspace{8pt}{\popdisplay\fontsize{15}{17}\selectfont\color{popink}\MakeUppercase{#1}}\par
  \vspace{4pt}\noindent{\color{poporange}\rule{36pt}{3.6pt}}
  \end{minipage}\par\nopagebreak\vspace{8pt}
}
\newcounter{popsub}
\newcommand{\courssub}[1]{%
  \stepcounter{popsub}%
  \par\vspace{9pt}\noindent{\popxbold\fontsize{11.5}{13}\selectfont\color{popink}{\color{poporange}\thepopsec.\thepopsub}\hspace{6pt}#1}\par\nopagebreak\vspace{4pt}
}

% ============================================================
% Boîtes pédagogiques — même recette : fond blanc, bordure épaisse colorée,
% ombre dure encre, badge plein. Titre optionnel [#1].
% ============================================================
\tcbset{popbase/.style={breakable,enhanced,colback=white,boxrule=2.3pt,arc=9pt,
  shadow={3pt}{-3pt}{0pt}{popink},left=14pt,right=14pt,top=11pt,bottom=11pt,before skip=11pt,after skip=15pt,
  fontupper=\popsemi\fontsize{10.6}{14.5}\selectfont\color{popink},
  fontlower=\popsemi\fontsize{10.6}{14.5}\selectfont\color{popink}}}
% Coche dessinée (le glyphe ✓ n'existe pas dans les polices du document)
\newcommand{\popcheck}[1]{\tikz[baseline=-0.5ex]\draw[#1,line width=1.6pt,line cap=round,line join=round](-0.13,0.0)--(-0.04,-0.09)--(0.13,0.1);}
\providecommand{\coche}{\popcheck{popgreen}}
\providecommand{\resultat}[1]{\par\smallskip\noindent\fcolorbox{popgreen}{popgreenL}{\parbox{\dimexpr\linewidth-2\fboxsep-2\fboxrule\relax}{\textbf{Résultat.} #1}}\par\smallskip}
\newcommand{\popboxhead}[3]{\poppill{#1}{#2}{white}\ifx\relax#3\relax\else\hspace{6pt}{\popxbold\fontsize{10.6}{12}\selectfont\color{popink}#3}\fi\par\vspace{7pt}}

\newtcolorbox{aretenir}[1][]{popbase,colframe=popgreen,before upper={\popboxhead{À retenir}{popgreen}{#1}}}
% Texte en anglais (document de lecture, dialogue, extrait) : cadre
% d'étude. Un titre optionnel (source, auteur) s'affiche dans l'en-tête.
\newtcolorbox{textebox}[1][]{popbase,colframe=popblue,before upper={\popboxhead{Text}{popblue}{#1}\begin{otherlanguage}{english}},after upper={\end{otherlanguage}}}
% Marques ✓ / ✗ (forme correcte / fautive) souvent utilisées par les modèles
\providecommand{\cross}{\textcolor{poppink}{\ensuremath{\times}}}
\providecommand{\xmark}{\cross}\providecommand{\crossmark}{\cross}\providecommand{\cmark}{\ensuremath{\checkmark}}
% Ligne de réponse / justification à compléter (inventée par les modèles)
\providecommand{\justification}[1]{\par\noindent\textit{Justification~:} \dotfill\par}
% \textipa (package tipa non chargé) : la police ne couvre pas l'API -> simple passage
\providecommand{\textipa}[1]{#1}
% Soulignement (ulem non chargé) : \uline, \ul
\providecommand{\uline}[1]{\underline{#1}}\providecommand{\ul}[1]{\underline{#1}}
% \glossaire{\item ...} inventée par les modèles : liste de glossaire
\providecommand{\glossaire}[1]{\begin{itemize}#1\end{itemize}}
% \alert (beamer) inventée par les modèles : texte mis en évidence
\providecommand{\alert}[1]{\textbf{\textcolor{poppink}{#1}}}
% variantes ASCII d'ordinaux écrites en macro
\providecommand{\iere}{\textsuperscript{re}}\providecommand{\ieme}{\textsuperscript{e}}\providecommand{\ere}{\textsuperscript{re}}\providecommand{\eme}{\textsuperscript{e}}\providecommand{\er}{\textsuperscript{er}}\providecommand{\ie}{\textsuperscript{e}}
% Ordinaux français écrits en macro par les modèles : 1\ière, 3\ième
\newcommand{\ière}{\textsuperscript{re}}\newcommand{\ième}{\textsuperscript{e}}
% \exemple (inventée par les modèles) : étiquette « Exemple : »
\providecommand{\exemple}{\textbf{Exemple~:}\ }
% \square utilisable aussi hors mode math (cases à cocher)
\AtBeginDocument{\let\squareMath\square\renewcommand{\square}{\ensuremath{\squareMath}}}
% Mot ou phrase en anglais à l'intérieur d'un paragraphe français
\newcommand{\eng}[1]{\ifmmode\text{\textit{#1}}\else\begingroup\selectlanguage{english}\itshape #1\endgroup\fi}
\let\esp\eng % alias toléré
% flèches d'intonation inventées par les modèles
\providecommand{\textsearrow}{}\renewcommand{\textsearrow}{\ensuremath{\searrow}}\providecommand{\textnearrow}{}\renewcommand{\textnearrow}{\ensuremath{\nearrow}}
% \fr{..} : mot français (inventé par les modèles)
\providecommand{\fr}[1]{\textit{#1}}
\newtcolorbox{exemplebox}[1][]{popbase,colframe=poporange,before upper={\popboxhead{Exemple}{poporange}{#1}}}
\newtcolorbox{definition}[1][]{popbase,colframe=popblue,before upper={\popboxhead{Définition}{popblue}{#1}}}
\newtcolorbox{propriete}[1][]{popbase,colframe=poppurple,before upper={\popboxhead{Propriété}{poppurple}{#1}}}
\newtcolorbox{methode}[1][]{popbase,colframe=popgold,before upper={\popboxhead{Méthode}{popgold}{#1}}}
\newtcolorbox{attention}[1][]{popbase,colframe=poppink,before upper={\popboxhead{Attention}{poppink}{#1}}}
\newtcolorbox{experience}[1][]{popbase,colframe=popblue,colback=popblueL,before upper={\popboxhead{Expérience}{popblue}{#1}}}
% « Le savais-tu ? » : carte violette pleine (contexte, culture, Cameroun)
\newtcolorbox{savaistu}[1][]{popbase,colback=poppurple,colframe=popink,
  fontupper=\popsemi\fontsize{10.4}{14.2}\selectfont\color{white},
  before upper={\poppill{Le savais-tu\,?}{white}{poppurple}\ifx\relax#1\relax\else\hspace{6pt}{\popxbold\color{white}#1}\fi\par\vspace{7pt}}}
% Exercice résolu : énoncé en haut, \tcblower puis la solution sur fond crème
\newcounter{popexo}\newcounter{popexr}
\newtcolorbox{exoresolu}[1][]{popbase,colframe=poporange,colbacklower=popcream,
  segmentation style={popink,line width=1.2pt,dashed},
  before upper={\stepcounter{popexr}\popboxhead{Exercice résolu \thepopexr}{poporange}{#1}},
  before lower={\poppill{Solution}{popink}{popcream}\par\vspace{6pt}}}
% Exercice (non résolu en place) — la correction est dans la partie Corrigés
\newtcolorbox{exercice}[1][]{popbase,colframe=popink,
  before upper={\stepcounter{popexo}\popboxhead{Exercice \thepopexo}{popink}{#1}}}
% Corrigé
\newtcolorbox{corrige}[1][]{popbase,colframe=popgreen,colback=popgreenL,
  before upper={\popboxhead{Corrigé}{popgreen}{#1}}}
% Objectifs / prérequis / bilan
\newtcolorbox{objectifs}{popbase,colframe=popink,colback=poporangeL,
  before upper={\popboxhead{Objectifs de la leçon}{popink}{}}}
\newtcolorbox{prerequis}{popbase,colframe=popink,colback=popblueL,
  before upper={\popboxhead{Prérequis}{popblue}{}}}
\newtcolorbox{bilan}{popbase,colframe=popink,colback=white,boxrule=2.8pt,
  before upper={\popboxhead{Fiche bilan}{poporange}{L'essentiel à connaître pour l'examen}}}
% Figure encadrée avec légende
\newtcolorbox{popfigure}[1][]{enhanced,breakable=false,colback=white,colframe=popline,boxrule=1.4pt,arc=8pt,
  left=10pt,right=10pt,top=10pt,bottom=8pt,before skip=10pt,after skip=12pt,halign=center,
  lower separated=false,halign lower=center,
  fontlower=\popsemi\fontsize{9}{11.5}\selectfont\color{popmuted},
  #1}
\newcommand{\legende}[1]{\par\vspace{6pt}{\popsemi\fontsize{9}{11.5}\selectfont\color{popmuted}#1}}

% ============================================================
% Signature OnBuch+
% ============================================================
\newcommand{\poplogoinline}[1]{{\poplogo\fontsize{#1}{#1}\selectfont\color{popink}OnBuch\color{poporange}+}}
\newcommand{\signatureonbuch}{%
  \begin{tcolorbox}[enhanced,colback=popink,colframe=popink,boxrule=2.2pt,arc=10pt,
      drop shadow=poporange,left=14pt,right=14pt,top=10pt,bottom=10pt,before skip=0pt,after skip=0pt]
    \tikz[baseline=-0.7ex]\node[circle,fill=popgreen,draw=popcream,line width=1.3pt,minimum size=22pt,inner sep=0]{\popcheck{white}};%
    \hspace{9pt}\begin{minipage}[c]{0.62\linewidth}
      {\popxbold\fontsize{12}{13}\selectfont\color{popcream}Signé\ \textbullet\ {\poplogo OnBuch\color{poporange}+}}\par\vspace{2pt}
      {\raggedright\popsemi\fontsize{8}{10.5}\selectfont\color{popline}\MakeUppercase{Équipe pédagogique OnBuch+}\\\MakeUppercase{Conforme au programme officiel MINESEC}\par}
    \end{minipage}\hfill
    {\poplogo\fontsize{22}{22}\selectfont\color{popcream}OnBuch\color{poporange}+}
  \end{tcolorbox}%
}

% ============================================================
% Page de garde
% #1 titre · #2 matière · #3 niveau · #4 module · #5 n° leçon · #6 durée ·
% #7 description / objectifs (texte ou itemize)
% ============================================================
\newcommand{\pagegarde}[7]{%
  \thispagestyle{empty}
  \newgeometry{a4paper,margin=1.7cm,top=1.6cm,bottom=1.4cm}%
  \begin{tikzpicture}[remember picture,overlay]
    \fill[poporange!18] ([xshift=-3.2cm,yshift=-3.6cm]current page.north east) circle (4.6cm);
    \fill[poppurple!14] ([xshift=2.4cm,yshift=7.6cm]current page.south west) circle (3.8cm);
  \end{tikzpicture}%
  {\poplogo\fontsize{30}{30}\selectfont\color{popink}OnBuch\color{poporange}+}\hfill
  \raisebox{0.6ex}{\poppill{Cours signé \textbullet\ Premium}{popink}{popcream}}\par
  \vspace{1pt}\popsquiggle{poppurple}\par
  \vspace{8pt}
  \begin{tcolorbox}[enhanced,colback=poporange,colframe=popink,boxrule=2.8pt,arc=13pt,
      drop shadow=popink,left=18pt,right=18pt,top=12pt,bottom=13pt,after skip=10pt]
    \poppill{#2 \textbullet\ #3}{popink}{popcream}\hspace{5pt}\poppill{#4}{white}{popink}\par\vspace{8pt}
    {\popdisplay\fontsize{14}{14}\selectfont\color{popink}LEÇON #5}\par\vspace{4pt}
    {\raggedright\hyphenpenalty=10000\exhyphenpenalty=10000\popdisplay\fontsize{25}{27}\selectfont\color{popcream}\MakeUppercase{#1}\par}
    \vspace{8pt}
    {\popxbold\fontsize{9.5}{9.5}\selectfont\color{popink}\MakeUppercase{Durée conseillée : #6}}
  \end{tcolorbox}
  \begin{tcolorbox}[enhanced,colback=white,colframe=popink,boxrule=2.2pt,arc=10pt,
      drop shadow=popink,left=16pt,right=16pt,top=10pt,bottom=10pt,after skip=8pt]
    \poppill{Dans cette leçon}{poppurple}{white}\par\vspace{5pt}
    \popsemi\fontsize{9.3}{12.6}\selectfont\color{popink}#7
  \end{tcolorbox}
  \noindent\begin{minipage}[t]{0.487\linewidth}
    \begin{tcolorbox}[enhanced,colback=popgreen,colframe=popink,boxrule=2.2pt,arc=10pt,
        drop shadow=popink,left=11pt,right=11pt,top=10pt,bottom=10pt,height=3.1cm,valign=center]
      \tcbox[colback=white,colframe=popink,boxrule=1.3pt,arc=5pt,boxsep=0pt,nobeforeafter,
        left=3pt,right=3pt,top=3pt,bottom=3pt,valign=center]{\qrcode[height=1.9cm]{https://play.google.com/store/apps/details?id=com.luvvix.onbuch&hl=fr}}%
      \hspace{8pt}\begin{minipage}[c]{0.5\linewidth}
        {\popdisplay\fontsize{11}{12.5}\selectfont\color{white}\MakeUppercase{Télécharge OnBuch}}\par\vspace{4pt}
        {\raggedright\popsemi\fontsize{8.4}{11}\selectfont\color{white}Gratuit \textbullet\ Google Play\\Tuteur IA, annales, concours, résultats\par}
      \end{minipage}
    \end{tcolorbox}
  \end{minipage}\hfill
  \begin{minipage}[t]{0.487\linewidth}
    \begin{tcolorbox}[enhanced,colback=poppurple,colframe=popink,boxrule=2.2pt,arc=10pt,
        drop shadow=popink,left=11pt,right=11pt,top=10pt,bottom=10pt,height=3.1cm,valign=center]
      \tikz[baseline=-0.7ex]\node[circle,fill=white,draw=popink,line width=1.3pt,minimum size=1.6cm,inner sep=0]{\popdisplay\fontsize{24}{24}\selectfont\color{poppurple}+};%
      \hspace{8pt}\begin{minipage}[c]{0.6\linewidth}
        {\popdisplay\fontsize{11}{12.5}\selectfont\color{white}\MakeUppercase{Passe à OnBuch+}}\par\vspace{4pt}
        {\raggedright\popsemi\fontsize{8.4}{11}\selectfont\color{white}Tous les cours signés, Tuteur IA illimité, fiches PDF — onglet OnBuch+ de l'app\par}
      \end{minipage}
    \end{tcolorbox}
  \end{minipage}
  \vspace{8pt}
  \popcontenu
  \vfill
  \signatureonbuch
  \restoregeometry
  \newpage
}

% Rangée « Ce cours contient » (page de garde)
\newcommand{\poptile}[3]{% #1 couleur #2 gros texte #3 légende
  \begin{tcolorbox}[enhanced,colback=#1,colframe=popink,boxrule=2pt,arc=9pt,shadow={2.5pt}{-2.5pt}{0pt}{popink},
      left=6pt,right=6pt,top=9pt,bottom=9pt,halign=center,valign=center,height=1.75cm,width=\linewidth,nobeforeafter]
    {\popdisplay\fontsize{12.5}{13}\selectfont\color{white}\MakeUppercase{#2}}\par\vspace{3pt}
    {\centering\popsemi\fontsize{7.8}{9.5}\selectfont\color{white}#3\par}
  \end{tcolorbox}}
\newcommand{\popcontenu}{%
  \par\noindent{\popxboldls\fontsize{7.5}{7.5}\selectfont\color{popmuted}CE COURS SIGNÉ CONTIENT}\par\vspace{7pt}
  \noindent\begin{minipage}[t]{0.235\linewidth}\poptile{popblue}{Cours}{complet, illustré, pas à pas}\end{minipage}\hfill
  \begin{minipage}[t]{0.235\linewidth}\poptile{poporange}{Méthodes}{et exercices résolus}\end{minipage}\hfill
  \begin{minipage}[t]{0.235\linewidth}\poptile{poppink}{Exercices}{type examen + corrigés}\end{minipage}\hfill
  \begin{minipage}[t]{0.235\linewidth}\poptile{popgreen}{Fiche bilan}{l'essentiel à retenir}\end{minipage}\par}

% Sommaire
\newtcolorbox{sommaire}{enhanced,colback=white,colframe=popink,boxrule=2.2pt,arc=10pt,
  drop shadow=popink,left=18pt,right=18pt,top=13pt,bottom=13pt,before skip=6pt,after skip=20pt,
  before upper={\poppill{Sommaire}{poppurple}{white}\par\vspace{9pt}\popsemi\fontsize{10.6}{14.5}\selectfont\color{popink}}}

% Signature de fin de cours — poussée tout en bas de la dernière page
\newcommand{\findecours}{%
  \par\vfill\nopagebreak
  \begin{center}\popsquiggle{poporange}\end{center}\vspace{4pt}
  \signatureonbuch
}
\AtBeginDocument{\def\llbracket{\mathopen{[\mkern-3.5mu[}}\def\rrbracket{\mathclose{]\mkern-3.5mu]}}} % intervalles d'entiers (glyphe ⟦ absent de la police)
% Glyphes absents de Fira Math : redéfinis à partir de glyphes disponibles
\AtBeginDocument{%
  \def\setminus{\mathbin{\backslash}}\let\smallsetminus\setminus
  \def\vdots{\mathord{\vbox{\baselineskip3.2pt\lineskiplimit0pt\kern4pt\hbox{.}\hbox{.}\hbox{.}}}}%
  \def\ddots{\mathinner{\mkern1mu\raise6.5pt\hbox{.}\mkern2mu\raise3.6pt\hbox{.}\mkern2mu\raise0.7pt\hbox{.}\mkern1mu}}%
  \def\triangle{\mathord{\scalebox{0.9}{$\symup{\Delta}$}}}%
  \def\nabla{\mathord{\raisebox{\depth}{\scalebox{1}[-1]{$\symup{\Delta}$}}}}%
  \def\checkmark{\popcheck{popdark}}%
  \def\star{\mathbin{\ast}}\def\bigstar{\mathord{\ast}}%
  \def\diamond{\mathbin{\scalebox{0.6}{$\Diamond$}}}%
  \def\bigcup{\mathop{\mathchoice{\raisebox{-0.3em}{\scalebox{1.7}{$\cup$}}}{\scalebox{1.25}{$\cup$}}{\cup}{\cup}}\displaylimits}%
  \def\bigcap{\mathop{\mathchoice{\raisebox{-0.3em}{\scalebox{1.7}{$\cap$}}}{\scalebox{1.25}{$\cap$}}{\cap}{\cap}}\displaylimits}%
  \def\lesseqqgtr{\mathrel{\lessgtr}}\def\bigoplus{\mathop{\oplus}\displaylimits}%
}
\providecommand{\ssi}{si et seulement si} % abréviation fréquente
\AtBeginDocument{\providecommand{\wideparen}{\overparen}} % arc de cercle

%% FIGFIX-BEGIN v4 — correctifs globaux des figures (docs/onbuch-generation/tools/figfix)
\usepackage{adjustbox}\usepackage{etoolbox}
% toute figure TikZ/pgfplots plus large que la ligne est réduite à la largeur disponible
\BeforeBeginEnvironment{tikzpicture}{\begin{adjustbox}{max width=\linewidth}}
\AfterEndEnvironment{tikzpicture}{\end{adjustbox}}
% légendes pgfplots lisibles par-dessus les courbes
\pgfplotsset{every axis/.append style={legend style={fill=white,fill opacity=0.92,text opacity=1,draw=black!15,inner sep=3pt}}}
% étiquettes d'arêtes (syntaxe edge["..."]) avec fond blanc
\tikzset{every edge quotes/.append style={fill=white,inner sep=1.2pt}}
%% FIGFIX-END
\begin{document}

\pagegarde{Lesson 42 — Vocabulary: Words and expressions related to training activities/courses on democracy}{Anglais}{Tle A}{Chapter 4 : Citizenship and Human Rights}{EN42}{2 h}{Cette leçon de vocabulaire présente les mots et expressions liés aux formations et cours sur la démocratie (workshops, seminars, civic education...). Elle permet aux élèves de Terminale A d'identifier, de prononcer et de réemployer ce lexique à l'oral comme à l'écrit, en évitant les faux amis et les calques du français.
\vspace{4pt}
\begin{multicols}{2}\small
\begin{itemize}
\item À la fin de cette leçon, je sais identifier le vocabulaire des activités de formation sur la démocratie dans un texte ou un dialogue.
\item Je sais classer ce lexique en champs : types de formations, acteurs, démarches, thèmes démocratiques.
\item Je sais employer les collocations correctes : attend a workshop, run a training session, raise awareness, etc.
\item Je sais éviter les faux amis : to attend ≠ assister (aider), formation ≠ training, etc.
\item Je sais former des mots de la même famille : train / trainer / trainee / training.
\item Je sais prononcer correctement les mots clés (workshop, seminar, democracy, citizen...).
\end{itemize}\end{multicols}}

\begin{sommaire}
\begin{multicols}{2}
\begin{enumerate}
\item Types de formations et de cours
\item Les acteurs de la formation
\item Thèmes démocratiques abordés en formation
\item Faux amis et pièges des francophones
\item Familles de mots, prononciation et orthographe
\item Le lexique en contexte : dialogue et production guidée
\item Activité d'intégration
\item Méthodes et exercices résolus
\item Exercices
\item Corrigés des exercices
\item Fiche bilan
\end{enumerate}
\end{multicols}
\end{sommaire}

\begin{objectifs}
\begin{itemize}
\item À la fin de cette leçon, je sais identifier le vocabulaire des activités de formation sur la démocratie dans un texte ou un dialogue.
\item Je sais classer ce lexique en champs : types de formations, acteurs, démarches, thèmes démocratiques.
\item Je sais employer les collocations correctes : attend a workshop, run a training session, raise awareness, etc.
\item Je sais éviter les faux amis : to attend ≠ assister (aider), formation ≠ training, etc.
\item Je sais former des mots de la même famille : train / trainer / trainee / training.
\item Je sais prononcer correctement les mots clés (workshop, seminar, democracy, citizen...).
\item Je sais réemployer ce lexique dans des phrases, un dialogue et un court texte.
\item Je sais rédiger une courte lettre de motivation pour participer à une formation sur la citoyenneté.
\end{itemize}
\end{objectifs}

\begin{prerequis}
\begin{itemize}
\item Vocabulaire de base de la citoyenneté vu en Première (citizen, vote, rights, election).
\item Les temps simples (present simple, past simple, future) pour décrire des activités.
\item La structure d'une lettre formelle simple (formule d'appel, introduction, conclusion).
\item Les verbes suivis de l'infinitif ou du gérondif (want to attend, enjoy taking part).
\end{itemize}
\end{prerequis}

\newpage

\coursec{Types de formations et de cours}

\begin{textebox}[An invitation from the British Council]
Imagine that your school in Buea has just received an invitation from the British Council: a three-day workshop on “Youth and Democratic Participation” will be held in Yaoundé. The headmaster asks for two motivated candidates from Terminale A. To apply, you must fill in a registration form, write a short motivation letter and take part in an interview — all in English. Do you know the difference between a “workshop”, a “seminar” and a “training course”? Can you say that you want to “attend” a course without falling into the trap of “assister à”? This lesson gives you all the words you need.
\end{textebox}

\textbf{Problématique.} Dans le monde anglophone, on ne dit pas simplement \eng{a course} pour tout : chaque type de formation a son mot précis, ses verbes associés et ses documents. Comment nommer correctement une formation sur la démocratie, dire que l'on s'y inscrit, qu'on y participe et qu'on la termine ? C'est l'objectif de cette première section : maîtriser le vocabulaire des \cle{types de formations} et des \cle{cours} liés à l'éducation civique.

\courssub{Observation : un témoignage authentique}

Lisons d'abord le témoignage de Clarisse, une élève de Terminale qui a participé à une formation à Yaoundé. Soulignez mentalement tous les mots liés à la formation.

\begin{textebox}[Clarisse's testimony]
Last July, I attended a five-day workshop on democratic governance organised by the Ministry of Youth in Yaoundé. More than fifty young leaders from all ten regions took part in the training. On the first day, we filled in a registration form and received a badge and a detailed programme. The workshop was very practical: in small groups, we designed a project to encourage young people to vote. Two trainers guided us, and a guest lecturer from the university gave a lecture on the history of elections in Africa. We also attended a one-day seminar where we debated freedom of expression. At the end of the course, every participant who completed the training received a certificate. It was much more useful than a simple conference, because we produced something concrete: an action plan for our schools. Next year, I want to enrol in a refresher course to update my knowledge, and perhaps join an online webinar on human rights.
\end{textebox}

\textbf{Glossaire.}

\begin{center}
\begin{tabularx}{\linewidth}{|X|X|X|}\hline
\rowcolor{popblueL} English & Nature & Français \\ \hline
to attend & verbe & assister à, suivre (un cours) \\ \hline
workshop & nom & atelier (formation pratique) \\ \hline
to take part in & verbe & participer à \\ \hline
registration form & nom & formulaire d'inscription \\ \hline
badge & nom & badge, insigne d'identification \\ \hline
trainer & nom & formateur \\ \hline
lecture & nom & cours magistral, conférence \\ \hline
certificate & nom & certificat, attestation \\ \hline
refresher course & nom & cours de remise à niveau \\ \hline
webinar & nom & webinaire (séminaire en ligne) \\ \hline
\end{tabularx}
\end{center}

\textbf{Questions de compréhension.}
\begin{enumerate}
\item How long did the workshop last? \emph{(Five days.)}
\item What did the participants receive on the first day? \emph{(A badge and a programme.)}
\item Why was the workshop more useful than a conference, according to Clarisse? \emph{(Because participants produced something concrete.)}
\end{enumerate}

\courssub{Les types de formations : huit mots à distinguer}

\begin{definition}[Les principaux types de formations]
Voici les huit termes essentiels, du plus pratique au plus théorique :
\begin{itemize}
\item \cle{training course} : un cours de formation, général et structuré, souvent sur plusieurs jours (une formation).
\item \cle{workshop} : a practical training session where participants work in groups and produce something concrete (un atelier).
\item \cle{seminar} : a small-group course where participants discuss a topic with a specialist (un séminaire, un groupe de discussion).
\item \cle{conference} : a large formal event where speakers address a big audience (une conférence, un congrès).
\item \cle{lecture} : a talk given by one expert to an audience that mostly listens and takes notes (un cours magistral).
\item \cle{webinar} : a seminar held on the internet (un webinaire).
\item \cle{refresher course} : a short course to update knowledge you already have (un cours de remise à niveau / de recyclage).
\item \cle{summer school} : intensive courses organised during the summer holidays, often at a university (une université d'été).
\end{itemize}
\end{definition}

\begin{propriete}[La distinction clé : workshop, seminar, conference, lecture]
Ces quatre mots ne sont \textbf{pas interchangeables}. Retenez deux critères : la \textbf{taille du groupe} et le \textbf{degré de participation}.
\begin{itemize}
\item \eng{lecture} : un expert parle, le public écoute (grand groupe, passif) — « cours magistral ».
\item \eng{conference} : plusieurs orateurs, très grand public, événement formel — « congrès ».
\item \eng{seminar} : petit groupe, discussion guidée par un spécialiste — « séminaire ».
\item \eng{workshop} : petit groupe, travail pratique, production concrète — « atelier ».
\end{itemize}
\end{propriete}

\begin{center}
\begin{tabularx}{\linewidth}{|X|X|X|}\hline
\rowcolor{popblueL} English & Français & Caractéristique \\ \hline
training course & formation (cours) & terme général, structuré \\ \hline
workshop & atelier & pratique, en groupes, production concrète \\ \hline
seminar & séminaire & petit groupe, discussion \\ \hline
conference & conférence, congrès & grand public, discours formels \\ \hline
lecture & cours magistral & un expert parle, on écoute \\ \hline
webinar & webinaire & séminaire en ligne \\ \hline
refresher course & cours de recyclage & remise à niveau courte \\ \hline
summer school & université d'été & cours intensifs pendant les vacances \\ \hline
\end{tabularx}
\end{center}

\begin{popfigure}
\begin{tikzpicture}[scale=1]
\draw[thick, popblue, -{Stealth}] (0,0) -- (10.5,0);
\node[align=center] at (10.5,-0.7) {\small \textbf{pratique}};
\node[align=center] at (0,-0.7) {\small \textbf{théorique}};
\node[draw=popdark, fill=popblueL, rounded corners, align=center, minimum width=2.6cm] at (1.3,1.2) {\textbf{lecture}\\ \small cours magistral};
\node[draw=popdark, fill=poporangeL, rounded corners, align=center, minimum width=2.6cm] at (5,1.2) {\textbf{seminar}\\ \small discussion};
\node[draw=popdark, fill=popgreenL, rounded corners, align=center, minimum width=2.6cm] at (8.7,1.2) {\textbf{workshop}\\ \small atelier pratique};
\fill[popblue] (1.3,0) circle (0.09);
\fill[poporange] (5,0) circle (0.09);
\fill[popgreen] (8.7,0) circle (0.09);
\draw[dashed, popmuted] (1.3,0) -- (1.3,0.7);
\draw[dashed, popmuted] (5,0) -- (5,0.7);
\draw[dashed, popmuted] (8.7,0) -- (8.7,0.7);
\end{tikzpicture}
\legende{Du plus théorique au plus pratique : lecture, seminar, workshop.}
\end{popfigure}

\courssub{Carte mentale du lexique de la formation}

\begin{popfigure}
\begin{tikzpicture}[mindmap, grow cyclic, every node/.style={concept}, concept color=popblueL, text=popdark, root concept/.append style={concept color=popblue, text=white, font=\small\bfseries, align=center}, level 1/.append style={level distance=3.4cm, sibling angle=90, font=\small\bfseries}, level 2/.append style={level distance=2.4cm, font=\footnotesize, concept color=popcream}]
\node[root concept]{TRAINING ON DEMOCRACY}
  child[concept color=poporangeL] { node {types}
    child { node {workshop} }
    child { node {seminar} }
    child { node {lecture} }
  }
  child[concept color=popgreenL] { node {people}
    child { node {trainer} }
    child { node {trainee} }
    child { node {participant} }
  }
  child[concept color=poppinkL] { node {actions}
    child { node {to attend} }
    child { node {to enrol in} }
    child { node {to complete} }
  }
  child[concept color=popgoldL] { node {documents}
    child { node {certificate} }
    child { node {registration form} }
    child { node {handout} }
  };
\end{tikzpicture}
\legende{Carte mentale : les quatre familles du vocabulaire de la formation.}
\end{popfigure}

\courssub{Expressions de durée et d'organisation}

\begin{propriete}[Dire la durée : l'adjectif composé avec trait d'union]
Devant un nom, la durée devient un \textbf{adjectif composé} : \emph{nombre + trait d'union + unité au SINGULIER}.
\begin{itemize}
\item \eng{a three-day workshop} « un atelier de trois jours » (jamais \emph{a three-days workshop}).
\item \eng{a one-week training course} « une formation d'une semaine ».
\item \eng{a five-day seminar} « un séminaire de cinq jours ».
\item \eng{a two-hour lecture} « un cours de deux heures ».
\end{itemize}
Sans trait d'union, l'unité reprend son \emph{s} : \eng{The workshop lasted three days.} « L'atelier a duré trois jours. »
\end{propriete}

\begin{propriete}[Présentiel ou en ligne ?]
\begin{itemize}
\item \eng{a face-to-face seminar} « un séminaire en présentiel ».
\item \eng{an online session} « une séance en ligne ».
\item \eng{a distance-learning course} « une formation à distance ».
\end{itemize}
Notez les traits d'union dans \eng{face-to-face} et \eng{distance-learning} devant le nom.
\end{propriete}

\begin{exemplebox}[Exemples traduits]
\begin{itemize}
\item \eng{Our school organised a two-day workshop on elections.} « Notre école a organisé un atelier de deux jours sur les élections. »
\item \eng{The face-to-face seminar takes place in Douala; the online session is on Friday.} « Le séminaire en présentiel a lieu à Douala ; la séance en ligne est vendredi. »
\item \eng{She followed a six-week online training course.} « Elle a suivi une formation en ligne de six semaines. »
\end{itemize}
\end{exemplebox}

\courssub{Les verbes associés et les documents}

\begin{propriete}[Cinq verbes à maîtriser, avec leurs prépositions]
\begin{itemize}
\item \cle{to attend} + nom (SANS préposition) : \eng{to attend a workshop} « assister à un atelier ». Attention : pas de \emph{to} !
\item \cle{to take part in} + nom : \eng{to take part in a seminar} « participer à un séminaire ».
\item \cle{to enrol in} + nom (britannique ; américain : \emph{enroll}) : \eng{to enrol in a course} « s'inscrire à une formation ».
\item \cle{to register for} + nom : \eng{to register for a webinar} « s'inscrire à un webinaire ».
\item \cle{to complete} + nom : \eng{to complete a course} « suivre une formation jusqu'au bout, la valider ».
\end{itemize}
\end{propriete}

\begin{definition}[Les documents de la formation]
\begin{itemize}
\item \cle{registration form} : le formulaire d'inscription (\eng{to fill in a registration form} « remplir un formulaire »).
\item \cle{certificate} : le certificat, l'attestation délivrée en fin de formation.
\item \cle{programme} (britannique) / \cle{schedule} : le programme, l'emploi du temps de la formation.
\item \cle{handout} : le polycopié, le document distribué aux participants.
\item \cle{badge} : le badge porté par les participants.
\end{itemize}
\end{definition}

\begin{exemplebox}[Exemples traduits]
\begin{itemize}
\item \eng{She enrolled in a seminar on human rights.} « Elle s'est inscrite à un séminaire sur les droits de l'homme. »
\item \eng{Participants received a certificate at the end of the course.} « Les participants ont reçu un certificat à la fin de la formation. »
\item \eng{Please fill in the registration form before Monday.} « Veuillez remplir le formulaire d'inscription avant lundi. »
\item \eng{The trainer gave us a handout on voting procedures.} « Le formateur nous a distribué un polycopié sur les procédures de vote. »
\end{itemize}
\end{exemplebox}

\begin{attention}[Les pièges des francophones]
\begin{itemize}
\item Ne dites \textbf{jamais} \eng{a formation} pour traduire « une formation » : dites \eng{a training course} ou \eng{a training}. En anglais, \emph{formation} signifie « formation de rochers » ou « formation d'une équipe » (\eng{rock formation}, \eng{team formation}).
\item \eng{to attend} ne prend \textbf{pas} de préposition : \eng{I attended a conference}, et non \emph{I attended to a conference}.
\item \eng{to assist} signifie « aider » (\eng{to assist a customer} « aider un client »), pas « assister à ». Faux ami classique !
\item « S'inscrire à » se dit \eng{to enrol in} ou \eng{to register for}, jamais \emph{to inscribe}.
\item \eng{a lecture} n'est pas « une lecture » (qui se dit \emph{reading}) mais « un cours magistral ».
\end{itemize}
\end{attention}

\begin{exoresolu}[Complétez avec le mot correct]
Énoncé. Complétez chaque phrase avec le mot qui convient : \emph{workshop, conference, enrolled, certificate, three-day, attend}.
\begin{enumerate}
\item Last month, I \dots\ in a refresher course on citizenship.
\item The \dots\ workshop brought together thirty young leaders in Bamenda.
\item In the \dots\ , we built a real project in groups of five.
\item More than two thousand people came to the international \dots\ on democracy.
\item After completing the course, every trainee received a \dots\ .
\item Will you \dots\ the webinar on human rights tomorrow?
\end{enumerate}
\tcblower
Solution détaillée.
\begin{enumerate}
\item \textbf{enrolled} : « s'inscrire à » se dit \eng{to enrol in} ; le contexte est au passé (\emph{last month}).
\item \textbf{three-day} : adjectif composé devant le nom, avec trait d'union et \emph{day} au singulier.
\item \textbf{workshop} : on travaille en groupes et on produit quelque chose de concret — c'est la définition de l'atelier.
\item \textbf{conference} : un très grand public (deux mille personnes) écoute des discours — c'est le congrès.
\item \textbf{certificate} : le document remis en fin de formation.
\item \textbf{attend} : « assister à » sans préposition ; ici après \emph{will}, l'infinitif sans \emph{to}.
\end{enumerate}
\end{exoresolu}

\begin{exercice}[À vous de jouer]
Traduisez en anglais :
\begin{enumerate}
\item J'ai participé à un séminaire d'une semaine sur la démocratie à Yaoundé.
\item Elle s'est inscrite à un cours de remise à niveau en ligne.
\item Tous les stagiaires ont rempli le formulaire d'inscription et reçu un badge.
\end{enumerate}
\end{exercice}

\begin{corrige}[Exercice « À vous de jouer »]
\begin{enumerate}
\item \eng{I took part in a one-week seminar on democracy in Yaoundé.} (adjectif composé : \emph{one-week}, sans \emph{s}).
\item \eng{She enrolled in an online refresher course.}
\item \eng{All the trainees filled in the registration form and received a badge.}
\end{enumerate}
\end{corrige}

\begin{savaistu}[Summer schools et civics courses]
Au Royaume-Uni et aux États-Unis, les \eng{summer schools} et les \eng{civics courses} (cours d'éducation civique) sont très courants pour former les jeunes à la vie démocratique : on y apprend comment fonctionnent le Parlement, les élections et les droits des citoyens. Beaucoup d'universités américaines organisent chaque été des programmes intensifs où des lycéens simulent des débats parlementaires — une excellente préparation à la citoyenneté active, comme les ateliers que le Cameroun organise pour ses jeunes leaders.
\end{savaistu}

\begin{aretenir}[L'essentiel de la section]
\begin{itemize}
\item Huit types de formations : \eng{training course, workshop, seminar, conference, lecture, webinar, refresher course, summer school}.
\item Échelle : \eng{lecture} (théorique) $\rightarrow$ \eng{seminar} (discussion) $\rightarrow$ \eng{workshop} (pratique).
\item Durée : \eng{a three-day workshop} (trait d'union, singulier).
\item Verbes : \eng{to attend} (sans préposition), \eng{to take part in}, \eng{to enrol in}, \eng{to register for}, \eng{to complete}.
\item Documents : \eng{registration form, certificate, programme, handout, badge}.
\item Piège majeur : « une formation » = \eng{a training course}, jamais \emph{a formation}.
\end{itemize}
\end{aretenir}

\coursec{Les acteurs de la formation}

Dans la section précédente, nous avons découvert les différents \cle{types de formations et de cours} : \eng{workshop}, \eng{seminar}, \eng{training course}, \eng{lecture}. Mais une formation ne se déroule jamais toute seule : derrière chaque \eng{workshop} sur la démocratie, il y a des \cle{personnes} et des \cle{institutions} qui organisent, financent, animent et suivent la formation. C'est précisément le vocabulaire de ces \cle{acteurs} (\eng{the people and bodies involved}) que nous allons étudier maintenant.

\courssub{Observation : un texte pour découvrir les acteurs}

Lisons d'abord ce court extrait d'article de presse. Soulignez mentalement tous les mots qui désignent des personnes ou des institutions.

\begin{textebox}[A civic education workshop in Bamenda — press report]
The workshop was facilitated by two trainers from a local NGO. The keynote speaker, a law professor from the University of Buea, delivered a lecture on citizens' rights. About forty participants, mostly youth leaders and teachers, attended the three-day training course. The event was organised by a civil society network and funded by the British Council, with technical support from the National Commission on Human Rights and Freedoms. At the closing ceremony, the guest of honour, a representative of the Ministry of Territorial Administration, congratulated the organisers and encouraged the trainees to share their new knowledge in their communities. One experienced facilitator explained that active participation was the key to success: “We do not impose ideas; we help people discuss and decide together.”
\end{textebox}

\textbf{Glossaire :}

\begin{center}
\begin{tabularx}{\linewidth}{|l|l|X|}
\hline
\rowcolor{popblueL} \textbf{English} & \textbf{Nature} & \textbf{Français} \\
\hline
to facilitate & verb & animer, faciliter (une discussion) \\
\hline
keynote speaker & noun & intervenant principal (qui ouvre les débats) \\
\hline
to deliver a lecture & verb phrase & faire un cours magistral, donner une conférence \\
\hline
to attend & verb & assister à, participer à \\
\hline
civil society & noun & la société civile \\
\hline
to fund & verb & financer \\
\hline
guest of honour & noun & invité d'honneur \\
\hline
trainee & noun & stagiaire, personne en formation \\
\hline
\end{tabularx}
\end{center}

\textbf{Questions de compréhension :}
\begin{enumerate}
\item Who facilitated the workshop?
\item What did the keynote speaker do?
\item Who organised the event, and who funded it?
\item What does the facilitator mean by “we do not impose ideas”?
\end{enumerate}

\begin{corrige}[Questions de compréhension]
\begin{enumerate}
\item \eng{The workshop was facilitated by two trainers from a local NGO.} « L'atelier a été animé par deux formateurs d'une ONG locale. »
\item \eng{The keynote speaker delivered a lecture on citizens' rights.} « L'intervenant principal a donné une conférence sur les droits des citoyens. »
\item \eng{A civil society network organised the event, and the British Council funded it.} « Un réseau de la société civile a organisé l'événement, et le British Council l'a financé. »
\item \eng{He means that a facilitator does not force his own opinions on the group; he only helps people exchange ideas and make their own decisions.} « Il veut dire qu'un animateur n'impose pas ses opinions au groupe ; il aide seulement les gens à échanger des idées et à prendre leurs propres décisions. »
\end{enumerate}
\end{corrige}

\courssub{Les personnes : trainer, trainee, speaker et les autres}

\begin{definition}[Les personnes impliquées dans une formation]
\begin{itemize}
\item \cle{trainer} : a person who teaches practical skills during a training course (un formateur).
\item \cle{facilitator} : a person who helps a group discuss and work together, without imposing ideas (un animateur).
\item \cle{trainee} : a person who is being trained (un stagiaire, une personne en formation).
\item \cle{participant} : a person who takes part in an event (un participant).
\item \cle{organiser} (AmE \eng{organizer}) : the person or body that plans and prepares the event (l'organisateur).
\item \cle{speaker} : a person who gives a talk or a speech (un intervenant, un orateur).
\item \cle{lecturer} : a university teacher who gives lectures (un maître de conférences, un chargé de cours).
\item \cle{guest of honour} : the important official invited to open or close a ceremony (l'invité d'honneur).
\end{itemize}
\end{definition}

\begin{propriete}[Formation des mots : suffixes de personnes]
Observez la régularité de la formation de ces noms de métiers et de rôles :
\begin{itemize}
\item verbe + \cle{-er / -or} = la personne qui fait l'action : \eng{train} $\rightarrow$ \eng{trainer} ; \eng{organise} $\rightarrow$ \eng{organiser} ; \eng{facilitate} $\rightarrow$ \eng{facilitator} ; \eng{lecture} $\rightarrow$ \eng{lecturer}.
\item verbe + \cle{-ee} = la personne qui \emph{subit} l'action : \eng{train} $\rightarrow$ \eng{trainee} (celui qu'on forme) ; \eng{employ} $\rightarrow$ \eng{employee} ; \eng{interview} $\rightarrow$ \eng{interviewee}.
\end{itemize}
\eng{The trainer trains the trainee.} « Le formateur forme le stagiaire. » — La paire \eng{-er / -ee} est un outil précieux pour enrichir votre vocabulaire. Attention à l'orthographe : \eng{trainee} prend deux « e » (comme \eng{employee}), mais \eng{trainer} n'en prend qu'un.
\end{propriete}

\begin{attention}[Faux ami : “animator”]
« Un animateur » de formation ou de débat se dit \eng{facilitator} ou \eng{trainer}, jamais \eng{animator}. En anglais, \cle{animator} désigne un dessinateur de films d'animation (comme chez Disney). De même, \eng{animation} signifie « le dessin animé », pas « l'animation » d'un atelier. Pour « l'animation » d'une session, utilisez \eng{facilitation}.
\end{attention}

\courssub{Les institutions : NGO, civil society, the Ministry...}

\begin{definition}[Les institutions actrices de la formation civique]
\begin{itemize}
\item \cle{NGO} (non-governmental organisation) : une ONG, organisation indépendante de l'État. On prononce les lettres : « ène-dji-o ».
\item \cle{civil society} : la société civile, c'est-à-dire l'ensemble des associations, syndicats, ONG et mouvements de citoyens.
\item \cle{the Ministry} : le ministère (ex. \eng{the Ministry of Secondary Education} = le MINESEC).
\item \cle{the British Council} : l'organisme britannique de coopération culturelle et éducative, très actif au Cameroun.
\item \cle{the UN} (the United Nations) : l'ONU, avec ses agences comme \eng{UNESCO} ou \eng{UNDP}.
\item \cle{the National Commission on Human Rights and Freedoms} : la Commission nationale des droits de l'homme et des libertés (CNDHL), au Cameroun.
\end{itemize}
\end{definition}

\begin{propriete}[Article et majuscules avec les institutions]
\begin{itemize}
\item Les sigles se prononcent lettre par lettre et prennent généralement l'article \eng{the} : \eng{the UN}, \eng{the British Council}. Mais \eng{an NGO} (sans \eng{the}) quand on parle d'une ONG quelconque : \eng{She works for an NGO in Douala.} « Elle travaille pour une ONG à Douala. »
\item Les noms officiels d'institutions prennent la \cle{majuscule} à chaque mot important : \eng{the Ministry of Territorial Administration}, \eng{the University of Buea}. Les petits mots (\eng{of}, \eng{the}, \eng{and}) restent en minuscules.
\item \eng{civil society} est un nom commun abstrait : pas de majuscule, pas d'article : \eng{Civil society plays a key role in democracy.} « La société civile joue un rôle clé dans la démocratie. »
\end{itemize}
\end{propriete}

\begin{savaistu}[Le British Council au Cameroun]
Le \eng{British Council} est présent au Cameroun depuis des décennies. Il y soutient l'enseignement de l'anglais, organise des examens internationaux comme l'IELTS et finance des programmes de formation des enseignants et des jeunes leaders. Beaucoup de lycéens camerounais connaissent aussi ses bibliothèques et ses clubs d'anglais (\eng{English clubs}) à Yaoundé et à Douala.
\end{savaistu}

\courssub{Collocations et verbes d'action des acteurs}

En anglais, certains adjectifs « vont avec » certains noms, et certains verbes « vont avec » certaines fonctions. Ce sont les \cle{collocations}. Les apprendre par blocs évite les calques du français.

\begin{center}
\begin{tabularx}{\linewidth}{|X|X|X|}
\hline
\rowcolor{popblueL} \textbf{Collocation} & \textbf{Traduction} & \textbf{Remarque} \\
\hline
a qualified trainer & un formateur qualifié & \eng{qualified} = qui a un diplôme, une qualification \\
\hline
an experienced facilitator & un animateur expérimenté & \eng{experienced} se prononce « iks-pi-rienst » \\
\hline
an active participant & un participant actif & \eng{active}, pas \eng{actual} (faux ami : \eng{actual} = réel) \\
\hline
a keynote speaker & l'intervenant principal & \eng{keynote} = « note clé », discours d'ouverture \\
\hline
a resource person & une personne-ressource & expert invité pour ses connaissances \\
\hline
\end{tabularx}
\end{center}

\begin{center}
\begin{tabularx}{\linewidth}{|X|X|X|}
\hline
\rowcolor{popgreenL} \textbf{Verbe d'action} & \textbf{Français} & \textbf{Exemple} \\
\hline
to organise a workshop & organiser un atelier & The NGO organised a workshop in Buea. \\
\hline
to run a session & animer, mener une session & She runs training sessions every month. \\
\hline
to facilitate a debate & animer un débat & He facilitated the group discussion. \\
\hline
to deliver a lecture & donner une conférence & The professor delivered a lecture on rights. \\
\hline
to fund a programme & financer un programme & The British Council funds the project. \\
\hline
to sponsor an event & parrainer un événement & A local bank sponsored the seminar. \\
\hline
to attend a course & suivre une formation & Forty trainees attended the course. \\
\hline
\end{tabularx}
\end{center}

\begin{exemplebox}[Exemples traduits]
\eng{The NGO runs training sessions in rural areas.} « L'ONG anime des sessions de formation en zone rurale. »

\eng{Who funds this programme?} « Qui finance ce programme ? »

\eng{An experienced facilitator guided the group work.} « Un animateur expérimenté a guidé le travail de groupe. »

\eng{The keynote speaker delivered a lecture on voting rights.} « L'intervenant principal a donné une conférence sur le droit de vote. »

\eng{The training was sponsored by the British Council.} « La formation a été parrainée par le British Council. »
\end{exemplebox}

\begin{attention}[Pièges de traduction]
\begin{itemize}
\item « Assister à une formation » = \eng{to attend a training course}, jamais \eng{to assist a training}. \eng{To assist} signifie « aider » : \eng{She assisted the trainer.} « Elle a aidé le formateur. »
\item « Animer une session » = \eng{to run / to facilitate a session}. \eng{To animate} existe mais est rare et signifie « donner vie à ».
\item « Un intervenant » = \eng{a speaker}, pas \eng{an intervener}.
\item « Suivre une formation » = \eng{to attend a course} ou \eng{to take a course}, jamais \eng{to follow a training} (calque du français).
\end{itemize}
\end{attention}

\courssub{Trainer, teacher ou lecturer ? Choisir le bon mot}

Ces trois mots se traduisent tous par « celui qui enseigne », mais ils ne sont pas interchangeables.

\begin{methode}[Comment choisir entre trainer, facilitator, speaker et lecturer]
\begin{enumerate}
\item Posez-vous la question : \emph{que fait exactement la personne ?}
\item Si elle enseigne un \cle{savoir-faire pratique} en formation professionnelle ou civique $\rightarrow$ \eng{trainer} (formateur technique).
\item Si elle \cle{aide un groupe à discuter} sans imposer ses idées $\rightarrow$ \eng{facilitator} (animateur de débat).
\item Si elle \cle{fait un discours} lors d'une cérémonie ou d'une conférence $\rightarrow$ \eng{speaker} (intervenant).
\item Si elle \cle{enseigne à l'université} et fait des cours magistraux $\rightarrow$ \eng{lecturer} (professeur d'université).
\item Si elle enseigne à l'école ou au collège/lycée $\rightarrow$ \eng{teacher}.
\end{enumerate}
\end{methode}

\begin{center}
\begin{tabularx}{\linewidth}{|l|l|X|}
\hline
\rowcolor{popblueL} \textbf{Mot} & \textbf{Français} & \textbf{Contexte d'emploi} \\
\hline
trainer & formateur & training courses, workshops, sport \\
\hline
facilitator & animateur & debates, group work, seminars \\
\hline
teacher & enseignant & primary and secondary schools \\
\hline
lecturer & chargé de cours & universities, colleges \\
\hline
speaker & intervenant & conferences, ceremonies \\
\hline
\end{tabularx}
\end{center}

\begin{popfigure}
\begin{tikzpicture}[
  node distance=1.6cm,
  box/.style={draw=popblue, fill=popblueL, rounded corners, thick, align=center, minimum width=3.4cm, minimum height=1cm, font=\small\bfseries},
  lab/.style={font=\small\itshape, text=popdark}
]
\node[box, align=center] (org){ORGANISER\\ (NGO)};
\node[box, below left=1.6cm and -0.5cm of org, align=center] (trainer){TRAINER /\\ FACILITATOR};
\node[box, below right=1.6cm and -0.5cm of org, align=center] (part){PARTICIPANTS /\\ TRAINEES};
\draw[-{Stealth}, very thick, poporange] (org) -- node[lab, above left, align=center]{funds and\\ organises} (trainer);
\draw[-{Stealth}, very thick, popgreen] (trainer) -- node[lab, below, align=center]{trains} (part);
\draw[-{Stealth}, very thick, poppurple] (part) -- node[lab, above right, align=center]{attend the\\ sessions of} (org);
\end{tikzpicture}
\legende{Les acteurs de la formation : qui finance, qui forme, qui participe.}
\end{popfigure}

\courssub{Exercice résolu et entraînement}

\begin{exoresolu}[Complétez avec le bon acteur]
Énoncé — Complétez chaque phrase avec \eng{trainer}, \eng{facilitator}, \eng{trainee}, \eng{keynote speaker} ou \eng{organiser}.

1. The \dots\ of the workshop booked the hall and sent the invitations.
2. During the debate, the \dots\ made sure everybody could speak.
3. Every \dots\ received a certificate at the end of the course.
4. The \dots\ showed us how to fill in a voter registration form.
5. The \dots\ , a professor from the University of Buea, opened the conference.

\tcblower
Solution détaillée :
\begin{enumerate}
\item \eng{organiser} — c'est la personne qui planifie l'événement (réserver la salle, envoyer les invitations).
\item \eng{facilitator} — son rôle est d'aider le groupe à discuter, de donner la parole à tous.
\item \eng{trainee} — celui qui suit la formation et reçoit le certificat (suffixe \eng{-ee} = personne qui subit l'action).
\item \eng{trainer} — il enseigne un savoir-faire pratique (remplir un formulaire d'inscription électorale).
\item \eng{keynote speaker} — c'est l'intervenant principal qui ouvre la conférence par un discours.
\end{enumerate}
\end{exoresolu}

\begin{exercice}[À vous de jouer]
Traduisez en anglais :
\begin{enumerate}
\item L'ONG organise des ateliers dans les zones rurales.
\item Qui a financé ce programme de formation ?
\item La conférencière a parlé des droits des citoyens.
\item Un animateur expérimenté a guidé nos discussions de groupe.
\item Quarante stagiaires ont suivi la formation de trois jours.
\end{enumerate}
\end{exercice}

\begin{corrige}[Exercice : traductions]
\begin{enumerate}
\item \eng{The NGO organises workshops in rural areas.}
\item \eng{Who funded this training programme?}
\item \eng{The lecturer spoke about citizens' rights.}
\item \eng{An experienced facilitator guided our group discussions.}
\item \eng{Forty trainees attended the three-day training course.}
\end{enumerate}
\end{corrige}

\begin{aretenir}[L'essentiel de la section]
\begin{itemize}
\item Les personnes : \eng{trainer} (formateur), \eng{facilitator} (animateur), \eng{trainee} (stagiaire), \eng{participant}, \eng{organiser}, \eng{keynote speaker}, \eng{lecturer}, \eng{guest of honour}.
\item Les institutions : \eng{an NGO}, \eng{civil society}, \eng{the Ministry}, \eng{the British Council}, \eng{the UN}, \eng{the National Commission on Human Rights and Freedoms}.
\item Verbes d'action : \eng{to organise}, \eng{to run a session}, \eng{to facilitate}, \eng{to deliver a lecture}, \eng{to fund}, \eng{to sponsor}, \eng{to attend}.
\item Pièges : « animateur » $\neq$ \eng{animator} ; « assister à » = \eng{to attend}, pas \eng{to assist} ; « suivre une formation » = \eng{to attend a course}, pas \eng{to follow a training}.
\end{itemize}
\end{aretenir}

\coursec{Types de formations et de cours}

Avant d'étudier les thèmes démocratiques, il faut maîtriser le vocabulaire qui permet de décrire \emph{le cadre} de la formation : les types d'événements, leur organisation et leur déroulement. Ce lexique est indispensable pour comprendre un programme, une invitation ou un rapport d'activité.

\courssub{Les principaux types d'événements de formation}

\begin{definition}[Types de formations]
En contexte camerounais et international, on distingue plusieurs formats :
\begin{itemize}
\item \eng{a training course} / \eng{a training programme} : une formation (sur plusieurs jours/semaines), structurée, avec un curriculum.
\item \eng{a workshop} : un atelier (travaux pratiques, participatif, souvent 1 à 3 jours).
\item \eng{a seminar} : un séminaire (exposés d'experts, discussions, public plus large).
\item \eng{a capacity-building workshop} : un atelier de renforcement des capacités (terme très utilisé par les ONG et les partenaires au développement).
\item \eng{a civic education session} : une session d'éducation civique (souvent à l'école ou dans les centres de jeunes).
\item \eng{a refresher course} : une session de recyclage / de mise à niveau (pour actualiser des connaissances).
\end{itemize}
\end{definition}

\begin{exemplebox}[Exemples en contexte]
\begin{itemize}
\item \eng{The NGO organised a three-day workshop on good governance in Buea.} « L'ONG a organisé un atelier de trois jours sur la bonne gouvernance à Buea. »
\item \eng{The Ministry of Youth Affairs runs civic education programmes in secondary schools.} « Le Ministère de la Jeunesse mène des programmes d'éducation civique dans les lycées. »
\item \eng{Local councillors attended a refresher course on the rule of law.} « Les conseillers municipaux ont suivi une session de recyclage sur l'État de droit. »
\end{itemize}
\end{exemplebox}

\courssub{Organisation et déroulement : le lexique logistique}

\begin{center}
\begin{tabularx}{\linewidth}{|X|X|X|}\hline
\rowcolor{popblueL} English & Français & Collocation fréquente \\ \hline
\eng{venue} & le lieu (de l'événement) & \eng{the training venue} \\ \hline
\eng{schedule} / \eng{timetable} & l'emploi du temps, le programme horaire & \eng{a tight schedule} \\ \hline
\eng{agenda} & l'ordre du jour & \eng{to set the agenda} \\ \hline
\eng{session} & la séance, la session & \eng{a morning session} \\ \hline
\eng{module} & le module & \eng{Module 1: Human Rights} \\ \hline
\eng{facilitator} & l'animateur, le facilitateur & \eng{the lead facilitator} \\ \hline
\eng{resource person} & la personne ressource, l'expert & \eng{invite a resource person} \\ \hline
\eng{participant} / \eng{trainee} & le participant / le stagiaire & \eng{select participants} \\ \hline
\eng{beneficiary} & le bénéficiaire & \eng{direct beneficiaries} \\ \hline
\eng{certificate of attendance} & l'attestation de participation & \eng{issue certificates} \\ \hline
\eng{evaluation form} & la fiche d'évaluation & \eng{fill in the evaluation form} \\ \hline
\end{tabularx}
\end{center}

\begin{attention}[Facilitator vs Teacher / Trainer]
\eng{A facilitator} (animateur/facilitateur) guide les échanges et la réflexion du groupe (approche participative). \eng{A trainer} (formateur) transmet des compétences techniques précises. \eng{A teacher} (enseignant) suit un programme scolaire. Dans les ateliers démocratiques, on préfère \eng{facilitator}. Ne dites pas \eng{animator} (qui désigne un dessinateur de dessins animés ou un animateur de club de vacances).
\end{attention}

\courssub{Verbes d'action pour décrire une formation}

\begin{center}
\begin{tabularx}{\linewidth}{|X|X|}\hline
\rowcolor{popblueL} Verbe (forme base / prétérit / participe) & Sens / Collocation \\ \hline
\eng{organise / organised / organised} & organiser (\eng{organise a workshop}) \\ \hline
\eng{conduct / conducted / conducted} & mener, animer (\eng{conduct a session}) \\ \hline
\eng{facilitate / facilitated / facilitated} & faciliter, animer (\eng{facilitate a discussion}) \\ \hline
\eng{deliver / delivered / delivered} & dispenser, présenter (\eng{deliver a presentation}) \\ \hline
\eng{attend / attended / attended} & assister à, participer à (\eng{attend a course}) \\ \hline
\eng{enrol (Br) / enrolled / enrolled} & s'inscrire (\eng{enrol in a programme}) \\ \hline
\eng{register / registered / registered} & s'inscrire (sur une liste) (\eng{register for the workshop}) \\ \hline
\eng{participate / participated / participated} & participer (\eng{participate actively}) \\ \hline
\eng{benefit / benefited / benefited} & bénéficier de (\eng{benefit from the training}) \\ \hline
\eng{assess / assessed / assessed} & évaluer (\eng{assess the outcomes}) \\ \hline
\end{tabularx}
\end{center}

\begin{propriete}[Attend : pas de préposition]
\eng{Attend} est transitif direct en anglais : \eng{attend a meeting}, \eng{attend the workshop}. On ne dit jamais \eng{attend to a meeting} (sauf sens « s'occuper de »). Erreur fréquente : \eng{*I attended to the training}. Correct : \eng{I attended the training.}
\end{propriete}

\begin{experience}[Jeu de rôle : Organiser un atelier]
\textbf{Consigne :} En binôme. L'élève A est un \eng{programme officer} d'une ONG à Douala. L'élève B est le \eng{principal} d'un lycée. A appelle B pour organiser un \eng{civic education workshop} pour les élèves de Terminale.
\textbf{Mots à utiliser :} \eng{venue, schedule, facilitator, participants, module, certificate of attendance, register}.
\textbf{Durée :} 3 minutes par binôme.
\end{experience}

\coursec{Les acteurs de la formation}

Une formation implique des rôles précis. Les identifier permet de comprendre qui fait quoi dans un rapport ou un article de presse.

\courssub{Tableau des acteurs}

\begin{center}
\begin{tabularx}{\linewidth}{|X|X|X|}\hline
\rowcolor{popblueL} English & Français & Rôle typique \\ \hline
\eng{organiser} / \eng{coordinator} & l'organisateur / le coordinateur & Planifie, gère la logistique, le budget. \\ \hline
\eng{facilitator} / \eng{trainer} & l'animateur / le formateur & Anime les séances, transmet le savoir. \\ \hline
\eng{resource person} / \eng{expert} & la personne ressource / l'expert & Apporte une expertise pointue sur un thème. \\ \hline
\eng{participant} / \eng{trainee} / \eng{beneficiary} & le participant / le stagiaire / le bénéficiaire & Suit la formation, met en œuvre les acquis. \\ \hline
\eng{mentor} / \eng{coach} & le mentor / le coach & Accompagne sur le long terme (post-formation). \\ \hline
\eng{stakeholder} & la partie prenante & Tout acteur concerné (État, ONG, communauté). \\ \hline
\end{tabularx}
\end{center}

\begin{exemplebox}[Acteurs en phrases]
\begin{itemize}
\item \eng{The coordinator ensured the venue was ready before the trainees arrived.} « Le coordinateur a veillé à ce que le lieu soit prêt avant l'arrivée des stagiaires. »
\item \eng{Stakeholders from the Ministry and civil society attended the opening ceremony.} « Les parties prenantes du Ministère et de la société civile ont assisté à la cérémonie d'ouverture. »
\item \eng{Each trainee was assigned a mentor for the follow-up phase.} « Chaque stagiaire s'est vu assigner un mentor pour la phase de suivi. »
\end{itemize}
\end{exemplebox}

\begin{attention}[Stakeholder : un mot-clé du développement]
\eng{Stakeholder} (partie prenante) désigne toute personne ou groupe ayant un intérêt dans le projet : État, collectivités, ONG, chefs traditionnels, jeunes, femmes. Traduction approximative : « acteur concerné ». Ne pas confondre avec \eng{shareholder} (actionnaire, entreprise).
\end{attention}

\coursec{Thèmes démocratiques abordés en formation}

Dans les sections précédentes, nous avons identifié les \cle{types de formations} et les \cle{acteurs}. Mais que leur enseigne-t-on exactement ? Quels sujets sont au programme d'un \eng{training course on democracy} ? C'est précisément l'objet de cette section : le vocabulaire des \cle{thèmes démocratiques}, les \cle{valeurs} qui les sous-tendent, les \cle{collocations} qui permettent d'en parler correctement, et les mots qui décrivent les \cle{résultats} d'une formation.

\courssub{Observer d'abord : un texte de formation}

Lisons d'abord ce court extrait du rapport d'un formateur après un atelier organisé à Yaoundé. Soulignez mentalement tous les mots qui désignent des thèmes ou des valeurs démocratiques.

\begin{textebox}[A trainer's report — Yaoundé]
During the training, participants discussed good governance, transparency and the fight against corruption. They learnt that every citizen has a say in public life. On the first day, the facilitator introduced the rule of law: the principle that everyone, including leaders, must obey the law. The trainees then worked in groups on human rights and freedom of speech. Many young women said they wanted to make their voices heard in their communities.

On the second day, the workshop focused on elections and citizenship. The participants learnt how to register to vote and why peaceful elections matter. They also talked about gender equality and the need to empower young people. At the end of the course, each group prepared an action plan: some promised to raise awareness about voting in their schools, others decided to stand up for the rights of children in their villages. The trainer was pleased: the trainees had gained knowledge, skills and a real commitment to democracy.
\end{textebox}

\textbf{Glossaire}

\begin{center}
\begin{tabularx}{\linewidth}{|X|X|X|}\hline
\rowcolor{popblueL} Mot anglais & Nature & Traduction française \\ \hline
\eng{good governance} & nom & la bonne gouvernance \\ \hline
\eng{transparency} & nom & la transparence \\ \hline
\eng{the rule of law} & nom & l'État de droit \\ \hline
\eng{to obey the law} & verbe & respecter la loi \\ \hline
\eng{freedom of speech} & nom & la liberté d'expression \\ \hline
\eng{to register to vote} & verbe & s'inscrire sur les listes électorales \\ \hline
\eng{gender equality} & nom & l'égalité des sexes \\ \hline
\eng{to empower} & verbe & donner les moyens d'agir, autonomiser \\ \hline
\eng{an action plan} & nom & un plan d'action \\ \hline
\eng{commitment} & nom & l'engagement \\ \hline
\end{tabularx}
\end{center}

\textbf{Questions de compréhension}

\begin{enumerate}
\item What principle means that even leaders must obey the law?
\item What did each group prepare at the end of the course?
\item Name two themes discussed on the second day.
\end{enumerate}

\courssub{Le lexique des grands thèmes démocratiques}

\begin{definition}[Les thèmes d'une formation à la démocratie]
Un \eng{training course on democracy} (une formation à la démocratie) porte généralement sur neuf grands thèmes : \eng{democracy} (la démocratie), \eng{citizenship} (la citoyenneté), \eng{human rights} (les droits de l'homme), \eng{good governance} (la bonne gouvernance), \eng{elections} (les élections), \eng{the rule of law} (l'État de droit), \eng{freedom of speech} (la liberté d'expression), \eng{gender equality} (l'égalité des sexes) et \eng{peace education} (l'éducation à la paix).
\end{definition}

\begin{definition}[The rule of law]
\eng{The rule of law} désigne \emph{the principle that everyone, including leaders, must obey the law} — le principe selon lequel tout le monde, y compris les dirigeants, doit obéir à la loi. En français : \cle{l'État de droit}. Attention à l'article : on dit \eng{the rule of law}, jamais \eng{a rule of law}.
\end{definition}

\begin{exemplebox}[Les thèmes en contexte]
\begin{itemize}
\item \eng{The workshop deals with citizenship and human rights.} « L'atelier porte sur la citoyenneté et les droits de l'homme. »
\item \eng{Elections must be free, fair and peaceful.} « Les élections doivent être libres, équitables et pacifiques. »
\item \eng{Freedom of speech allows citizens to criticise bad governance.} « La liberté d'expression permet aux citoyens de critiquer la mauvaise gouvernance. »
\item \eng{Peace education is taught in many schools in Bamenda.} « L'éducation à la paix est enseignée dans de nombreuses écoles de Bamenda. »
\item \eng{Gender equality means that boys and girls have the same opportunities.} « L'égalité des sexes signifie que les garçons et les filles ont les mêmes chances. »
\end{itemize}
\end{exemplebox}

\begin{attention}[Human rights : toujours au pluriel]
On dit \eng{human rights} au pluriel : \eng{to defend human rights}, \eng{human rights education}. Le singulier \eng{a human right} n'existe que pour désigner UN droit précis : \eng{Education is a human right.} « L'éducation est un droit humain. » De même, \eng{freedom of speech} est invariable et sans article : jamais \eng{the freedom of speech} dans le sens général.
\end{attention}

\courssub{Le lexique des valeurs démocratiques}

Une formation à la démocratie ne transmet pas seulement des notions : elle transmet des \cle{valeurs} (\eng{values}). Voici les six valeurs essentielles à connaître.

\begin{center}
\begin{tabularx}{\linewidth}{|X|X|X|}\hline
\rowcolor{popblueL} English & Français & Exemple de collocation \\ \hline
\eng{tolerance} & la tolérance & \eng{to promote tolerance} \\ \hline
\eng{dialogue} & le dialogue & \eng{to encourage dialogue} \\ \hline
\eng{accountability} & la redevabilité, l'obligation de rendre compte & \eng{to demand accountability} \\ \hline
\eng{transparency} & la transparence & \eng{to ensure transparency} \\ \hline
\eng{participation} & la participation & \eng{to increase participation} \\ \hline
\eng{responsibility} & la responsabilité & \eng{to take responsibility} \\ \hline
\end{tabularx}
\end{center}

\begin{exemplebox}[Les valeurs en phrases]
\begin{itemize}
\item \eng{Leaders must show transparency and accountability.} « Les dirigeants doivent faire preuve de transparence et de redevabilité. »
\item \eng{Dialogue is better than violence.} « Le dialogue vaut mieux que la violence. »
\item \eng{Every citizen should take responsibility for the community.} « Chaque citoyen devrait prendre ses responsabilités envers la communauté. »
\item \eng{The training encourages the participation of women in public life.} « La formation encourage la participation des femmes à la vie publique. »
\end{itemize}
\end{exemplebox}

\begin{attention}[Accountability : un mot sans équivalent direct]
\eng{Accountability} n'a pas de traduction française unique : selon le contexte, il signifie « la redevabilité », « l'obligation de rendre compte » ou « la responsabilité (devant le peuple) ». Ne le confondez pas avec \eng{responsibility}, qui est plus général. Et méfiez-vous de \eng{dialogue} : il s'écrit pareil en français, mais se prononce « daïa-log » (/ˈdaɪəlɒɡ/) en anglais.
\end{attention}

\courssub{Carte mentale : organiser le lexique autour de DEMOCRACY}

Pour mémoriser ce vocabulaire, rien de tel qu'une carte mentale. Le mot central \eng{DEMOCRACY} se ramifie en quatre branches : les droits (\eng{rights}), les devoirs (\eng{duties}), les valeurs (\eng{values}) et les institutions (\eng{institutions}).

\begin{popfigure}
\begin{center}
\begin{tikzpicture}[mindmap, grow cyclic,
every node/.style={concept, align=center, font=\small},
root concept/.append style={concept color=popblue, font=\bfseries},
level 1/.append style={level distance=3.2cm, sibling angle=90},
level 2/.append style={level distance=2.4cm, sibling angle=45, font=\scriptsize}]
\node[root concept]{DEMOCRACY}
  child[concept color=poporange] { node {rights} 
      child { node {human rights} }
      child { node {freedom of speech} } }
  child[concept color=popgreen] { node {duties}
      child { node {obey the law} }
      child { node {vote} } }
  child[concept color=poppurple] { node {values}
      child { node {tolerance} }
      child { node {transparency} } }
  child[concept color=popgold] { node {institutions}
      child { node {elections} }
      child { node {parliament} } };
\end{tikzpicture}
\end{center}
\legende{Carte mentale du champ lexical de la démocratie : quatre branches (droits, devoirs, valeurs, institutions).}
\end{popfigure}

\begin{methode}[Construire votre propre carte mentale lexicale]
\begin{enumerate}
\item Écrivez le mot central au milieu d'une page (ici \eng{DEMOCRACY}).
\item Choisissez 3 ou 4 catégories (branches principales) : \eng{rights}, \eng{duties}, \eng{values}, \eng{institutions}.
\item Ajoutez sous chaque branche 2 ou 3 mots appris dans la leçon.
\item Complétez la carte après chaque leçon de vocabulaire : une carte vivante est un outil de révision idéal pour le baccalauréat.
\end{enumerate}
\end{methode}

\courssub{Les collocations essentielles}

En anglais, un nom appelle un verbe précis : on ne peut pas combiner n'importe quel verbe avec n'importe quel nom. Ces combinaisons figées s'appellent des \cle{collocations}. Les voici, à apprendre par cœur.

\begin{center}
\begin{tabularx}{\linewidth}{|X|X|}\hline
\rowcolor{popblueL} Collocation anglaise & Traduction française \\ \hline
\eng{to promote democracy} & promouvoir la démocratie \\ \hline
\eng{to defend human rights} & défendre les droits de l'homme \\ \hline
\eng{to respect the rule of law} & respecter l'État de droit \\ \hline
\eng{to raise awareness} & sensibiliser \\ \hline
\eng{to empower young people} & donner aux jeunes les moyens d'agir \\ \hline
\eng{to fight corruption} & lutter contre la corruption \\ \hline
\eng{to hold free elections} & organiser des élections libres \\ \hline
\eng{to protect citizens} & protéger les citoyens \\ \hline
\end{tabularx}
\end{center}

\begin{propriete}[Collocation figée : to RAISE awareness, jamais « to rise awareness »]
Pour dire « sensibiliser », l'anglais utilise la collocation figée \eng{to raise awareness}. Deux verbes se ressemblent mais ne s'emploient pas de la même façon :
\begin{itemize}
\item \eng{to raise} (raise, raised, raised) est \cle{transitif} : il exige un complément d'objet. \eng{The trainer raised awareness about voting.} « Le formateur a sensibilisé au vote. »
\item \eng{to rise} (rise, rose, risen) est \cle{intransitif} : il n'a jamais de complément d'objet. \eng{Prices are rising.} « Les prix augmentent. »
\end{itemize}
On dira donc toujours \eng{to raise awareness}, jamais \eng{to rise awareness}.
\end{propriete}

\begin{attention}[« Sensibiliser » ne se traduit pas par « to sensitize »]
Le verbe \eng{to sensitize} existe en anglais médical/technique (rendre sensible), mais pas dans le sens « faire prendre conscience » en langage courant : c'est un calque du français. Pour traduire « sensibiliser le public à la démocratie », dites :
\begin{itemize}
\item \eng{to raise awareness about democracy} (la solution la plus fréquente) ;
\item \eng{to sensitise the public to democracy} (anglais britannique ; \eng{to sensitize} en anglais américain — à éviter dans un contexte camerounais francophone).
\end{itemize}
Autre piège : « une sensibilisation » ne se dit pas \eng{a sensitisation} dans la langue courante ; préférez \eng{an awareness campaign} « une campagne de sensibilisation » ou \eng{an awareness-raising campaign}.
\end{attention}

\courssub{Expressions idiomatiques : donner la parole aux citoyens}

Trois expressions idiomatiques reviennent constamment dans les textes sur la citoyenneté. Elles ne se traduisent pas mot à mot.

\begin{definition}[Trois expressions idiomatiques à connaître]
\begin{itemize}
\item \eng{to have a say (in something)} : avoir son mot à dire (dans quelque chose). \eng{Every citizen has a say in public life.} « Chaque citoyen a son mot à dire dans la vie publique. »
\item \eng{to make one's voice heard} : se faire entendre. \eng{Young people must make their voices heard.} « Les jeunes doivent se faire entendre. »
\item \eng{to stand up for one's rights} : défendre ses droits. \eng{She stood up for her rights.} « Elle a défendu ses droits. »
\end{itemize}
\end{definition}

\begin{exemplebox}[Exemples traduits]
\begin{itemize}
\item \eng{Young people must make their voices heard.} « Les jeunes doivent se faire entendre. »
\item \eng{The course empowers women to take part in elections.} « La formation donne aux femmes les moyens de participer aux élections. »
\item \eng{Villagers stood up for their rights during the meeting.} « Les villageois ont défendu leurs droits pendant la réunion. »
\item \eng{Do students have a say in school decisions?} « Les élèves ont-ils leur mot à dire dans les décisions de l'école ? »
\end{itemize}
\end{exemplebox}

\begin{attention}[Attention à l'accord de « one's »]
Dans \eng{to make one's voice heard} et \eng{to stand up for one's rights}, le possessif s'accorde avec le sujet : \eng{I make my voice heard}, \eng{you make your voice heard}, \eng{they make their voices heard}. Erreur fréquente des francophones : écrire \eng{they make their voice heard} avec \eng{voice} au singulier alors que le sujet est pluriel — préférez le pluriel \eng{their voices}.
\end{attention}

\courssub{Le vocabulaire des résultats d'une formation}

À la fin d'une formation, on évalue ce que les participants ont gagné. Cinq noms décrivent ces résultats.

\begin{center}
\begin{tabularx}{\linewidth}{|X|X|X|}\hline
\rowcolor{popblueL} English & Français & Collocation typique \\ \hline
\eng{skills} & les compétences & \eng{to acquire new skills} \\ \hline
\eng{knowledge} (indénombrable) & les connaissances & \eng{to gain knowledge} \\ \hline
\eng{awareness} & la conscience, la sensibilisation & \eng{to raise awareness} \\ \hline
\eng{commitment} & l'engagement & \eng{to show commitment} \\ \hline
\eng{action plan} & le plan d'action & \eng{to draw up an action plan} \\ \hline
\end{tabularx}
\end{center}

\begin{exemplebox}[Les résultats en phrases]
\begin{itemize}
\item \eng{The trainees gained knowledge about elections.} « Les stagiaires ont acquis des connaissances sur les élections. »
\item \eng{Each group drew up an action plan for its community.} « Chaque groupe a élaboré un plan d'action pour sa communauté. »
\item \eng{The participants showed great commitment to peace.} « Les participants ont montré un grand engagement en faveur de la paix. »
\end{itemize}
\end{exemplebox}

\begin{attention}[Knowledge est indénombrable]
\eng{Knowledge} ne prend jamais de \eng{s} et jamais l'article indéfini \eng{a}. On dit \eng{He gained a lot of knowledge} « Il a acquis beaucoup de connaissances », jamais \eng{he gained knowledges}. Pour parler d'une connaissance précise, utilisez \eng{a piece of knowledge} ou changez de mot (\eng{a fact}, \eng{a skill}).
\end{attention}

\coursec{Faux amis et pièges des francophones}

Cette section regroupe les erreurs les plus fréquentes observées dans les copies du baccalauréat camerounais sur ce thème.

\begin{center}
\begin{tabularx}{\linewidth}{|X|X|X|}\hline
\rowcolor{popblueL} Faux ami / Piège & Erreur fréquente & Forme correcte \\ \hline
\eng{sensitize} & \eng{*The NGO sensitized the youth.} & \eng{The NGO raised awareness among the youth.} \\ \hline
\eng{knowledges} & \eng{*They acquired many knowledges.} & \eng{They acquired a lot of knowledge.} \\ \hline
\eng{rise awareness} & \eng{*We must rise awareness.} & \eng{We must raise awareness.} \\ \hline
\eng{a human right} (général) & \eng{*Human right is important.} & \eng{Human rights are important.} \\ \hline
\eng{the freedom of speech} & \eng{*The freedom of speech is a right.} & \eng{Freedom of speech is a right.} \\ \hline
\eng{animator} & \eng{*The animator spoke well.} & \eng{The facilitator / trainer spoke well.} \\ \hline
\eng{assist to} & \eng{*He assisted to the workshop.} & \eng{He attended the workshop.} \\ \hline
\eng{benefit of} & \eng{*They benefited of the training.} & \eng{They benefited from the training.} \\ \hline
\eng{claim} (revendiquer) & \eng{*They claimed their rights.} & \eng{They demanded / stood up for their rights.} \\ \hline
\eng{control} (contrôler) & \eng{*The trainer controlled the list.} & \eng{The trainer checked / monitored the list.} \\ \hline
\end{tabularx}
\end{center}

\begin{propriete}[Claim vs Demand vs Stand up for]
\eng{To claim} signifie « affirmer, prétendre » (souvent sans preuve) ou « réclamer (un bagage, une indemnité) ». Pour « revendiquer un droit », utilisez \eng{to demand} (exiger) ou l'idiome \eng{to stand up for one's rights}.
\end{propriete}

\begin{propriete}[Control vs Check / Monitor]
\eng{To control} signifie « commander, diriger, maîtriser » (ex. \eng{control the budget}, \eng{control the crowd}). Pour « vérifier une liste, surveiller un examen », utilisez \eng{to check} ou \eng{to monitor}.
\end{propriete}

\coursec{Familles de mots, prononciation et orthographe}

Maîtriser la dérivation permet de passer du nom au verbe, à l'adjectif ou à l'adverbe sans faute d'orthographe ni erreur de catégorie grammaticale.

\courssub{Tableau de dérivation (Word families)}

\begin{center}
\begin{tabularx}{\linewidth}{|X|X|X|X|}\hline
\rowcolor{popblueL} Verbe & Nom (chose/concept) & Nom (personne) & Adjectif / Adverbe \\ \hline
\eng{democratise} & \eng{democracy} & \eng{democrat} & \eng{democratic} / \eng{democratically} \\ \hline
\eng{elect} & \eng{election} & \eng{voter / candidate} & \eng{electoral} / \eng{electorally} \\ \hline
\eng{govern} & \eng{government / governance} & \eng{governor} & \eng{governmental} \\ \hline
\eng{participate} & \eng{participation} & \eng{participant} & \eng{participatory} / \eng{actively} \\ \hline
\eng{empower} & \eng{empowerment} & — & \eng{empowered} \\ \hline
\eng{transparise} (rare) & \eng{transparency} & — & \eng{transparent} / \eng{transparently} \\ \hline
\eng{account for} & \eng{accountability} & — & \eng{accountable} \\ \hline
\eng{corrupt} & \eng{corruption} & \eng{corrupt official} & \eng{corrupt} / \eng{corruptly} \\ \hline
\eng{legislate} & \eng{legislation} & \eng{legislator} & \eng{legislative} \\ \hline
\eng{advocate} & \eng{advocacy} & \eng{advocate} & \eng{advocacy} (adj. inv.) \\ \hline
\end{tabularx}
\end{center}

\courssub{Prononciation : accents toniques et pièges}

En anglais, l'accent tonique (stress) change souvent de place dans la famille de mots. Marquez l'accent sur la syllabe en \textbf{gras}.

\begin{center}
\begin{tabularx}{\linewidth}{|X|X|}\hline
\rowcolor{popblueL} Mot & Prononciation approximative (français) \\ \hline
\eng{de\bf{moc}racy} & « di-mo-kra-si » (accent sur 2e) \\ \hline
\eng{demo\bf{cra}tic} & « de-mo-kra-tik » (accent sur 3e) \\ \hline
\eng{citi\bf{zen}ship} & « si-ti-zen-ship » (accent sur 2e) \\ \hline
\eng{go\bf{vern}ance} & « go-ve-nans » (accent sur 2e) \\ \hline
\eng{e\bf{lec}tion} & « i-lek-chon » (accent sur 2e, \eng{c} = /k/) \\ \hline
\eng{trans\bf{pa}ren\bf{cy}} & « trans-pa-ren-si » (accent sur 2e et 4e) \\ \hline
\eng{ac\bf{coun}ta\bf{bi}li\bf{ty}} & « a-koun-ta-bi-li-ti » (accents secondaires) \\ \hline
\eng{par\bf{ti}ci\bf{pa}tion} & « pa-ti-si-peï-chon » (accent sur 3e) \\ \hline
\eng{em\bf{pow}er} & « em-pa-o » (diphtongue /aʊ/) \\ \hline
\end{tabularx}
\end{center}

\begin{attention}[Pièges d'orthographe]
\begin{itemize}
\item \eng{Democracy} : un seul \eng{m}, \eng{c} avant \eng{r}.
\item \eng{Governance} : garde le \eng{n} de \eng{govern}, pas \eng{*governence}.
\item \eng{Transparency} : \eng{paren\bf{c}y}, pas \eng{*transparence} (français).
\item \eng{Accountability} : double \eng{c}, double \eng{t}, \eng{bi}, \eng{li}, \eng{ty}.
\item \eng{Legislation} : \eng{sl}, pas \eng{*legistlation}.
\item \eng{Advocacy} : \eng{c} après \eng{v}, pas \eng{*advocasy}.
\end{itemize}
\end{attention}

\coursec{Le lexique en contexte : dialogue et production guidée}

Cette section met en œuvre le vocabulaire dans des situations de communication authentiques : compréhension orale/écrite (dialogue), expression orale (débat/jeu de rôle) et expression écrite (rédaction de rapport).

\courssub{Compréhension : Dialogue entre deux stagiaires}

\begin{textebox}[Dialogue : After the workshop in Douala]
\textbf{Context:} Two trainees, Achille and Mireille, are talking near the Wouri river after a three-day workshop on \eng{youth participation in local governance}.

\eng{Achille:} That was an intense workshop, wasn't it? I really liked the module on \eng{accountability}. 

\eng{Mireille:} Me too. The facilitator explained that \eng{accountability} means leaders must answer to the people. It's not just \eng{transparency} — showing what you do — but \eng{answering for} what you do.

\eng{Achille:} Exactly. And the session on \eng{action plans} was practical. Our group drew up a plan to \eng{raise awareness} about \eng{voter registration} in our neighbourhoods before the next council elections.

\eng{Mireille:} Good idea. Young people often don't \eng{register to vote} because they think their vote doesn't count. We need to \eng{empower} them to \eng{make their voices heard}.

\eng{Achille:} And to \eng{stand up for} their rights if the process isn't \eng{free and fair}. I feel I've gained real \eng{skills} in \eng{advocacy}.

\eng{Mireille:} Same here. Let's keep in touch and \eng{monitor} the implementation of our plans.
\end{textebox}

\textbf{Glossaire du dialogue}

\begin{center}
\begin{tabularx}{\linewidth}{|X|X|}\hline
\rowcolor{popblueL} Expression & Traduction \\ \hline
\eng{answer to / answer for} & rendre des comptes à / répondre de \\ \hline
\eng{voter registration} & l'inscription sur les listes électorales \\ \hline
\eng{free and fair} & libres et équitables (collocation fixe pour élections) \\ \hline
\eng{advocacy} & le plaidoyer, la défense d'une cause \\ \hline
\eng{monitor} & suivre, surveiller (la mise en œuvre) \\ \hline
\eng{keep in touch} & rester en contact \\ \hline
\end{tabularx}
\end{center}

\courssub{Expression orale : Débat structuré}

\begin{experience}[Débat : « La jeunesse camerounaise et la démocratie »]
\textbf{Consigne :} En groupes de 4. Préparez 5 minutes d'arguments pour ou contre la motion : \eng{``Young Cameroonians have the power to change local governance through civic engagement.''}
\textbf{Rôles :} 1 \eng{Chairperson} (modérateur), 2 \eng{Speakers for}, 1 \eng{Speaker against} (ou 2/2).
\textbf{Lexique imposé (minimum 5 par intervenant) :} \eng{empowerment, accountability, voter registration, raise awareness, stand up for, make one's voice heard, action plan, participation, transparency, stakeholder}.
\textbf{Évaluation :} Fluidité, justesse du vocabulaire, interaction (répondre aux arguments adverses).
\end{experience}

\courssub{Expression écrite : Rédaction d'un rapport de formation}

\begin{methode}[Rédiger un \eng{training report} (compte-rendu de formation)]
Un rapport suit une structure standard. Respectez-la pour le bac (Composition / Guided Writing).
\begin{enumerate}
\item \textbf{Header (En-tête) :} \eng{REPORT ON THE WORKSHOP ON [THEME]} \\ \eng{Date: ... | Venue: ... | Facilitator: ... | Participants: ...}
\item \textbf{Introduction (1 paragraphe) :} But, organisateur, date, lieu, nombre de participants. \eng{The workshop aimed to... / was organised by...}
\item \textbf{Content / Proceedings (2-3 paragraphes) :} Thèmes abordés par jour/session. Verbes au prétérit. \eng{On Day 1, the facilitator introduced... Participants discussed... Groups worked on...}
\item \textbf{Outcomes / Results (1 paragraphe) :} Ce qui a été produit. \eng{The trainees acquired skills in... They drew up action plans... They showed commitment to...}
\item \textbf{Recommendations / Way Forward (1 paragraphe) :} \eng{It is recommended that... / Participants suggested... / Follow-up activities should include...}
\item \textbf{Conclusion (1 phrase) :} \eng{The workshop was a success.}
\item \textbf{Signature :} \eng{Report written by [Name], [Function].}
\end{enumerate}
\end{methode}

\begin{exemplebox}[Modèle de rapport (extrait)]
\eng{REPORT ON THE CIVIC EDUCATION WORKSHOP} \\
\eng{Date: 12--14 March 2025 | Venue: Council Hall, Buea | Facilitator: Mr. N. Ekema} \\
\eng{Introduction: The three-day workshop aimed to empower youth leaders on good governance and the rule of law. Thirty participants attended.} \\
\eng{Proceedings: On Day 1, the facilitator introduced the rule of law and human rights. Participants debated freedom of speech. On Day 2, the focus was on elections and gender equality. Trainees learnt how to register to vote. On Day 3, groups designed action plans to raise awareness in schools.} \\
\eng{Outcomes: Participants gained knowledge on citizenship. They acquired advocacy skills. Each group submitted an action plan.} \\
\eng{Recommendations: Organise refresher courses quarterly. Involve more rural stakeholders.}
\end{exemplebox}

\courssub{Exercice résolu et entraînement}

\begin{exoresolu}[Complétez avec la collocation correcte]
Énoncé — Complétez chaque phrase avec l'une des expressions suivantes : \eng{raise awareness}, \eng{stand up for}, \eng{have a say}, \eng{fight corruption}, \eng{empower}.
\begin{enumerate}
\item The NGO wants to \ldots\ young people in rural areas.
\item Journalists help to \ldots\ by reporting on dishonest leaders.
\item Citizens should \ldots\ in local decisions.
\item The club organised a campaign to \ldots\ about human rights.
\item It is important to \ldots\ your rights when they are threatened.
\end{enumerate}
\tcblower
Solution détaillée :
\begin{enumerate}
\item \eng{empower} — « donner les moyens d'agir » : la phrase parle de jeunes en zone rurale, c'est la collocation \eng{to empower young people}.
\item \eng{fight corruption} — les journalistes dénoncent les dirigeants malhonnêtes : \eng{to fight corruption} « lutter contre la corruption ».
\item \eng{have a say} — suivi de la préposition \eng{in} : \eng{to have a say in decisions} « avoir son mot à dire dans les décisions ».
\item \eng{raise awareness} — collocation figée avec \eng{about} ; jamais \eng{rise awareness}, car \eng{rise} est intransitif.
\item \eng{stand up for} — \eng{to stand up for your rights} « défendre vos droits ».
\end{enumerate}
\end{exoresolu}

\begin{exercice}[À vous de jouer]
\begin{enumerate}
\item Traduisez en anglais : « La formation sensibilise les jeunes à l'État de droit. »
\item Traduisez en anglais : « Chaque citoyen a son mot à dire dans la vie publique. »
\item Complétez : \eng{The workshop promotes gender \ldots\ and peace \ldots .}
\item Trouvez l'erreur et corrigez : \eng{The trainer rose awareness about elections and the trainees gained many knowledges.}
\item Rédigez trois phrases sur une formation imaginaire à Douala, en utilisant \eng{empower}, \eng{commitment} et \eng{action plan}.
\end{enumerate}
\end{exercice}


\begin{corrige}[Exercice « À vous de jouer »]
\begin{enumerate}
\item \eng{The training raises young people's awareness of the rule of law.} (ou \eng{The course sensitises young people to the rule of law.})
\item \eng{Every citizen has a say in public life.}
\item \eng{The workshop promotes gender equality and peace education.}
\item Deux erreurs : \eng{rose awareness} doit devenir \eng{raised awareness} (\eng{raise} est transitif, prétérit \eng{raised}) ; \eng{many knowledges} doit devenir \eng{a lot of knowledge} (\eng{knowledge} est indénombrable). Phrase correcte : \eng{The trainer raised awareness about elections and the trainees gained a lot of knowledge.}
\item Exemple de production : \eng{The training in Douala empowers young traders to defend their rights. The participants showed real commitment to democracy. At the end, each group drew up an action plan for its neighbourhood.}
\end{enumerate}
\end{corrige}

\begin{aretenir}[L'essentiel de la leçon 42]
\begin{itemize}
\item \textbf{Types de formation :} \eng{training course, workshop, seminar, capacity-building, civic education session}.
\item \textbf{Acteurs :} \eng{organiser, facilitator, resource person, participant/trainee, stakeholder, mentor}.
\item \textbf{Thèmes démocratiques :} \eng{democracy, citizenship, human rights, good governance, elections, the rule of law, freedom of speech, gender equality, peace education}.
\item \textbf{Valeurs :} \eng{tolerance, dialogue, accountability, transparency, participation, responsibility}.
\item \textbf{Collocations figées :} \eng{promote democracy, defend human rights, respect the rule of law, raise awareness, empower young people, fight corruption, hold free elections, protect citizens}.
\item \textbf{Idiomatismes :} \eng{have a say (in), make one's voice heard, stand up for one's rights}.
\item \textbf{Résultats :} \eng{skills, knowledge (U), awareness, commitment, action plan (draw up)}.
\item \textbf{Faux amis :} \eng{sensitize} (non), \eng{knowledges} (non), \eng{rise awareness} (non), \eng{animator} (non), \eng{assist to} (non), \eng{claim} (non pour droit), \eng{control} (non pour vérifier).
\item \textbf{Familles de mots :} \eng{democracy/democratic, election/electoral, governance/govern, participation/participatory, transparency/transparent, accountability/accountable, legislation/legislative}.
\end{itemize}
\end{aretenir}

\coursec{Faux amis et pièges des francophones}

Dans la section précédente, nous avons exploré les thèmes démocratiques abordés pendant les formations : \eng{Elections}, \eng{Human Rights}, \eng{Good Governance}, \eng{Civic Education}. Vous disposez maintenant d'un lexique riche. Mais attention : c'est précisément sur ce vocabulaire de la formation et de la citoyenneté que les francophones commettent le plus d'erreurs, car le français et l'anglais partagent de nombreux mots d'origine latine au sens différent. Cette section vous apprend à repérer ces \cle{faux amis} (\eng{false friends} ou \eng{false cognates}) et à éviter les \cle{calques} du français.

\courssub{Qu'est-ce qu'un faux ami ?}

\begin{definition}[Faux ami]
Un \cle{faux ami} est un mot anglais qui ressemble à un mot français par sa forme mais qui a un sens différent. Exemple : \eng{to assist} ressemble à « assister », mais signifie « aider ». Un \cle{calque} est une construction française traduite mot à mot en anglais, ce qui produit une phrase incorrecte : \eng{to assist to a course} est un calque de « assister à un cours ».
\end{definition}

Observez d'abord ces deux phrases authentiques :

\begin{exemplebox}[Observation : deux verbes, deux sens]
\eng{“Fifty trainees attended the course on democracy.”} = « Cinquante stagiaires ont assisté à la formation sur la démocratie. » (ils étaient présents)\par
\eng{“The trainer assisted a disabled participant during the workshop.”} = « Le formateur a aidé un participant handicapé pendant l'atelier. » (il lui a apporté son aide)
\end{exemplebox}

Même famille de mots, sens totalement différents. C'est tout l'enjeu de cette section.

\courssub{To attend ou to assist ? Le piège n° 1}

\begin{propriete}[To attend = assister à (être présent)]
\eng{To attend} + nom (sans préposition !) signifie « assister à, être présent à » : \eng{to attend a course}, \eng{to attend a seminar}, \eng{to attend school} (aller à l'école), \eng{to attend a meeting}.\par
\eng{To assist} (+ nom ou \eng{somebody in doing something}) signifie « aider » : \eng{“The facilitator assisted the participants.”} = « L'animateur a aidé les participants. »
\end{propriete}

\begin{attention}[Le piège n° 1 des candidats camerounais au Bac]
Chaque année, des milliers de candidats écrivent : \eng{“I assisted to the seminar.”} — FAUX, doublement faux !\par
1. \eng{To assist} ne signifie pas « assister » mais « aider » ;\par
2. \eng{To attend} se construit SANS préposition : on ne dit jamais \eng{attend to} pour « assister à ».\par
Dites : \eng{“I attended the seminar.”} = « J'ai assisté au séminaire. »\par
Attention aussi : \eng{to attend to something} existe, mais signifie « s'occuper de quelque chose » : \eng{“Please attend to this matter.”} = « Veuillez vous occuper de cette affaire. »
\end{attention}

\courssub{Formation, conférence, stage, animateur : quatre mots piégés}

\begin{propriete}[Quatre faux amis majeurs du lexique de la formation]
\begin{itemize}
\item « une formation » = \eng{a training} ou \eng{a training course}. Le mot anglais \eng{formation} désigne une formation géologique (\eng{rock formation}) ou la mise en place de quelque chose (\eng{the formation of a new party} = « la création d'un nouveau parti »).
\item « une conférence » (exposé) = \eng{a lecture} ; « une conférence » (grand colloque) = \eng{a conference}. \eng{A lecture} est un cours magistral : \eng{“The professor gave a lecture on human rights.”}
\item « un stage » = \eng{an internship} (américain) ou \eng{a placement / a work placement} (britannique). \eng{A stage} en anglais = « une scène » (théâtre) ou « une étape » : \eng{the first stage of the project}.
\item « un animateur » (de formation, d'atelier) = \eng{a facilitator} ou \eng{a trainer}. \eng{An animator} = un dessinateur de films d'animation.
\end{itemize}
\end{propriete}

\begin{exemplebox}[Exemples traduits]
\eng{“She is doing an internship at an NGO in Yaoundé.”} = « Elle fait un stage dans une ONG à Yaoundé. »\par
\eng{“The play was performed on an open-air stage.”} = « La pièce a été jouée sur une scène en plein air. »\par
\eng{“A good facilitator encourages everyone to speak.”} = « Un bon animateur encourage chacun à prendre la parole. »\par
\eng{“Walt Disney was a famous animator.”} = « Walt Disney était un célèbre dessinateur d'animation. »
\end{exemplebox}

\courssub{Sensibiliser, débat et autres pièges}

\begin{propriete}[Sensibiliser = to raise awareness]
« Sensibiliser le public à la démocratie » se dit \eng{to raise public awareness of democracy} ou \eng{to sensitise/sensitize people to democracy} (verbe rare et technique). Le verbe \eng{to sensibilize} n'existe pas en anglais courant. La forme la plus naturelle et la plus sûre est \cle{to raise awareness}.\par
\eng{“The NGO organised a campaign to raise awareness of voting rights.”} = « L'ONG a organisé une campagne pour sensibiliser aux droits de vote. »
\end{propriete}

\begin{propriete}[Un débat = a debate : prononciation]
\eng{A debate} se prononce « di-béït » : le \textbf{b} se prononce, contrairement à \eng{doubt} (« doute », prononcé « daout », b muet) et \eng{debt} (« dette », « dèt », b muet). Verbe : \eng{to debate an issue} = « débattre d'une question ».\par
\eng{“The trainees debated the issue of press freedom.”} = « Les stagiaires ont débattu de la question de la liberté de la presse. »
\end{propriete}

\begin{savaistu}[Actually ne veut pas dire « actuellement » !]
\eng{Actually} signifie « en fait, réellement » : \eng{“The course is actually free.”} = « En fait, la formation est gratuite. » Pour dire « actuellement », utilisez \eng{currently} ou \eng{at present} : \eng{“She is currently attending a training course in Buea.”} = « Elle suit actuellement une formation à Buea. » Même piège avec \eng{eventually} (« finalement ») et « éventuellement » (\eng{possibly}) !
\end{savaistu}

\courssub{Calques à éviter : tableau récapitulatif}

\begin{center}
\begin{tabularx}{\linewidth}{|X|X|X|}
\hline
\rowcolor{popblueL} Français & Calque à éviter (faux) & Forme correcte \\
\hline
assister à un cours & to assist to a course & to attend a course \\
\hline
faire une formation & to make a formation & to do / to follow a training course \\
\hline
suivre une formation & to follow a formation & to attend / to take a training course \\
\hline
animer un atelier & to animate a workshop & to facilitate / to run a workshop \\
\hline
sensibiliser les jeunes & to sensibilize the youth & to raise young people's awareness \\
\hline
faire un stage & to make a stage & to do an internship / a placement \\
\hline
une salle de formation & a formation room & a training room \\
\hline
un certificat de formation & a formation certificate & a training certificate \\
\hline
\end{tabularx}
\end{center}

\begin{popfigure}
\begin{tikzpicture}[
  faux/.style={rectangle, rounded corners, draw=poppink, fill=poppinkL, thick, align=center, text width=5.2cm, font=\small},
  ok/.style={rectangle, rounded corners, draw=popgreen, fill=popgreenL, thick, align=center, text width=5.2cm, font=\small},
  head/.style={font=\bfseries\small, align=center}
]
\node[head, text=poppink] at (-3.1,0) {Ce que dit le francophone (FAUX)};
\node[head, text=popgreen] at (3.1,0) {Ce qu'il faut dire (CORRECT)};
\node[faux] (f1) at (-3.1,-1.3) {\eng{“I assisted to the seminar.”}};
\node[ok] (o1) at (3.1,-1.3) {\eng{“I attended the seminar.”}};
\node[faux] (f2) at (-3.1,-2.7) {\eng{“She made a formation last year.”}};
\node[ok] (o2) at (3.1,-2.7) {\eng{“She did a training course last year.”}};
\node[faux] (f3) at (-3.1,-4.1) {\eng{“He animates the workshop.”}};
\node[ok] (o3) at (3.1,-4.1) {\eng{“He facilitates the workshop.”}};
\node[faux] (f4) at (-3.1,-5.5) {\eng{“We sensibilize the citizens.”}};
\node[ok] (o4) at (3.1,-5.5) {\eng{“We raise citizens' awareness.”}};
\foreach \a/\b in {f1/o1, f2/o2, f3/o3, f4/o4}
  \draw[-{Stealth}, very thick, poporange] (\a.east) -- (\b.west);
\end{tikzpicture}
\legende{Du calque français à la phrase anglaise correcte : quatre transformations à mémoriser.}
\end{popfigure}

\courssub{Méthode : comment déjouer un faux ami}

\begin{methode}[Trois réflexes anti-piège]
\begin{enumerate}
\item \textbf{Méfiez-vous des mots trop ressemblants.} Si un mot anglais ressemble exactement au français (\eng{assist}, \eng{formation}, \eng{stage}, \eng{actually}), allumez un signal d'alarme.
\item \textbf{Vérifiez le sens anglais dans le contexte.} Demandez-vous : « Cette phrase aurait-elle un sens pour un Anglais qui ne parle pas français ? » \eng{“I assisted to the seminar”} n'a aucun sens pour lui.
\item \textbf{Reformulez avec un verbe simple.} En cas de doute, remplacez par un mot sûr : \eng{help} (aider), \eng{go to} (aller à), \eng{follow / do} (suivre), \eng{run} (animer). Mieux vaut une phrase simple et correcte qu'un mot savant et faux.
\end{enumerate}
\end{methode}

\begin{exoresolu}[Corrige les erreurs]
Énoncé — Les phrases suivantes ont été écrites par des élèves francophones. Corrige chaque phrase.\par
1. \eng{“I assisted to a conference on human rights.”}\par
2. \eng{“My brother is making a formation in Douala.”}\par
3. \eng{“The animator of the workshop explained the rules.”}\par
4. \eng{“This campaign sensibilizes young voters.”}
\tcblower
Solution détaillée :\par
1. \eng{“I attended a lecture on human rights.”} — « assister à » = \eng{attend}, sans préposition ; une « conférence-exposé » = \eng{a lecture}.\par
2. \eng{“My brother is doing a training course in Douala.”} — « faire une formation » = \eng{do a training course} ; \eng{formation} est un faux ami.\par
3. \eng{“The facilitator of the workshop explained the rules.”} — un « animateur » d'atelier = \eng{a facilitator} ; \eng{animator} = dessinateur.\par
4. \eng{“This campaign raises young voters' awareness.”} — « sensibiliser » = \eng{to raise awareness} ; \eng{sensibilize} n'existe pas.
\end{exoresolu}

\begin{exercice}[À toi de jouer]
Traduis en anglais correct : 1. « Cinquante participants ont assisté au séminaire. » 2. « Elle suit actuellement un stage dans une ONG. » 3. « Le formateur a aidé les stagiaires. » 4. « Cette émission sensibilise le public aux élections. »
\end{exercice}

\begin{corrige}[Exercice : traductions]
1. \eng{“Fifty participants attended the seminar.”} 2. \eng{“She is currently doing an internship at an NGO.”} 3. \eng{“The trainer assisted the trainees.”} 4. \eng{“This programme raises public awareness of elections.”}
\end{corrige}

\begin{aretenir}[L'essentiel de la section]
\eng{attend} = assister à (sans préposition) ; \eng{assist} = aider. « Une formation » = \eng{a training course}, « un stage » = \eng{an internship}, « un animateur » = \eng{a facilitator}, « sensibiliser » = \eng{to raise awareness}, « actuellement » = \eng{currently}. En cas de doute, reformulez avec un verbe simple.
\end{aretenir}

\coursec{Familles de mots, prononciation et orthographe}

Après avoir identifié les \textbf{faux amis} et les pièges de traduction qui guettent l'apprenant francophone, nous devons maintenant structurer notre lexique pour le rendre productif. La maîtrise d'un thème ne se limite pas à connaître une liste de mots isolés ; elle exige de comprendre comment les mots se construisent (\textbf{morphologie}), comment ils se prononcent (\textbf{phonologie}) et comment ils s'écrivent (\textbf{orthographe}). C'est la clé pour passer de la compréhension passive à l'expression active, à l'oral comme à l'écrit.

\courssub{La dérivation nominale et adjectivale : observer les familles}

En anglais, comme en français, un mot racine (souvent un verbe) génère une famille entière par adjonction de \textbf{préfixes} (devant) et de \textbf{suffixes} (derrière). Observer ces régularités permet de deviner le sens d'un mot inconnu et de varier son vocabulaire dans une rédaction.

\begin{textebox}{Observation : Extraits d'un rapport de formation à Buea}
The \textbf{workshop} on \textbf{democratic governance} gathered fifty \textbf{participants} from civil society organisations. The main \textbf{trainer}, Mr. Atanga, emphasised the need for \textbf{civic education} in schools. He explained that \textbf{citizenship} is not only about \textbf{voting} but requires active \textbf{participation}. The \textbf{trainees} worked on \textbf{awareness}-raising strategies for the upcoming elections. Good \textbf{governance} depends on \textbf{responsible} leaders and an \textbf{informed} citizenry.
\end{textebox}

\courssub{Famille de \eng{train} (former / formation)}

\begin{definition}[Famille lexicale : TRAIN]
La racine \eng{train} (verbe \textbf{régulier} : \eng{train, trained, trained}) génère des noms d'agent, de patient et d'action.
\end{definition}

\begin{center}
\begin{tabularx}{\linewidth}{|X|X|X|}
\hline
\rowcolor{popblueL} \textbf{Forme} & \textbf{Nature / Sens} & \textbf{Traduction / Contexte} \\
\hline
\eng{to train} & Verbe (régulier) & Former, entraîner, instruire. \\
\hline
\eng{training} & Nom (non comptable) & Formation, entraînement. \eng{In-service training} = formation continue. \\
\hline
\eng{a trainer} & Nom (agent, suffixe \eng{-er}) & Formateur, entraîneur (celui qui fait l'action). \\
\hline
\eng{a trainee} & Nom (patient, suffixe \eng{-ee}) & Stagiaire, apprenant (celui qui \textbf{subit} l'action). \\
\hline
\eng{to retrain} & Verbe (préfixe \eng{re-}) & Réformer, recycler (changer de métier). \\
\hline
\eng{retraining} & Nom & Reconversion, recyclage professionnel. \\
\hline
\end{tabularx}
\end{center}

\begin{popfigure}
\begin{tikzpicture}[
    mindmap,
    grow cyclic,
    every node/.style={concept, execute at begin node=\hskip0pt},
    root concept/.append style={concept color=popblue, font=\large\bfseries, text width=2.5cm, align=center},
    level 1 concept/.append style={concept color=popblue!70!poppurple, font=\normalsize\bfseries, text width=2.2cm, align=center, sibling angle=90, level distance=4cm},
    level 2 concept/.append style={concept color=popgreen!60!popblue, font=\small, text width=2cm, align=center, sibling angle=45, level distance=3.5cm}
]
\node[root concept, align=center]{train\\ (verb)}
    child[concept color=poporange] { node[level 1 concept, align=center]{trainer\\ (agent -er) }
        child { node[level 2 concept] {co-trainer} }
        child { node[level 2 concept] {master trainer} }
    }
    child[concept color=popgreen] { node[level 1 concept, align=center]{trainee\\ (patient -ee) }
        child { node[level 2 concept] {trainee teacher} }
        child { node[level 2 concept] {graduate trainee} }
    }
    child[concept color=poppurple] { node[level 1 concept, align=center]{training\\ (action -ing) }
        child { node[level 2 concept] {on-the-job training} }
        child { node[level 2 concept] {training session} }
        child { node[level 2 concept] {training manual} }
    }
    child[concept color=poppink] { node[level 1 concept, align=center]{retrain\\ (re- + verb) }
        child { node[level 2 concept] {retraining} }
        child { node[level 2 concept] {retraining programme} }
    };
\end{tikzpicture}
\legende{Arbre morphologique de la famille TRAIN : dérivation par suffixes agentifs (-er), patientifs (-ee) et nominaux (-ing).}
\end{popfigure}

\begin{propriete}[Les suffixes d'agent et de patient : \eng{-er/-or} vs \eng{-ee}]
C'est une règle morphologique majeure pour éviter les confusions.
\begin{itemize}
    \item \textbf{\eng{-er} / \eng{-or}} (souvent pour les verbes d'origine latine) $\rightarrow$ \textbf{Celui qui FAIT l'action} (Agent).
    \item \textbf{\eng{-ee}} (emprunté au français participe passé \eng{-é}) $\rightarrow$ \textbf{Celui qui SUBIT l'action} (Patient / Bénéficiaire).
\end{itemize}
\begin{center}
\begin{tabularx}{\linewidth}{|X|X|X|}
\hline
\rowcolor{popblueL} \textbf{Verbe} & \textbf{Agent (\eng{-er/-or})} & \textbf{Patient (\eng{-ee})} \\
\hline
\eng{to train} & \eng{trainer} (formateur) & \eng{trainee} (stagiaire) \\
\hline
\eng{to employ} & \eng{employer} (employeur) & \eng{employee} (employé) \\
\hline
\eng{to interview} & \eng{interviewer} (journaliste) & \eng{interviewee} (personne interviewée) \\
\hline
\eng{to appoint} & \eng{appointer} (rare) / \eng{appointing authority} & \eng{appointee} (nommé) \\
\hline
\eng{to mentor} & \eng{mentor} & \eng{mentee} (mentoré) \\
\hline
\end{tabularx}
\end{center}
\eng{The trainer trains the trainees.} = « Le formateur forme les stagiaires. »
\end{propriete}

\begin{attention}[Prononciation et confusion \eng{trainer} vs \eng{trainee}]
C'est l'erreur n°1 des francophones.
\begin{itemize}
    \item \eng{Trainer} [\eng{'treɪnə(r)}] $\rightarrow$ se note « \textbf{tréï-neu} » (accent sur la 1ère syllabe, \eng{er} muet/\eng{ə}).
    \item \eng{Trainee} [\eng{treɪ'niː}] $\rightarrow$ se note « \textbf{tréï-ni} » (accent sur la \textbf{dernière} syllabe, \eng{ee} = \eng{/iː/} long).
\end{itemize}
\emph{Astuce :} Le suffixe \eng{-ee} attire toujours l'accent tonique sur la dernière syllabe (\eng{employ\textbf{ee}}, \eng{refug\textbf{ee}}, \eng{absent\textbf{ee}}).
\end{attention}

\courssub{Famille de \eng{participate} (participer)}

\begin{definition}[Famille lexicale : PARTICIPATE]
Verbe régulier \eng{participate, participated, participated}. Attention à l'orthographe du son \eng{/ʃ/} (\eng{ti} ou \eng{ci}) et à la préposition \eng{in}.
\end{definition}

\begin{center}
\begin{tabularx}{\linewidth}{|X|X|X|}
\hline
\rowcolor{popblueL} \textbf{Forme} & \textbf{Nature / Morphologie} & \textbf{Traduction / Collocation} \\
\hline
\eng{to participate} & Verbe + \eng{in} (pas \eng{to}) & Participer \textbf{à}. \eng{Participate in a workshop}. \\
\hline
\eng{participation} & Nom (\eng{-tion}) & Participation. \eng{Active participation}, \eng{voter participation}. \\
\hline
\eng{participant} & Nom / Adj (\eng{-ant}) & Participant (n.), participant (adj). \eng{Participant observation}. \\
\hline
\eng{participatory} & Adjectif (\eng{-ory}) & Participatif. \eng{Participatory approach}, \eng{participatory democracy}. \\
\hline
\eng{non-participation} & Nom (préfixe \eng{non-}) & Non-participation, abstention. \\
\hline
\end{tabularx}
\end{center}

\begin{exemplebox}[Collocations utiles pour l'expression écrite]
\begin{itemize}
    \item \eng{To \textbf{participate actively in} the debate.} (Participer activement au débat.)
    \item \eng{\textbf{Ensure broad participation} of youth.} (Assurer une large participation des jeunes.)
    \item \eng{A \textbf{participatory approach} to governance.} (Une approche participative de la gouvernance.)
    \item \eng{\textbf{Key participants} included local leaders.} (Les participants clés incluaient des leaders locaux.)
\end{itemize}
\end{exemplebox}

\courssub{Famille de \eng{democracy} (démocratie)}

\begin{definition}[Famille lexicale : DEMOCRACY]
Nom d'origine grecque \eng{democracy} (demos + kratos). L'accent tonique se déplace selon les suffixes (règle des mots en \eng{-cy} / \eng{-crat} / \eng{-cratic}).
\end{definition}

\begin{center}
\begin{tabularx}{\linewidth}{|X|X|X|}
\hline
\rowcolor{popblueL} \textbf{Forme} & \textbf{Prononciation (approx. fr.)} & \textbf{Traduction / Usage} \\
\hline
\eng{democracy} & \eng{di-'mo-kra-si} (\textbf{di-mo-kra-si}) & Démocratie (nom). \eng{Consolidate democracy}. \\
\hline
\eng{democrat} & \eng{'de-mo-krat} (\textbf{dé-mo-krat}) & Démocrate (personne). \\
\hline
\eng{democratic} & \eng{de-mo-'kra-tik} (\textbf{di-mo-kra-tik}) & Démocratique (adj). Accent sur \textbf{kra}. \eng{Democratic elections}. \\
\hline
\eng{democratically} & \eng{de-mo-'kra-ti-kli} & Démocratiquement. \eng{Elected democratically}. \\
\hline
\eng{democratise} (\eng{US} \eng{democratize}) & \eng{di-'mo-kra-taiz} & Démocratiser (verbe). \eng{To democratise the party}. \\
\hline
\eng{democratisation} & \eng{di-mo-kra-tai-'zei-sn} & Démocratisation (processus). \\
\hline
\end{tabularx}
\end{center}

\begin{aretenir}[Règle d'accentuation : suffixe \eng{-ic} / \eng{-ical}]
Les adjectifs en \eng{-ic} (\eng{democratic}, \eng{civic}, \eng{historic}, \eng{economic}) portent l'accent sur la syllabe \textbf{immédiatement précédant} le suffixe \eng{-ic}.
\begin{itemize}
    \item \eng{de-mo-\textbf{CRA}-tic} (pas \eng{de-MO-cra-tic}).
    \item \eng{ci-\textbf{VI}-c} (pas \eng{CI-vi-c}).
\end{itemize}
L'adverbe \eng{-ically} conserve cet accent : \eng{de-mo-\textbf{CRA}-ti-cal-ly}.
\end{aretenir}

\courssub{Famille de \eng{citizen} (citoyen)}

\begin{definition}[Famille lexicale : CITIZEN]
Attention à l'orthographe \eng{ti} = \eng{/ʃ/} (son « ch ») et à la distinction \eng{civic} / \eng{civics}.
\end{definition}

\begin{center}
\begin{tabularx}{\linewidth}{|X|X|X|}
\hline
\rowcolor{popblueL} \textbf{Forme} & \textbf{Prononciation (approx. fr.)} & \textbf{Traduction / Nuance} \\
\hline
\eng{citizen} & \eng{'si-ti-zn} (\textbf{si-ti-zeun}) & Citoyen. \eng{Senior citizen} (personne âgée). \\
\hline
\eng{citizenship} & \eng{'si-ti-zn-ship} & Citoyenneté (statut, droits, devoirs). \eng{Active citizenship}. \\
\hline
\eng{civic} & \eng{'si-vik} (\textbf{si-vik}) & Civique (adj : relatif au citoyen/ville). \eng{Civic duty}, \eng{civic centre}. \\
\hline
\eng{civics} & \eng{'si-viks} (\textbf{si-viks}) & \textbf{Éducation civique} (matière scolaire, toujours au pluriel en forme). \\
\hline
\eng{civilian} & \eng{si-'vil-ian} & Civil (non militaire). \eng{Civilian population}. \\
\hline
\eng{civil} & \eng{'si-vl} & Civil (adj : non militaire, poli, droit civil). \eng{Civil society}, \eng{civil rights}. \\
\hline
\end{tabularx}
\end{center}

\begin{attention}[Faux ami \eng{civics} $\neq$ \eng{civic} / \eng{civil}]
\begin{itemize}
    \item \eng{Civics} (avec \textbf{s} final) = \textbf{Éducation civique} (nom de matière). \eng{“Civics is on the timetable.”} (L'éducation civique est à l'emploi du temps.)
    \item \eng{Civic} (adj) = Civique. \eng{“Civic education”} = Éducation civique (adj + nom).
    \item \eng{Civil} (adj) = Civil (opposé à militaire) ou poli. \eng{“Civil society”} = Société civile. \eng{“Civil war”} = Guerre civile.
\end{itemize}
Ne dites pas \eng{“I study civic”} mais \eng{“I study civics”} ou \eng{“I study civic education”}.
\end{attention}

\courssub{Famille de \eng{govern} (gouverner)}

\begin{definition}[Famille lexicale : GOVERN]
Verbe régulier. Distinction cruciale entre \eng{government} (institution/équipe) et \eng{governance} (processus/manière de gouverner).
\end{definition}

\begin{center}
\begin{tabularx}{\linewidth}{|X|X|X|}
\hline
\rowcolor{popblueL} \textbf{Forme} & \textbf{Prononciation (approx. fr.)} & \textbf{Traduction / Contexte démocratique} \\
\hline
\eng{to govern} & \eng{'ge-vn} (\textbf{geu-veun}) & Gouvernor, diriger. \\
\hline
\eng{government} & \eng{'ge-vn-mnt} (\textbf{geu-veun-meunt}) & Gouvernement (exécutif), gouvernement (institution). \eng{Local government}. \\
\hline
\eng{governance} & \eng{'ge-v-nns} (\textbf{geu-veu-nanss}) & \textbf{Gouvernance} (qualité de la gestion, transparence, reddition de comptes). \eng{Good governance}. \\
\hline
\eng{governor} & \eng{'ge-v-n} (\textbf{geu-veu-neu}) & Gouverneur (chef d'État fédéré, ex: USA, Nigéria), directeur. \\
\hline
\eng{governing} & \eng{'ge-v-n-ing} & Adj : au pouvoir. \eng{The governing party}. Participe : \eng{Governing is hard}. \\
\hline
\eng{ungovernable} & \eng{un-'ge-v-n-bl} & Ingouvernable. \\
\hline
\end{tabularx}
\end{center}

\begin{savaistu}[\eng{Government} vs \eng{Governance} : la nuance académique]
En science politique et dans les formations sur la démocratie (PNUD, UE, ONG) :
\begin{itemize}
    \item \textbf{Government} = \textbf{Qui ?} L'équipe au pouvoir, les ministres, l'appareil d'État. C'est un \textbf{nom comptable} (\eng{a government}, \eng{governments}).
    \item \textbf{Governance} = \textbf{Comment ?} L'art de gouverner : transparence, responsabilité (\eng{accountability}), état de droit (\eng{rule of law}), participation. C'est un \textbf{nom non comptable}.
\end{itemize}
\eng{“A new \textbf{government} must improve \textbf{governance}.”} = « Un nouveau gouvernement doit améliorer la gouvernance. »
\end{savaistu}

\courssub{Prononciation : les pièges sonores du thème}

Le passage de l'écrit à l'oral est souvent source d'erreurs. Voici les prononciations cibles (notation francisée pour repérage rapide) et les règles associées.

\begin{center}
\begin{tabularx}{\linewidth}{|X|X|X|}
\hline
\rowcolor{popblueL} \textbf{Mot} & \textbf{Notation francisée} & \textbf{Point de vigilance} \\
\hline
\eng{workshop} & \textbf{‘ouèrk-chop} & \eng{w} = \eng{ouè} ; \eng{sh} = \eng{ch} ; \eng{o} court. Pas « wor-k-chop ». \\
\hline
\eng{seminar} & \textbf{‘sé-mi-nar} & Accent sur 1ère syllabe. \eng{a} final long \eng{ar}. Pas « sé-mi-nère ». \\
\hline
\eng{democracy} & \textbf{di-‘mo-kra-si} & Accent sur \textbf{mo}. \eng{o} court. \eng{cy} = \eng{si}. \\
\hline
\eng{citizen} & \textbf{‘si-ti-zeun} & \eng{ti} = \eng{si} (palatalisation). \eng{en} = \eng{zeun} (\eng{/zn/}). \\
\hline
\eng{governance} & \textbf{‘geu-veu-nanss} & \eng{go} = \eng{geu} (\eng{/gʌ/}) ; \eng{vernance} = \eng{veu-nanss}. Pas « gou-ve-ran-ce ». \\
\hline
\eng{awareness} & \textbf{a-‘wéa-nèss} & Accent sur \textbf{éa}. \eng{ware} se prononce \eng{wéa} (diphthongue). \\
\hline
\eng{accountability} & \textbf{a-kaoun-ta-bi-li-ti} & Accent principal sur \textbf{bi}. Secondaire sur \emph{kaoun}. \\
\hline
\eng{transparency} & \textbf{trans-‘pa-rèn-si} & Accent sur \textbf{pa}. \eng{par} = \eng{pa} (ouvert). \\
\hline
\end{tabularx}
\end{center}

\begin{methode}[Comment travailler sa prononciation en autonomie]
\begin{enumerate}
    \item \textbf{Écoutez} le mot dans un dictionnaire en ligne (Oxford, Cambridge, Collins) sans lire l'écrit.
    \item \textbf{Répétez} en exagérant l'accent tonique (tape dans les mains sur la syllabe forte).
    \item \textbf{Enregistrez}-vous sur votre téléphone et comparez.
    \item \textbf{Ancrez} le mot dans une phrase courte : \eng{“Good \textbf{governance} is key.”} (Répétez la phrase 3 fois).
    \item \textbf{Vérifiez} l'orthographe des sons difficiles (\eng{ti/ci/si} = \eng{/ʃ/} ; \eng{aw} = \eng{/ɔː/} ou \eng{/æ/}).
\end{enumerate}
\end{methode}

\courssub{Orthographe : doubles consonnes, suffixes et pièges écrits}

L'orthographe anglaise est morphologique (elle marque l'histoire du mot) plus que phonétique. Trois règles couvrent 90\% des erreurs sur ce thème.

\begin{propriete}[Règle 1 : Le doublement de la consonne finale (CVC + suffixe vocalique)]
Un verbe monosyllabe (ou accentué sur la dernière syllabe) se terminant par \textbf{Consonne - Voyelle - Consonne} double la consonne finale avant un suffixe commençant par une voyelle (\eng{-ing, -ed, -er, -est}).
\begin{center}
\begin{tabular}{|l|l|l|l|}
\hline
\rowcolor{popblueL} \textbf{Base (CVC)} & \textbf{+ ing} & \textbf{+ ed} & \textbf{+ er} \\
\hline
\eng{plan} & \eng{planning} & \eng{planned} & \eng{planner} \\
\eng{refer} (accent fin) & \eng{referring} & \eng{referred} & \eng{referrer} \\
\eng{occur} & \eng{occurring} & \eng{occurred} & — \\
\hline
\end{tabular}
\end{center}
\emph{Application au thème :} \eng{to plan} $\rightarrow$ \eng{a planning meeting} (réunion de planification). \eng{To refer} $\rightarrow$ \eng{a referral} (renvoi).
\medskip
\textbf{Attention :} \eng{Train} se termine par \textbf{deux consonnes} (\eng{in}). \textbf{Pas de doublement} : \eng{training} (un seul \eng{n}), \eng{trained}, \eng{trainer}. Écrire \eng{trainning} est une \textbf{faute grave}.
\end{propriete}

\begin{propriete}[Règle 2 : Le \eng{-e} muet final tombe devant suffixe vocalique]
\eng{Participate} $\rightarrow$ \eng{participating}, \eng{participation} (pas \eng{participateing}).
\eng{Govern} $\rightarrow$ \eng{governing}, \eng{governance} (pas \eng{governeance}).
\eng{Democratise} $\rightarrow$ \eng{democratising}, \eng{democratisation}.
\emph{Exception :} \eng{Changeable, courageous, noticeable} (gardent le \eng{e} pour garder le son \eng{/dʒ/} ou \eng{/s/} doux). Pas de cas dans notre liste.
\end{propriete}

\begin{propriete}[Règle 3 : Suffixes \eng{-ance/-ence}, \eng{-ant/-ent}, \eng{-ancy/-ency}]
Pas de règle simple universelle, mais des familles :
\begin{itemize}
    \item Verbes en \eng{-ate} $\rightarrow$ \eng{-ance/-ant} : \eng{participate $\rightarrow$ participation, participant}. \eng{Tolerate $\rightarrow$ tolerance, tolerant}.
    \item Verbes en \eng{-ere/-er} $\rightarrow$ \eng{-ance/-ant} : \eng{appear $\rightarrow$ appearance}.
    \item \eng{Govern} $\rightarrow$ \eng{Governance} (son \eng{/nəns/}), pas \eng{government} (son \eng{/nmənt/}).
    \item \eng{Responsible} (adj, \eng{-ible}) $\rightarrow$ \eng{Responsibility} (nom, \eng{-ility}). \textbf{Pas de \eng{a} après le \eng{s}}.
\end{itemize}
\end{propriete}

\begin{exoresolu}[Correction d'erreurs fréquentes (niveau Bac)]
\textbf{Corrigez les erreurs dans les phrases suivantes (orthographe, forme du mot, préposition).}
\begin{enumerate}
    \item The \underline{trainner} explained the program to the \underline{trainees}.
    \item We need more \underline{participents} for the \underline{workshop} on \underline{democratic} governance.
    \item \underline{Civics} education teaches \underline{citizenship} duties.
    \item The \underline{governement} must ensure \underline{good governence}.
    \item She wants to \underline{participate to} the \underline{training} session.
    \item It is a \underline{responsability} of every \underline{citizen}.
\end{enumerate}

\tcblower
\textbf{Solution détaillée :}
\begin{enumerate}
    \item \textbf{trainer} (un seul \eng{n}, agent en \eng{-er}) / \textbf{trainees} (correct, patient en \eng{-ee}, accent sur \eng{ee}).
    \item \textbf{participants} (suffixe \eng{-ant}, pas \eng{-ent}) / \textbf{workshop} (correct) / \textbf{democratic} (adj, correct).
    \item \textbf{Civic education} ou \textbf{Civics} (seul). \eng{Civics education} est un pléonasme (éducation éducation civique). On dit \eng{Civics} (matière) ou \eng{Civic education}. / \textbf{citizenship} (correct).
    \item \textbf{government} (\eng{n} avant \eng{m}, son \eng{/nmənt/}) / \textbf{governance} (\eng{ance}, son \eng{/nəns/}).
    \item \textbf{participate \textbf{in}} (verbe + \eng{in}, jamais \eng{to}) / \textbf{training} (un seul \eng{n}).
    \item \textbf{responsibility} (\eng{-ility}, pas \eng{-ability}) / \textbf{citizen} (correct).
\end{enumerate}
\end{exoresolu}

\courssub{Méthode : Enrichir sa production écrite par la variation morphologique}

Dans une composition (\eng{essay}), une lettre formelle ou un article, la répétition du même mot (\eng{participate, participate, participate}) pénalise la note de \textbf{vocabulaire} et de \textbf{cohérence}. La dérivation est l'outil n°1 pour varier.

\begin{methode}[Technique du \eng{Word Family Substitution} (Remplacement par la famille)]
\textbf{Objectif :} Transformer une phrase répétitive en phrase variée sans changer le sens.
\textbf{Étapes :}
\begin{enumerate}
    \item \textbf{Identifiez} le mot-clé répété (souvent un verbe : \eng{participate, train, govern, develop}).
    \item \textbf{Choisissez} un autre membre de la famille (nom, adjectif, adverbe) selon la structure syntaxique manquante.
    \item \textbf{Reconstruisez} la phrase (souvent : Verbe + Objet $\rightarrow$ Adjectif + Nom ; ou Verbe $\rightarrow$ Nom + Verbe support \eng{make/do/have/take}).
\end{enumerate}
\end{methode}

\begin{exemplebox}[Application : \eng{Participate}]
\textbf{Phrase faible (répétition) :}
\eng{“Citizens must \textbf{participate} in elections. If they \textbf{participate} actively, democracy works. Without \textbf{participate}, the system fails.”} (Erreur : \eng{participate} nom impossible).

\textbf{Phrase forte (variation morphologique) :}
\eng{“Citizens must \textbf{participate} in elections. \textbf{Active participation} strengthens democracy. Without \textbf{citizen participation}, the system fails.”}
\begin{itemize}
    \item Occurrence 1 : Verbe \eng{participate}.
    \item Occurrence 2 : Nom \eng{participation} + Adj \eng{active} (collocation forte).
    \item Occurrence 3 : Nom composé \eng{citizen participation} (précision sémantique).
\end{itemize}
\end{exemplebox}

\begin{exemplebox}[Application : \eng{Train / Govern / Aware}]
\begin{tabularx}{\linewidth}{|X|X|}
\hline
\rowcolor{popblueL} \textbf{Base verbale} & \textbf{Variations nominales / adjectivales pour l'écrit} \\
\hline
\eng{to train} & \eng{provide \textbf{training}}, \eng{a \textbf{trainer}'s guide}, \eng{\textbf{well-trained} staff}, \eng{\textbf{retraining} programmes}. \\
\hline
\eng{to govern} & \eng{\textbf{government} policy}, \eng{\textbf{governance} principles}, \eng{a \textbf{governing} body}, \eng{\textbf{ungovernable} regions}. \\
\hline
\eng{to be aware} & \eng{raise \textbf{awareness}}, \eng{\textbf{awareness-raising} campaign}, \eng{\textbf{aware} of the risks}, \eng{\textbf{unaware} citizens}. \\
\hline
\end{tabularx}
\end{exemplebox}

\courssub{Exercices d'application (Entraînement Bac)}

\begin{exercice}[Morphologie et Orthographe : Complétion guidée]
Complétez le tableau en donnant les formes manquantes. Attention à l'orthographe et à la catégorie grammaticale.
\begin{center}
\begin{tabular}{|l|l|l|l|l|}
\hline
\rowcolor{popblueL} \textbf{Verbe} & \textbf{Nom (Action/État)} & \textbf{Nom (Agent/Patient)} & \textbf{Adjectif} & \textbf{Adverbe} \\
\hline
\eng{participate} & \eng{participation} & \eng{participant} & \eng{participatory} & \eng{participatorily} (rare) \\
\hline
\eng{democratise} & \eng{\ldots} & \eng{democrat} & \eng{\ldots} & \eng{\ldots} \\
\hline
\eng{\ldots} & \eng{citizenship} & \eng{citizen} & \eng{civic} & \eng{civically} \\
\hline
\eng{govern} & \eng{government / governance} & \eng{governor} & \eng{governmental} & \eng{governmentally} \\
\hline
\eng{train} & \eng{training} & \eng{trainer / trainee} & \eng{trained / training (adj)} & — \\
\hline
\eng{be aware} & \eng{awareness} & — & \eng{aware / unaware} & \eng{awarely} (rare) \\
\hline
\end{tabular}
\end{center}
\textbf{Consigne :} Remplissez les cases pointillées (\eng{\ldots}).
\end{exercice}

\begin{exercice}[Traduction thématique : Du français vers l'anglais]
Traduisez en anglais en utilisant \textbf{impérativement} le mot de la famille indiqué entre parenthèses.
\begin{enumerate}
    \item La \textbf{formation} des formateurs est une priorité. (\eng{TRAIN} $\rightarrow$ nom action)
    \item Les \textbf{stagiaires} ont apprécié la méthode participative. (\eng{TRAIN} $\rightarrow$ nom patient)
    \item Nous devons \textbf{démocratiser} l'accès à l'information. (\eng{DEMOCRACY} $\rightarrow$ verbe)
    \item C'est une décision prise \textbf{démocratiquement}. (\eng{DEMOCRACY} $\rightarrow$ adverbe)
    \item L'\textbf{éducation civique} est au programme. (\eng{CITIZEN} $\rightarrow$ nom matière scolaire)
    \item Les \textbf{citoyens} ont des devoirs civiques. (\eng{CITIZEN} $\rightarrow$ nom agent)
    \item La \textbf{gouvernance} locale s'améliore. (\eng{GOVERN} $\rightarrow$ nom processus)
    \item Le \textbf{gouverneur} de la région a inauguré le centre. (\eng{GOVERN} $\rightarrow$ nom agent)
    \item La \textbf{sensibilisation} (\eng{awareness}) des jeunes est cruciale. (\eng{AWARE} $\rightarrow$ nom)
    \item Il est \textbf{responsable} de ses actes. (\eng{RESPONSIBLE} $\rightarrow$ adjectif) / Sa \textbf{responsabilité} est engagée. (\eng{RESPONSIBLE} $\rightarrow$ nom en \eng{-ility})
\end{enumerate}
\end{exercice}


\begin{corrige}[Exercice 1 : Tableau complété]
\begin{center}
\begin{tabular}{|l|l|l|l|l|}
\hline
\rowcolor{popblueL} \textbf{Verbe} & \textbf{Nom (Action/État)} & \textbf{Nom (Agent/Patient)} & \textbf{Adjectif} & \textbf{Adverbe} \\
\hline
\eng{participate} & \eng{participation} & \eng{participant} & \eng{participatory} & \eng{participatorily} \\
\hline
\eng{democratise} & \eng{democratisation} & \eng{democrat} & \eng{democratic} & \eng{democratically} \\
\hline
\eng{— (pas de verbe direct)} & \eng{citizenship} & \eng{citizen} & \eng{civic} & \eng{civically} \\
\hline
\eng{govern} & \eng{government / governance} & \eng{governor} & \eng{governmental} & \eng{governmentally} \\
\hline
\eng{train} & \eng{training} & \eng{trainer / trainee} & \eng{trained / training} & — \\
\hline
\eng{be aware} & \eng{awareness} & — & \eng{aware / unaware} & \eng{awarely} \\
\hline
\end{tabular}
\end{center}
\end{corrige}

\begin{corrige}[Exercice 2 : Traduction]
\begin{enumerate}
    \item \textbf{Training} of trainers is a priority. (\eng{Trainer training} aussi possible).
    \item The \textbf{trainees} appreciated the \textbf{participatory} method.
    \item We must \textbf{democratise} access to information.
    \item It is a \textbf{democratically} made decision. / The decision was made \textbf{democratically}.
    \item \textbf{Civics} is on the curriculum. / \textbf{Civic education} is on the curriculum.
    \item \textbf{Citizens} have \textbf{civic} duties.
    \item Local \textbf{governance} is improving.
    \item The \textbf{governor} of the region inaugurated the centre.
    \item Raising \textbf{awareness} among youth is crucial. / Youth \textbf{awareness} is crucial.
    \item He is \textbf{responsible} for his actions. / His \textbf{responsibility} is engaged.
\end{enumerate}
\end{corrige}

\begin{aretenir}[Check-list de révision pour le Bac (Lexique Démocratie/Formation)]
\begin{itemize}
    \item \textbf{Orthographe :} \eng{Training} (1 \eng{n}), \eng{Participant} (\eng{-ant}), \eng{Government} (\eng{n} avant \eng{m}), \eng{Governance} (\eng{ance}), \eng{Responsibility} (\eng{-ility}), \eng{Democratically} (\eng{-ally}).
    \item \textbf{Prononciation :} \eng{Trainee} (accent fin), \eng{Democratic} (accent \eng{CRA}), \eng{Citizen} (\eng{ti}=si), \eng{Governance} (\eng{geu-veu-nanss}), \eng{Workshop} (\eng{ouèrk-chop}).
    \item \textbf{Grammaire :} \eng{Participate \textbf{in}}, \eng{Aware \textbf{of}}, \eng{Responsible \textbf{for}}, \eng{Train \textbf{in} / \textbf{for}}.
    \item \textbf{Distinction :} \eng{Government} (qui) vs \eng{Governance} (comment) ; \eng{Trainer} (fait) vs \eng{Trainee} (subit) ; \eng{Civics} (matière) vs \eng{Civic} (adj).
    \item \textbf{Variation :} Savoir passer de \eng{participate} $\rightarrow$ \eng{participation} $\rightarrow$ \eng{participatory} pour enrichir sa copie.
\end{itemize}
\end{aretenir}

\coursec{Le lexique en contexte : dialogue et production guidée}

Dans la section précédente, nous avons étudié les \cle{familles de mots} (\eng{train, trainer, trainee, training}), la prononciation et l'orthographe du lexique des formations sur la démocratie. Mais un mot n'est vraiment acquis que lorsqu'on sait l'\emph{employer} dans une situation réelle. C'est l'objectif de cette dernière section : passer du dictionnaire à la communication, à l'oral comme à l'écrit, avec les automatismes et la précision exigés au Baccalauréat.

\courssub{Dialogue modèle : talking about a workshop}

Lisons d'abord ce dialogue entre deux élèves de Tle A, le lundi matin, dans la cour du lycée.

\begin{textebox}[Dialogue — After the workshop]
\textbf{A:} Did you attend the workshop on democracy last Saturday?\\
\textbf{B:} Yes, I did. It was organised by a local NGO and facilitated by two excellent trainers.\\
\textbf{A:} Where did it take place?\\
\textbf{B:} At the community hall, near the main market. About fifty participants took part, mostly students and youth leaders.\\
\textbf{A:} What did you learn?\\
\textbf{B:} We discussed human rights and how young people can make their voices heard. The facilitators also raised our awareness of the importance of voting.\\
\textbf{A:} That sounds really useful!\\
\textbf{B:} It was. At the end, we received a certificate of participation.
\end{textebox}

\textbf{Glossaire :} \eng{NGO} (n., \emph{non-governmental organisation}, ONG) ; \eng{facilitated by} (animée par) ; \eng{youth leaders} (n. pl., responsables de jeunesse) ; \eng{make their voices heard} (faire entendre leur voix) ; \eng{raised our awareness of} (nous a sensibilisés à) ; \eng{certificate of participation} (n., certificat de participation).

\courssub{Observer : les questions qui structurent le dialogue}

Pour parler d'une formation, l'anglais utilise trois questions-types qu'il faut mémoriser avec leurs réponses :

\begin{center}
\begin{tabularx}{\linewidth}{|X|X|X|}
\hline
\rowcolor{popblueL} Question & Réponse type & Remarque \\
\hline
\eng{Where did it take place?} & \eng{It took place at the community hall.} & \eng{take place} + lieu (\eng{at} lieu précis, \eng{in} ville/pays/mois) \\
\hline
\eng{Who organised it?} & \eng{It was organised by a local NGO.} & passif avec \eng{by} + agent \\
\hline
\eng{What did you learn?} & \eng{We learned about human rights.} & \eng{learn about} + thème \\
\hline
\eng{Who facilitated the sessions?} & \eng{Two trainers facilitated them.} & \eng{facilitate} = animer \\
\hline
\end{tabularx}
\end{center}

\begin{attention}[Le passif, roi du compte rendu]
Pour raconter une formation, l'anglais privilégie le \cle{passif} : \eng{The workshop \textbf{was organised} by an NGO}, \eng{participants \textbf{were trained} in civic education}, \eng{certificates \textbf{were awarded}}. Le francophone traduit souvent mot à mot par l'actif (\eng{An NGO organised the workshop}), ce qui est correct mais moins naturel dans un compte rendu. Entraîne-toi à transformer : \eng{They organised a seminar} $\rightarrow$ \eng{A seminar was organised}.
\end{attention}

\courssub{Réemployer le lexique : phrases modèles à trous}

Voici les cinq collocations-clés de la leçon, à réemployer telles quelles :

\begin{itemize}
\item \eng{\textbf{attend} a workshop / a course / a seminar} — assister à (jamais \eng{assist to} !)
\item \eng{the \textbf{facilitator} guided the discussion} — l'animateur a guidé le débat
\item \eng{\textbf{raise awareness of} human rights} — sensibiliser aux droits de l'homme
\item \eng{receive a \textbf{certificate} (of participation)} — recevoir un certificat
\item \eng{\textbf{take part in} a training session} — participer à une session de formation
\end{itemize}

\begin{exoresolu}[Phrases à trous]
Complète avec : \eng{attend, facilitator, raise awareness, certificate, take part}.

1. Last month, I decided to \ldots\ a training course on citizenship.\\
2. The \ldots\ asked us many questions about voting.\\
3. The campaign aims to \ldots\ \ldots\ of children's rights.\\
4. All students who \ldots\ \ldots\ in the workshop will receive a \ldots .

\tcblower
\textbf{Solution :}
1. \eng{attend} — assister à une formation (pas de préposition après \eng{attend}).\\
2. \eng{facilitator} — celui qui anime la session pose les questions.\\
3. \eng{raise awareness} — collocation figée : \eng{raise awareness of + nom}.\\
4. \eng{take part} (collocation \eng{take part in}) ; \eng{certificate} (on \emph{reçoit} un certificat : \eng{receive a certificate}).
\end{exoresolu}

\courssub{Texte de lecture : a training course in Bamenda}

\begin{textebox}[Reading — Civic education at GBHS Bamenda]
Last Friday, Government Bilingual High School Bamenda hosted a one-day training course on democracy and civic responsibility. The event was organised by the school's English club, in partnership with a local NGO called Youth for Peace. About eighty students from Form Five to Upper Sixth took part in the sessions.

The morning workshop was facilitated by Mrs Ndi, a lawyer, and Mr Tabi, a journalist. They first explained the basic principles of democracy: free elections, the separation of powers and the rule of law. Then the participants were divided into small groups to discuss real-life situations, such as corruption in public services and discrimination at school. Each group had to propose concrete solutions and present them to the audience.

In the afternoon, the trainers focused on the role of young people in public life. They raised the students' awareness of their right to express their opinions peacefully and of their duty to respect other people's rights. A mock election was even organised to show how a transparent vote works.

At the end of the day, every participant received a certificate of attendance from the organisers. “I learned that democracy is not only about voting,” said Larissa, a Lower Sixth student. “It is also about listening, debating and respecting others.” The head teacher announced that similar workshops would be held every term.
\end{textebox}

\textbf{Glossaire :} \eng{hosted} (v., a accueilli / organiser) ; \eng{in partnership with} (loc. prép., en partenariat avec) ; \eng{the rule of law} (n., l'État de droit) ; \eng{real-life situations} (n. pl., situations de la vie réelle) ; \eng{mock election} (n., élection simulée) ; \eng{certificate of attendance} (n., certificat de présence) ; \eng{every term} (loc. adv., chaque trimestre).

\textbf{Questions de compréhension :}

\textbf{A. Vrai ou faux ?} Justifie par une citation du texte.

1. The training course lasted three days.\\
2. Mrs Ndi is a journalist.\\
3. A mock election was organised in the afternoon.

\textbf{B. QCM.} Choisis la bonne réponse.

4. The event was organised by : (a) the Ministry of Education (b) the English club and an NGO (c) the head teacher alone.\\
5. In groups, students discussed : (a) football results (b) real-life problems like corruption (c) their holidays.

\textbf{C. Réponses courtes.}

6. Who facilitated the morning workshop?\\
7. What did every participant receive at the end of the day?\\
8. According to Larissa, what is democracy about?

\begin{corrige}[Exercice de compréhension]
1. \textbf{False} — \eng{“a one-day training course”}. 2. \textbf{False} — \eng{“Mrs Ndi, a lawyer”} ; c'est Mr Tabi qui est journaliste. 3. \textbf{True} — \eng{“A mock election was even organised”} (après-midi). 4. \textbf{(b)}. 5. \textbf{(b)}. 6. \eng{Mrs Ndi, a lawyer, and Mr Tabi, a journalist.} 7. \eng{A certificate of attendance.} 8. \eng{Listening, debating and respecting others — not only voting.}
\end{corrige}

\begin{attention}[Accorde les temps dans un compte rendu au passé]
Un compte rendu de formation se raconte au \cle{preterit}. Ne mélange pas les temps : \eng{The workshop \textbf{was organised} by an NGO, participants \textbf{discussed} human rights, and we \textbf{received} a certificate.} Erreur fréquente du francophone : repasser au présent en cours de récit (\eng{The workshop was organised and the facilitator explains\ldots}) — faute sanctionnée au Bac.
\end{attention}

\courssub{Production guidée : “A training course I would like to attend”}

\begin{methode}[Rédiger ton paragraphe en 3 étapes]
\begin{enumerate}
\item \textbf{Name the course} : nomme la formation (type, lieu, organisateur) — \eng{I would like to attend a workshop on\ldots, organised by\ldots, in Yaoundé.}
\item \textbf{Say what you would learn} : indique les thèmes abordés — \eng{The course covers good governance, human rights and\ldots}
\item \textbf{Explain why} : explique ta motivation (valeurs, projets) — \eng{because I want to serve my community / become a youth leader.}
\end{enumerate}
Puis relis-toi : faux amis (\eng{formation} (FR) $\neq$ \eng{formation} (EN, géologie/militaire) $\rightarrow$ \eng{training course}, \eng{workshop}), collocations (\eng{attend} sans préposition, \eng{take part in}), temps (conditionnel \eng{would like} pour un souhait).
\end{methode}

\begin{exemplebox}[Phrases modèles traduites]
\eng{I would like to enrol in a training course on good governance because I want to serve my community.} « Je voudrais m'inscrire à une formation sur la bonne gouvernance parce que je veux servir ma communauté. »

\eng{Such a course would teach me how elections work and how citizens can hold leaders accountable.} « Une telle formation m'apprendrait comment fonctionnent les élections et comment les citoyens peuvent demander des comptes aux dirigeants. »

\eng{At the end of the training, I hope to receive a certificate and share my knowledge with my classmates.} « À la fin de la formation, j'espère recevoir un certificat et partager mes connaissances avec mes camarades. »
\end{exemplebox}

\begin{popfigure}
\begin{tikzpicture}[node distance=0.5cm]
\node[draw=popblue, fill=popblueL, rounded corners, minimum width=3.2cm, minimum height=1.6cm, align=center] (plan) {\textbf{PLAN}\\ idées : type, lieu,\\ thèmes, motivation};
\node[draw=poporange, fill=poporangeL, rounded corners, minimum width=3.2cm, minimum height=1.6cm, align=center, right=1.1cm of plan] (draft) {\textbf{DRAFT}\\ brouillon :\\ 6--8 lignes,\\ temps corrects};
\node[draw=popgreen, fill=popgreenL, rounded corners, minimum width=3.2cm, minimum height=1.6cm, align=center, right=1.1cm of draft] (check) {\textbf{CHECK}\\ relecture : faux amis,\\ collocations,\\ orthographe};
\draw[-{Stealth}, very thick, popdark] (plan) -- (draft);
\draw[-{Stealth}, very thick, popdark] (draft) -- (check);
\end{tikzpicture}
\legende{Les trois étapes de la production guidée : planifier, rédiger, relire.}
\end{popfigure}

\begin{exercice}[À toi d'écrire]
Rédige un paragraphe de 6 à 8 lignes intitulé \eng{A training course I would like to attend}. Utilise au moins quatre des mots suivants : \eng{attend, enrol, facilitator, raise awareness, take part in, certificate, good governance, human rights}.
\end{exercice}

\begin{savaistu}[Complément Bac]
Au Baccalauréat A, un sujet d'expression écrite peut te demander de décrire une activité citoyenne de ton école (club, campagne de sensibilisation, élection des délégués) ou de rédiger un \emph{essay} sur \eng{citizenship}, \eng{youth and democracy}. Tout le lexique de cette leçon — \eng{workshop, facilitator, raise awareness, take part, certificate, good governance, make one's voice heard} — est directement réutilisable et te rapporte des points de vocabulaire. Garde cette liste dans ton carnet de révisions !
\end{savaistu}

\newpage

\coursec{Activité d'intégration}

\courssub{Objectifs de l'activité}

À l'issue de cette activité d'intégration, l'élève sera capable de :
\begin{itemize}
\item réemployer le lexique des formations et cours sur la démocratie (\eng{training course}, \eng{workshop}, \eng{seminar}, \eng{facilitator}, \eng{raise awareness}, \eng{empower}, \eng{certificate}, \eng{take part in}...) dans une tâche authentique ;
\item comprendre un document authentique (une annonce d'appel à candidatures) et en extraire les informations utiles ;
\item rédiger une lettre de motivation courte, correcte et bien structurée en anglais ;
\item éviter les faux amis étudiés dans la leçon (\eng{to attend} et non \eng{to assist}, \eng{training} et non \eng{formation}, \eng{actually} et non \eng{actuellement}...) ;
\item s'exprimer oralement pendant deux minutes avec une prononciation soignée et une bonne fluidité ;
\item évaluer la production de ses camarades à l'aide d'une grille simple.
\end{itemize}

\courssub{Situation}

Une ONG camerounaise, \emph{Youth for Democracy Cameroon} (YDC), organise à Yaoundé un camp national de formation des jeunes à la citoyenneté. Elle publie un appel à candidatures dans les lycées. Votre classe a décidé de répondre à cet appel : chaque binôme préparera un dossier complet et passera un entretien de sélection devant un jury.

\begin{textebox}[Call for applications — Youth for Democracy Cameroon]
\textbf{NATIONAL YOUTH TRAINING CAMP ON DEMOCRACY AND HUMAN RIGHTS}\\
\textbf{Yaoundé, 15--22 August}

Youth for Democracy Cameroon (YDC) is pleased to announce its annual training camp for young citizens aged 17 to 20. Are you passionate about democracy, human rights and community life? Then this camp is for you!

\textbf{Programme:} interactive workshops on voting and elections, a seminar on human rights with experienced facilitators, group discussions on gender equality, and a final debate on youth participation in public life.

\textbf{Our objectives:} to raise awareness of citizens' rights and duties, to empower young people to take part in decision-making, and to train future community leaders.

\textbf{Conditions:} participants must attend all sessions. A certificate of participation will be awarded at the end of the camp. Accommodation and meals are free.

\textbf{How to apply:} fill in the application form below and write a short motivation letter (10--12 lines) explaining why you want to join the camp. Shortlisted candidates will be invited to a two-minute interview.

\textbf{Deadline:} 30 June. Send your application to: YDC, P.O. Box 1234, Yaoundé.
\end{textebox}

\textbf{Glossaire :} \emph{to be pleased} (être heureux de), \emph{aged} (âgé de), \emph{accommodation} (le logement), \emph{shortlisted} (présélectionné), \emph{deadline} (la date limite), \emph{P.O. Box} (boîte postale).

\courssub{Consignes}

Travail par binômes. Suivez les étapes dans l'ordre.

\begin{enumerate}
\item \textbf{Compréhension du document.} Lisez l'appel à candidatures et répondez en anglais : (a) Who organises the camp? (b) Where and when does it take place? (c) Name two activities offered. (d) What will participants receive at the end? (e) What must candidates send before 30 June?
\item \textbf{Fiche d'inscription.} Remplissez en anglais la fiche suivante pour les deux membres du binôme : \emph{Full name / Age / School and class / Town / Course you wish to attend (workshop, seminar or debate) / Why this course? (one sentence)}. Attention à l'ordre des mots et aux majuscules des noms propres.
\item \textbf{Lettre de motivation.} Rédigez ensemble une lettre de motivation de 10 à 12 lignes. Vous devez utiliser \textbf{au moins 8 mots ou expressions de la leçon}, par exemple : \eng{attend}, \eng{training course}, \eng{facilitator}, \eng{raise awareness}, \eng{human rights}, \eng{certificate}, \eng{take part in}, \eng{empower}, \eng{workshop}, \eng{citizenship}. Soulignez-les dans votre lettre. Vérifiez qu'aucun faux ami ne s'est glissé dans votre texte.
\item \textbf{Relecture croisée.} Échangez votre lettre avec un autre binôme. Traquez les erreurs : faux amis, ordre des mots, oubli du \eng{-s} à la troisième personne, prépositions (\eng{take part \textbf{in}}, \eng{interested \textbf{in}}, \eng{apply \textbf{for}}).
\item \textbf{Entretien oral.} Chaque binôme passe devant le jury un mini-entretien de deux minutes : chaque élève parle environ une minute pour présenter sa motivation. Préparez-vous à répondre à des questions comme : \eng{Why do you want to attend this training camp?} ou \eng{What will you do with your certificate?}
\item \textbf{Évaluation.} Le jury (trois élèves + le professeur) note chaque binôme selon la grille : lexique correct et varié (4 pts), absence de faux amis (2 pts), prononciation (2 pts), fluidité et aisance (2 pts). Total : 10 pts.
\end{enumerate}

\courssub{Ressources à mobiliser}

\begin{aretenir}[Boîte à outils pour réussir]
\begin{itemize}
\item \textbf{Lexique des formations :} \eng{a training course} (une formation), \eng{a workshop} (un atelier), \eng{a seminar} (un séminaire), \eng{a facilitator} (un animateur/facilitateur), \eng{a trainee} (un stagiaire), \eng{a certificate} (un certificat).
\item \textbf{Expressions clés :} \eng{to raise awareness of} (sensibiliser à), \eng{to empower young people} (donner aux jeunes les moyens d'agir), \eng{to take part in} (participer à), \eng{to apply for} (poser sa candidature à), \eng{to attend a course} (suivre une formation).
\item \textbf{Faux amis à éviter :} \eng{to attend} = assister à (et non « attendre » = \eng{to wait for}) ; \eng{training} = formation (\eng{formation} en anglais = formation géologique ou mise en forme) ; \eng{actually} = en fait (et non « actuellement » = \eng{currently}) ; \eng{a lecture} = une conférence (et non « une lecture »).
\item \textbf{Structures de la lettre de motivation :} \eng{Dear Sir or Madam,} / \eng{I am writing to apply for...} / \eng{I am very interested in...} / \eng{I would be grateful if you could...} / \eng{Yours faithfully,}.
\item \textbf{Prononciation :} \eng{facilitator} « fé-si-li-té-teur », \eng{awareness} « a-wère-nèss », \eng{certificate} « seur-ti-fi-kète », \eng{democracy} « dé-mo-kra-si » (accent sur la 2\up{e} syllabe).
\end{itemize}
\end{aretenir}

\courssub{Proposition de production et corrigé commenté}

\begin{exoresolu}[Modèle de lettre de motivation]
Rédigez une lettre de motivation de 10 à 12 lignes pour le camp de YDC, avec au moins 8 mots de la leçon soulignés.
\tcblower
\textbf{Production modèle :}

\eng{Dear Sir or Madam,}

\eng{I am writing to \textbf{apply for} the National Youth \textbf{Training Camp} on Democracy and \textbf{Human Rights} in Yaoundé. I am an eighteen-year-old student in Terminale A at Lycée de Nkol-Eton, and I am deeply interested in \textbf{citizenship} and community life.}

\eng{I would like to \textbf{attend} your \textbf{training course} because I believe young people must \textbf{take part in} public life. The \textbf{workshops} on voting and the seminar with experienced \textbf{facilitators} will help me understand my rights and duties as a citizen. After the camp, I want to \textbf{raise awareness} of human rights among the students of my school and \textbf{empower} them to become active citizens.}

\eng{I am ready to attend all the sessions, and I would be very proud to receive a \textbf{certificate} at the end of the camp.}

\eng{Thank you for considering my application.}

\eng{Yours faithfully,}\\
\eng{Manga Etoundi Amina}

\textbf{Commentaire des choix de langue :}
\begin{itemize}
\item La lettre suit la structure classique : formule d'appel, objet, présentation, motivation, projet, formule de politesse. \eng{Yours faithfully} s'emploie quand on ne connaît pas le nom du destinataire.
\item Dix mots de la leçon sont réemployés (en gras), dont les collocations exactes : \eng{apply \textbf{for}}, \eng{take part \textbf{in}}, \eng{raise awareness \textbf{of}}.
\item Aucun faux ami : \eng{attend} est bien utilisé pour « suivre une formation », \eng{training course} remplace le calque \eng{formation course}.
\item Grammaire : présent simple pour les faits et goûts (\eng{I believe}, \eng{I am interested}), futur avec \eng{will} pour les conséquences prévues, \eng{want to + base verbale}.
\item Le projet concret (sensibiliser les élèves de son lycée) rend la lettre crédible et personnelle : c'est ce qu'attend un jury.
\end{itemize}
\end{exoresolu}

\courssub{Bilan de l'activité}

Cette activité a permis de réinvestir l'ensemble des savoir-faire de la leçon dans une situation de communication authentique : comprendre un appel à candidatures, mobiliser le lexique des formations sur la démocratie à l'écrit (fiche d'inscription, lettre de motivation) et à l'oral (entretien), et surveiller les faux amis qui piègent les francophones. L'évaluation par les pairs a en outre développé l'écoute critique et l'esprit citoyen : sélectionner un candidat, c'est déjà exercer une responsabilité démocratique. Pour aller plus loin, la classe peut afficher les meilleures lettres au mur de la salle ou les envoyer réellement à une association locale de promotion de la citoyenneté.

\newpage

\coursec{Méthodes et exercices résolus}

\coursec{Lesson 42 — Vocabulary: Words and expressions related to training activities/courses on democracy}

\courssub{Méthodes et exercices résolus}

\begin{methode}[Identifier et classifier les types de formations démocratiques]
\textbf{Objectif :} Reconnaître le vocabulaire des formats de formation (atelier, séminaire, session, module, etc.) et les classer selon leur durée, leur public ou leur modalité (présentiel / en ligne).\par
\textbf{Démarche en 4 étapes :}
\begin{enumerate}
    \item \textbf{Repérer les indices contextuels :} Mots-clés de durée (\eng{week-long, two-day, intensive}), de modalité (\eng{online, in-person, hybrid, virtual}), de public (\eng{grassroots, youth leaders, civil society, officials}).
    \item \textbf{Isoler le nom du format :} \eng{workshop, seminar, training session, course, module, masterclass, boot camp, webinar, round table, focus group}.
    \item \textbf{Associer la collocation verbale :} \eng{attend a workshop, conduct a seminar, facilitate a session, run a training, deliver a module, participate in a webinar}.
    \item \textbf{Classer dans un tableau mental :} Durée courte/longue -- Interactif/Magistral -- Présentiel/Digital.
\end{enumerate}
\textbf{Structures à retenir :} \eng{to organise/hold/conduct a + [format]} ; \eng{to take part in / attend / enrol in a + [format]} ; \eng{a + [adjectif] + [format] on + [thème]}.
\end{methode}

\begin{exoresolu}[Compréhension et classification de formats de formation]
\textbf{Énoncé :} Lisez l'annonce ci-dessous et répondez aux questions en anglais.\par
\begin{textebox}[Appel à candidatures -- ONG "Youth for Democracy"]
\eng{The “Youth for Democracy” initiative announces a \textbf{five-day residential workshop} on “Electoral Observation and Civic Engagement” for young leaders aged 18--30 from the Northwest and Southwest Regions. The \textbf{intensive training}, scheduled for 15--19 July 2025 in Buea, combines \textbf{plenary sessions} with \textbf{breakout groups} and \textbf{simulation exercises}. Selected participants will benefit from \textbf{peer-to-peer learning} and \textbf{mentorship} by accredited \textbf{facilitators} from the Electoral Board (ELECAM). A \textbf{follow-up webinar} is planned for October. \textbf{Application deadline: 30 June.}}
\end{textebox}
\begin{enumerate}
    \item Identify \textbf{three} different training formats mentioned in the text.
    \item Find the words or expressions that correspond to these definitions:
    \begin{itemize}
        \item[a)] A meeting where all participants are together (opposite of breakout groups).
        \item[b)] Learning from people at the same level as you.
        \item[c)] Guidance by an experienced person.
    \end{itemize}
    \item Classify the “five-day residential workshop” according to: duration, modality, and accommodation.
    \item \textbf{Transform:} Rewrite the sentence “Selected participants will benefit from peer-to-peer learning” using the verb \eng{engage in}.
\end{enumerate}
\tcblower
\textbf{Solution détaillée :}
\begin{enumerate}
    \item \textbf{Three formats :} 1. \eng{Workshop} (atelier pratique, interactif, plusieurs jours). 2. \eng{Plenary sessions} / \eng{Breakout groups} (sous-formats au sein du workshop). 3. \eng{Webinar} (séminaire en ligne, suivi à distance). \textit{Note : “Simulation exercises” et “Mentorship” sont des méthodes/pédagogies, pas des formats d’événement à part entière.}
    \item \textbf{Vocabulaire :}
    \begin{itemize}
        \item[a)] \eng{Plenary sessions} (séances plénières).
        \item[b)] \eng{Peer-to-peer learning} (apprentissage par les pairs).
        \item[c)] \eng{Mentorship} (mentorat / tutorat).
    \end{itemize}
    \item \textbf{Classification du “five-day residential workshop” :}
    \begin{itemize}
        \item \textbf{Duration :} \eng{Short-term / Intensive} (5 jours).
        \item \textbf{Modality :} \eng{In-person / Face-to-face / Residential} (présentiel, hébergement sur place).
        \item \textbf{Accommodation :} \eng{Residential} (les participants logent sur place).
    \end{itemize}
    \item \textbf{Transformation :} \eng{Selected participants will \textbf{engage in} peer-to-peer learning.} (Structure : \eng{engage in + V-ing / nom}). \textit{Autre possibilité :} \eng{Selected participants will engage in learning from their peers.}
\end{enumerate}
\textbf{Conclusion :} Un même événement (le \eng{workshop}) contient plusieurs séquences pédagogiques (\eng{plenary, breakout, simulation}). Le vocabulaire précis permet de décrire l’architecture de la formation.
\end{exoresolu}

\begin{methode}[Identifier les acteurs et décrire leurs rôles dans une formation]
\textbf{Objectif :} Maîtriser le lexique des rôles (organisateur, formateur, participant, expert) et les verbes d’action associés pour décrire qui fait quoi.\par
\textbf{Démarche en 3 étapes :}
\begin{enumerate}
    \item \textbf{Distinguier les trois cercles :} \textbf{Commanditaires/Organisateurs} (\eng{organise, sponsor, fund, coordinate, oversee}) -- \textbf{Experts/Formateurs} (\eng{facilitate, train, coach, mentor, lecture, moderate, debrief}) -- \textbf{Bénéficiaires/Participants} (\eng{attend, enrol, participate, contribute, network, evaluate, implement}).
    \item \textbf{Noter la nuance \eng{teach} vs \eng{facilitate} :} En formation d’adultes (andragogie), on préfère \eng{facilitator} (qui guide le processus) à \eng{teacher} (qui transmet le savoir). \eng{Moderator} anime un débat/panel ; \eng{Trainer} forme à des compétences techniques.
    \item \textbf{Utiliser la voix passive pour le focus sur l’événement :} \eng{The workshop \textbf{is facilitated by} an expert.} / \eng{Participants \textbf{are selected based on} criteria.}
\end{enumerate}
\textbf{Collocations clés :} \eng{lead a session, chair a meeting, deliver a presentation, share best practices, build capacity, empower participants}.
\end{methode}

\begin{exoresolu}[Rôles, verbes d’action et transformation de phrases]
\textbf{Énoncé :} D’après le texte suivant, complétez le tableau, puis transformez les phrases.\par
\begin{textebox}[Extrait de rapport d’activité]
\eng{The National Democratic Institute (NDI) \textbf{sponsored} a \textbf{three-day Training of Trainers (ToT)} in Yaoundé. Master \textbf{trainers} \textbf{equipped} 25 \textbf{trainees} with skills to \textbf{cascade} the training in their communities. A local \textbf{consultant} \textbf{designed} the curriculum. The \textbf{participants} \textbf{drafted} action plans and \textbf{committed to} organising \textbf{step-down trainings} before December.}
\end{textebox}
\begin{enumerate}
    \item Complete the table with words from the text.
    \begin{center}
    \begin{tabular}{|l|l|l|}
    \hline
    \rowcolor{popblueL} \textbf{Rôle (Nom)} & \textbf{Verbe d’action (du texte)} & \textbf{Fonction principale} \\
    \hline
    Sponsor / Funder & \eng{sponsored} & Finance l’activité \\
    \hline
    \eng{...} & \eng{designed} & Conçoit le programme \\
    \hline
    \eng{Master trainers} & \eng{...} & Forment les formateurs \\
    \hline
    \eng{Trainees / Participants} & \eng{drafted, committed to} & \eng{...} \\
    \hline
    \end{tabular}
    \end{center}
    \item \textbf{Transform to Passive Voice} (focus on the action/receiver):
    \begin{itemize}
        \item[a)] Master trainers equipped 25 trainees.
        \item[b)] A local consultant designed the curriculum.
    \end{itemize}
    \item \textbf{Vocabulary check :} Explain the difference between \eng{ToT (Training of Trainers)} and \eng{Step-down training} (or \eng{cascade training}).
\end{enumerate}
\tcblower
\textbf{Solution détaillée :}
\begin{enumerate}
    \item \textbf{Tableau complété :}
    \begin{center}
    \begin{tabular}{|l|l|l|}
    \hline
    \rowcolor{popblueL} \textbf{Rôle} & \textbf{Verbe} & \textbf{Fonction} \\
    \hline
    Sponsor / Funder & \eng{sponsored} & Finance l’activité \\
    \hline
    \eng{Consultant / Curriculum developer} & \eng{designed} & Conçoit le programme \\
    \hline
    \eng{Master trainers / Facilitators} & \eng{equipped} & Forment les formateurs (transfert de compétences) \\
    \hline
    \eng{Trainees / Participants / Beneficiaries} & \eng{drafted, committed to} & Reçoivent la formation, produisent des livrables, s’engagent à répliquer \\
    \hline
    \end{tabular}
    \end{center}
    \item \textbf{Voix passive :}
    \begin{itemize}
        \item[a)] \eng{25 trainees \textbf{were equipped} by master trainers.} (Ou \eng{were equipped with skills...})
        \item[b)] \eng{The curriculum \textbf{was designed} by a local consultant.}
    \end{itemize}
    \item \textbf{Différence ToT vs Step-down :}
    \begin{itemize}
        \item \eng{Training of Trainers (ToT) :} Formation \textbf{initiale} destinée à des personnes qui deviendront formateurs (\eng{master trainers $\to$ trainees}). Objectif : \eng{build a pool of trainers}.
        \item \eng{Step-down training / Cascade training :} Formation \textbf{secondaire} organisée par les nouveaux formateurs (\eng{trainees}) pour leurs pairs ou leur communauté (\eng{trainees $\to$ community}). Objectif : \eng{multiply the impact / reach the grassroots}.
    \end{itemize}
\end{enumerate}
\textbf{Conclusion :} La terminologie \eng{ToT} et \eng{cascade/step-down} est standard dans la coopération internationale et les ONG. Le verbe \eng{equip} (avec \eng{with}) est crucial : \eng{equip someone with skills}.
\end{exoresolu}

\begin{methode}[Mobiliser le lexique thématique de la démocratie en contexte de formation]
\textbf{Objectif :} Utiliser le vocabulaire spécifique aux thèmes démocratiques (élections, gouvernance, droits, participation) dans des phrases complètes ou un paragraphe argumentatif.\par
\textbf{Démarche en 4 étapes :}
\begin{enumerate}
    \item \textbf{Brainstormer par champ lexical :} \eng{Elections} (ballot, polling station, voter registration, observation, integrity) -- \eng{Governance} (accountability, transparency, rule of law, separation of powers, decentralisation) -- \eng{Rights} (civil liberties, freedom of speech, assembly, association, minority rights) -- \eng{Participation} (civic engagement, advocacy, lobbying, public consultation, grassroots mobilisation).
    \item \textbf{Associer les verbes “puissants” (collocations) :} \eng{Promote accountability, uphold the rule of law, monitor elections, advocate for rights, foster dialogue, strengthen institutions, empower citizens, combat corruption}.
    \item \textbf{Structurer l’argumentation :} \eng{The training aims to + V-inf...} / \eng{Participants will be equipped to + V-inf...} / \eng{Key topics include + V-ing / nouns...}.
    \item \textbf{Vérifier les faux amis (voir Méthode 4) et l’orthographe des termes institutionnels.}
\end{enumerate}
\textbf{Expressions figées utiles :} \eng{free and fair elections, system of checks and balances, civic space, human rights defenders, capacity building, stakeholder engagement}.
\end{methode}

\begin{exoresolu}[Production écrite guidée : Objectifs d’un module de formation]
\textbf{Énoncé :} Vous êtes chargé de rédiger les \eng{Learning Objectives} (objectifs pédagogiques) pour un module de 3 jours intitulé \eng{“Strengthening Youth Participation in Local Governance”}. Utilisez le vocabulaire thématique imposé.\par
\textbf{Contraintes :}
\begin{itemize}
    \item Utilisez au moins 8 mots/expressions de la liste : \eng{accountability, transparency, advocacy, stakeholder, decentralisation, civic engagement, bylaws, public consultation, grassroots, empower}.
    \item Structurez avec : \eng{By the end of this module, participants will be able to:} + liste à puces (verbes d’action : \eng{analyse, design, conduct, facilitate, draft, evaluate}).
    \item 80--100 mots.
\end{itemize}
\tcblower
\textbf{Solution détaillée (Modèle de réponse) :}
\begin{textebox}[Learning Objectives -- Module: Strengthening Youth Participation in Local Governance]
\eng{By the end of this module, participants will be able to:}
\begin{itemize}
    \item \eng{\textbf{Analyse} the principles of \textbf{decentralisation} and the role of \textbf{local councils} in promoting \textbf{accountability} and \textbf{transparency}.}
    \item \eng{\textbf{Identify} key \textbf{stakeholders} and map \textbf{grassroots} entry points for youth \textbf{civic engagement}.}
    \item \eng{\textbf{Design} an \textbf{advocacy} strategy to influence \textbf{public consultation} processes on local \textbf{bylaws}.}
    \item \eng{\textbf{Facilitate} dialogue between youth groups and municipal authorities to \textbf{empower} young citizens in decision-making.}
    \item \eng{\textbf{Evaluate} the effectiveness of participatory mechanisms in their respective communities.}
\end{itemize}
\end{textebox}
\textbf{Justifications grammaticales et lexicales :}
\begin{itemize}
    \item \textbf{Structure objectif :} \eng{By the end of..., participants will be able to + BASE VERBALE} (standard Bloom’s taxonomy).
    \item \textbf{Collocations respectées :} \eng{promote accountability/transparency} (verbe + nom abstrait), \eng{influence processes}, \eng{map entry points}, \eng{empower citizens}.
    \item \textbf{Vocabulaire imposé intégré :} 9/8 mots utilisés (\eng{decentralisation, accountability, transparency, stakeholders, grassroots, civic engagement, advocacy, public consultation, bylaws, empower}).
    \item \textbf{Orthographe :} \eng{Bylaws} (un mot, souvent pluriel), \eng{Stakeholders} (pluriel), \eng{Grassroots} (adjectif invariable ici, ou nom \eng{the grassroots}).
\end{itemize}
\textbf{Conclusion :} Un objectif pédagogique bien rédigé utilise un verbe d’action observable (\eng{analyse, design, draft}) et non un verbe d’état (\eng{know, understand}).
\end{exoresolu}

\begin{methode}[Déjouer les faux amis, calques et pièges d’orthographe]
\textbf{Objectif :} Éviter les erreurs classiques des francophones sur le vocabulaire de la formation et de la démocratie.\par
\textbf{Démarche :} Pour chaque mot piège, appliquer la règle \textbf{STOP} : \textbf{S}ortir le mot -- \textbf{T}rouver le vrai sens anglais -- \textbf{O}bserver le mot anglais correct -- \textbf{P}roduire une phrase correcte.\par
\textbf{Tableau des 10 pièges majeurs :}
\begin{center}
\begin{tabularx}{\linewidth}{|X|X|X|}
\hline
\rowcolor{popblueL} \textbf{Mot / Expression (FR)} & \textbf{Faux ami / Erreur courante} & \textbf{Anglais correct + Exemple} \\
\hline
\eng{Une formation} & \eng{A formation} (géologie, militaire) & \eng{A training course / A workshop / A programme} \\
\hline
\eng{Un stage} & \eng{A stage} (scène, étage) & \eng{An internship / A work placement / A training course} \\
\hline
\eng{Assister à} & \eng{Assist} (aider) & \eng{Attend / Take part in / Participate in} \\
\hline
\eng{Un animateur} & \eng{An animator} (dessin animé, vie sociale) & \eng{A facilitator / A trainer / A moderator} \\
\hline
\eng{Une animation} & \eng{An animation} (dessin animé) & \eng{Facilitation / Ice-breakers / Energizers} \\
\hline
\eng{Un ordre du jour} & \eng{An order of the day} & \eng{An agenda} \\
\hline
\eng{Un compte-rendu} & \eng{A report} (trop large) & \eng{Minutes (of a meeting) / A proceedings report} \\
\hline
\eng{La gouvernance} & \eng{Gouvernance} (orthographe FR) & \eng{Governance} (pas de ‘u’) \\
\hline
\eng{Un électeur} & \eng{An elector} (grand électeur, rare) & \eng{A voter} (citoyen qui vote) \\
\hline
\eng{Sensibiliser} & \eng{Sensibilize} (n’existe pas) & \eng{Raise awareness (about) / Sensitize (terme technique ONG, mais \textbf{raise awareness} est plus courant)} \\
\hline
\end{tabularx}
\end{center}
\textbf{Règle d’or :} \eng{Assist = Help}. \eng{Attend = Être présent}. \eng{Facilitator $\neq$ Animator}.
\end{methode}

\begin{exoresolu}[Correction d’erreurs types “Bac” -- Faux amis et calques]
\textbf{Énoncé :} Les phrases suivantes contiennent des erreurs de vocabulaire (faux amis, calques, mauvais choix lexical). Corrigez-les et justifiez chaque correction.\par
\begin{enumerate}
    \item \eng{The NGO organised a \textbf{formation} on human rights for 50 teachers.}
    \item \eng{During the workshop, the \textbf{animator} asked us to \textbf{assist} the \textbf{simulation}.}
    \item \eng{The \textbf{stage} will take place in the conference hall for three days.}
    \item \eng{We need to \textbf{sensibilize} the population about the new \textbf{electoral code}.}
    \item \eng{The \textbf{minutes} of the meeting were written by the \textbf{rapporteur}.} (Phrase correcte ? Vérifiez \eng{rapporteur}).
\end{enumerate}
\tcblower
\textbf{Solution détaillée :}
\begin{enumerate}
    \item \textbf{Erreur :} \eng{Formation} (faux ami, sens géologique/militaire). \textbf{Correction :} \eng{The NGO organised a \textbf{training course / workshop / programme} on human rights...}
    \item \textbf{Erreurs :} \eng{Animator} (dessin animé), \eng{Assist} (aider au lieu de participer), \eng{Assist the simulation} (calque structurel). \textbf{Correction :} \eng{During the workshop, the \textbf{facilitator / trainer} asked us to \textbf{participate in / take part in / engage in} the \textbf{simulation exercise}.}
    \item \textbf{Erreur :} \eng{Stage} (scène, plateau). \textbf{Correction :} \eng{The \textbf{training course / workshop / session} will take place...} (Si stage pro : \eng{internship}).
    \item \textbf{Erreur :} \eng{Sensibilize} (n’existe pas en anglais courant). \textbf{Correction :} \eng{We need to \textbf{raise awareness among} the population about the new \textbf{electoral code}.} (\eng{Sensitize} existe en jargon ONG/développement mais \eng{raise awareness} est l’anglais standard).
    \item \textbf{Analyse :} \eng{Minutes} = Correct (terme standard pour PV). \eng{Rapporteur} = \textbf{Faux ami / Anglicisme déconseillé}. En anglais standard : \eng{The minute-taker} ou \eng{The secretary}. \eng{Rapporteur} s’utilise seulement dans des contextes très formels ONU/UE (\eng{Special Rapporteur}). \textbf{Version naturelle :} \eng{The \textbf{minutes} of the meeting were taken by the \textbf{secretary / minute-taker}.}
\end{enumerate}
\textbf{Conclusion :} Ces 5 erreurs couvrent 80\% des fautes de vocabulaire observées aux examens. Mémorisez les paires : \eng{Formation $\to$ Training}, \eng{Assister $\to$ Attend}, \eng{Animateur $\to$ Facilitator}, \eng{Stage $\to$ Internship/Course}, \eng{Sensibiliser $\to$ Raise awareness}.
\end{exoresolu}

\begin{methode}[Construire des familles de mots (dérivation) et gérer prononciation / orthographe]
\textbf{Objectif :} Dériver noms, verbes, adjectifs, adverbes à partir des racines du thème (democracy, participate, govern, elect, transparent) et anticiper les pièges de prononciation (accent tonique) et d’orthographe.\par
\textbf{Démarche en 3 temps :}
\begin{enumerate}
    \item \textbf{Identifier la racine et les suffixes productifs :} \eng{-tion/-sion} (nom), \eng{-ment} (nom), \eng{-ize/-ise} (verbe), \eng{-ive/-atory} (adjectif), \eng{-ly} (adverbe), \eng{re-/de-/en-} (préfixes).
    \item \textbf{Remplir le tableau de dérivation (Word Formation) :} Toujours vérifier l’existence du mot (certains adjectifs n’existent pas : \eng{*participative} $\to$ \eng{participatory}).
    \item \textbf{Marquer l’accent tonique (stress) :} Règle générale : suffixes \eng{-tion, -sion, -ic, -ity, -ive} $\to$ accent sur la syllabe \textbf{avant} le suffixe. Ex: \eng{deMOCracy $\to$ demoCRATic $\to$ demoCRATically}.
\end{enumerate}
\textbf{Tableau de référence (à mémoriser) :}
\begin{center}
\begin{tabular}{|l|l|l|l|l|}
\hline
\rowcolor{popblueL} \textbf{Verbe} & \textbf{Nom (Action/Concept)} & \textbf{Nom (Agent)} & \textbf{Adjectif} & \textbf{Adverbe} \\
\hline
\eng{Participate} & \eng{Participation} & \eng{Participant} & \eng{Participatory} & \eng{Participatorily} (rare) \\
\hline
\eng{Govern} & \eng{Governance / Government} & \eng{Governor} & \eng{Governmental} & \eng{Governmentally} (rare) \\
\hline
\eng{Democratize} & \eng{Democratization} & \eng{Democrat} & \eng{Democratic} & \eng{Democratically} \\
\hline
\eng{Elect} & \eng{Election} & \eng{Elector / Voter} & \eng{Electoral} & \eng{Electorally} \\
\hline
\eng{Observe} & \eng{Observation} & \eng{Observer} & \eng{Observant} & \eng{Observantly} (rare) \\
\hline
\eng{Transparent} (adj) & \eng{Transparency} & -- & \eng{Transparent} & \eng{Transparently} \\
\hline
\eng{Accountable} (adj) & \eng{Accountability} & -- & \eng{Accountable} & \eng{Accountably} \\
\hline
\eng{Empower} & \eng{Empowerment} & -- & \eng{Empowering} & \eng{Empoweringly} (rare) \\
\hline
\end{tabular}
\end{center}
\textbf{Prononciation clé (notation française) :} \eng{Democracy} « dé-mo-kra-si » ; \eng{Democratic} « dé-mo-kra-tik » (accent sur \textbf{cra}) ; \eng{Participatory} « par-ti-si-pa-to-ri » (accent sur \textbf{to}) ; \eng{Governance} « gou-ve-nance » (son \eng{v} pas \eng{g}) ; \eng{Accountability} « a-koun-ta-bi-li-ti » (accent sur \textbf{bi}).
\end{methode}

\begin{exoresolu}[Word Formation \& Stress Placement -- Type Examen]
\textbf{Énoncé :}
\begin{enumerate}
    \item Complete the sentences with the correct form of the word in CAPITALS.
    \begin{enumerate}
        \item The \eng{\_\_\_\_\_\_\_\_\_\_\_\_} of the electoral process is guaranteed by independent observers. (TRANSPARENT)
        \item Youth \eng{\_\_\_\_\_\_\_\_\_\_\_\_} in local decision-making remains low despite new laws. (PARTICIPATE)
        \item The workshop aims at the \eng{\_\_\_\_\_\_\_\_\_\_\_\_} of women leaders in rural areas. (EMPOWER)
        \item A good \eng{\_\_\_\_\_\_\_\_\_\_\_\_} must remain neutral and objective. (OBSERVE)
        \item The \eng{\_\_\_\_\_\_\_\_\_\_\_\_} of power is a key principle of democracy. (SEPARATE)
    \end{enumerate}
    \item \textbf{Prononciation :} Soulignez la syllabe accentuée (stress) dans ces mots : \eng{De-mo-cra-cy}, \eng{De-mo-cra-tic}, \eng{Par-ti-ci-pa-to-ry}, \eng{Gou-ver-nance}, \eng{Ac-coun-ta-bi-li-ty}.
    \item \textbf{Orthographe :} Choisissez l’orthographe correcte : \eng{Gouvernance / Governance} ; \eng{Accountability / Accountability} ; \eng{Democratisation / Democratization}.
\end{enumerate}
\tcblower
\textbf{Solution détaillée :}
\begin{enumerate}
    \item \textbf{Word Formation :}
    \begin{enumerate}
        \item \eng{Transparency} (Nom formé sur adj \eng{transparent} + \eng{-cy}). \textit{Piège : pas “transparence”.}
        \item \eng{Participation} (Nom sur verbe \eng{participate} + \eng{-tion}). \textit{Piège : adj = \eng{participatory}, pas *participative.}
        \item \eng{Empowerment} (Nom sur verbe \eng{empower} + \eng{-ment}). \textit{Concept clé ONG.}
        \item \eng{Observer} (Nom d’agent sur verbe \eng{observe} + \eng{-r}). \textit{Le nom d’action est \eng{observation}.}
        \item \eng{Separation} (Nom sur verbe \eng{separate} + \eng{-tion}). \textit{Notez la disparition du ‘e’ final : separat + ion.}
    \end{enumerate}
    \item \textbf{Stress (Accent tonique) -- Règle des suffixes :}
    \begin{itemize}
        \item \eng{De-mo-\textbf{cra}-cy} (suffixe \eng{-cy} $\to$ antépénultième).
        \item \eng{De-mo-\textbf{cra}-tic} (suffixe \eng{-ic} $\to$ pénultième).
        \item \eng{Par-ti-ci-pa-\textbf{to}-ry} (suffixe \eng{-ory} $\to$ pénultième).
        \item \eng{Gou-\textbf{ver}-nance} (suffixe \eng{-ance} $\to$ antépénultième, radical \eng{govern}).
        \item \eng{Ac-coun-ta-\textbf{bi}-li-ty} (suffixe \eng{-ity} $\to$ antépénultième).
    \end{itemize}
    \item \textbf{Orthographe :}
    \begin{itemize}
        \item \eng{Governance} (Pas de ‘u’ après ‘g’ -- \textbf{attention au calque FR}).
        \item \eng{Accountability} (Double ‘c’, double ‘t’, ‘a’ avant ‘bility’).
        \item \eng{Democratization} (Suffixation standard \eng{-ize} en anglais international / Oxford ; \eng{-ise} accepté en anglais britannique mais \eng{-ization} est la norme pour ce mot long). \textit{Jamais “Democratisation”.}
    \end{itemize}
\end{enumerate}
\end{exoresolu}

\begin{methode}[Produire un dialogue de négociation / inscription à une formation]
\textbf{Objectif :} Réemployer le lexique (formats, acteurs, thèmes, inscriptions) dans une interaction orale réaliste (téléphone, accueil, entretien).\par
\textbf{Démarche en 4 étapes (Jeu de rôle) :}
\begin{enumerate}
    \item \textbf{Assigner les rôles :} \textbf{A} = Candidat / Participant potentiel. \textbf{B} = Coordinateur de formation / Secrétariat.
    \item \textbf{Structurer l’échange (Script mental) :}
    \begin{itemize}
        \item \textbf{Accroche :} \eng{Hello, I’m calling about / I’d like to enquire about...}
        \item \textbf{Questions clés (A) :} Dates, venue, fees/scholarship, eligibility, certification, facilitators.
        \item \textbf{Réponses / Infos (B) :} \eng{The workshop runs from... to...} ; \eng{It is fully funded / residential} ; \eng{The deadline is...} ; \eng{You will receive a certificate of completion.}
        \item \textbf{Clôture :} \eng{How do I apply? / Send your CV and motivation letter to...} ; \eng{Thank you, I will apply.}
    \end{itemize}
    \item \textbf{Injecter le vocabulaire cible :} Minimum 5 mots de la leçon (\eng{facilitator, residential, plenary, breakout, curriculum, accreditation, deadline, enrol}).
    \item \textbf{Soigner la politesse et la phonologie :} \eng{Could you please...? / Would it be possible to...? / I look forward to hearing from you.}
\end{enumerate}
\textbf{Expressions fonctionnelles :} \eng{To be eligible for}, \eng{To meet the criteria}, \eng{To submit an application}, \eng{A fully-funded scholarship}, \eng{Logistics are covered}.
\end{methode}

\begin{exoresolu}[Production orale guidée : Dialogue d’inscription -- Évaluation Bac]
\textbf{Énoncé :} \textbf{Role-play.} Vous êtes \textbf{Candidate A}. Votre partenaire est \textbf{The Training Coordinator}.\par
\textbf{Your card :} You want to attend the “Youth Leadership \& Democratic Governance” workshop in Douala next month. Call the coordinator to get details. You must ask about: 1. Exact dates and venue. 2. Whether it is residential or not. 3. The profile of the facilitators. 4. The application procedure and deadline. 5. If there is a certificate.\par
\textbf{Rédigez votre partie du dialogue (repli écrit pour préparation).}
\tcblower
\textbf{Solution détaillée (Modèle de réplique Candidate A -- Niveau B2/Bac) :}
\begin{textebox}[Dialogue simulé -- Candidate A]
\eng{\textbf{Candidate A:} Good morning. My name is Achille Ngu. I’m calling to \textbf{enquire about} the “Youth Leadership and Democratic Governance” \textbf{workshop} advertised for next month in Douala.}
\eng{\textbf{Coordinator:} Good morning Achille. Certainly. It’s a \textbf{four-day residential training} from the 12th to the 15th of August at the City Hall conference centre.}
\eng{\textbf{Candidate A:} \textbf{Residential}? That means \textbf{accommodation} is provided?}
\eng{\textbf{Coordinator:} Yes, full board and \textbf{logistics} are covered for selected participants from outside Douala.}
\eng{\textbf{Candidate A:} Great. Could you tell me who the \textbf{facilitators} are? I’d like to know their \textbf{background} in \textbf{governance} issues.}
\eng{\textbf{Coordinator:} We have two \textbf{master trainers} from the National School of Administration (ENAM) and a \textbf{civic education} expert from an international NGO.}
\eng{\textbf{Candidate A:} Sounds excellent. What is the \textbf{application procedure} and the \textbf{deadline}?}
\eng{\textbf{Coordinator:} You must \textbf{submit} an online \textbf{application form} with a \textbf{motivation letter} and CV before July 30th. \textbf{Selection} is based on \textbf{merit} and \textbf{geographical balance}.}
\eng{\textbf{Candidate A:} Understood. And do participants receive a \textbf{certificate} at the end?}
\eng{\textbf{Coordinator:} Yes, a \textbf{Certificate of Completion} is awarded upon \textbf{full attendance} and submission of a \textbf{personal action plan}.}
\eng{\textbf{Candidate A:} Perfect. I have all the information I need. I will \textbf{apply} right away. Thank you very much for your help.}
\eng{\textbf{Coordinator:} You’re welcome. Good luck with your application.}
\end{textebox}
\textbf{Analyse lexicale (Mots cibles intégrés : 12/5 requis) :} \eng{enquire about, workshop, residential training, accommodation, logistics, facilitators, background, governance, master trainers, civic education, application procedure, deadline, submit, application form, motivation letter, selection, merit, certificate of completion, full attendance, action plan, apply.}
\textbf{Points de grammaire vérifiés :} \eng{Enquire about} (préposition), \eng{Submit + objet}, \eng{Selection is based on} (passif + préposition), \eng{Upon + V-ing} (conjonction temporelle formelle).
\end{exoresolu}

\begin{attention}[Les 5 erreurs qui coûtent des points au Bac -- Spécial Vocabulaire Démocratie / Formation]
\begin{enumerate}
    \item \textbf{\eng{Assist} au lieu de \eng{Attend / Take part in}.} \eng{“I assisted the workshop”} = “J’ai aidé l’atelier”. \textbf{Correct :} \eng{“I \textbf{attended} the workshop”} ou \eng{“I \textbf{took part in} the workshop”}.
    \item \textbf{\eng{Formation / Stage} calqués.} \eng{“I did a formation”} / \eng{“I did a stage”}. \textbf{Correct :} \eng{“I \textbf{attended a training course / workshop}”} ; \eng{“I \textbf{did an internship / work placement}”} (stage pro).
    \item \textbf{\eng{Animator} pour “Animateur de formation”.} \textbf{Correct :} \eng{\textbf{Facilitator}} (atelier), \eng{\textbf{Trainer}} (compétences), \eng{\textbf{Moderator}} (débat/panel). \eng{Animator} = dessinateur de cartoons.
    \item \textbf{\eng{Sensibilize} (néologisme) / \eng{Conscientize} (jargon rare).} \textbf{Correct :} \eng{\textbf{Raise awareness (about/on)}} (standard). Ex: \eng{“The NGO aims to \textbf{raise awareness about} voter registration.”}
    \item \textbf{Orthographe \eng{Gouvernance} / \eng{Démocratisation} / \eng{Particpation}.} \textbf{Correct :} \eng{\textbf{Governance}} (pas de u), \eng{\textbf{Democratization}} (suffixe \eng{-ization}), \eng{\textbf{Participation}} (double ‘t’, pas ‘c’). \textit{Astuce :} \eng{Govern} $\to$ \eng{Government} (u) mais \eng{Governance} (pas de u).
\end{enumerate}
\textbf{Check-list finale avant de rendre la copie :} Ai-je utilisé \eng{facilitator} ? Ai-je écrit \eng{raise awareness} ? Ai-je mis l’accent tonique sur la bonne syllabe à l’oral (\eng{demoCRATic}) ? Mes objectifs pédagogiques commencent-ils par un verbe d’action (\eng{analyse, design, draft}) ?
\end{attention}

\newpage

\coursec{Exercices}

\courssub{Je vérifie mes connaissances}

\begin{exercice}[QCM : le bon mot]
Pour chaque phrase, choisis la bonne réponse (a, b, c ou d).

\begin{enumerate}
\item I would like to \cle{\dots} a workshop on human rights next month.
\begin{itemize}
\item[a)] assist \item[b)] attend \item[c)] assist to \item[d)] join to
\end{itemize}
\item The \cle{\dots} guided the participants through the group activities.
\begin{itemize}
\item[a)] animator \item[b)] facilitator \item[c)] animator of formation \item[d)] trainerer
\end{itemize}
\item At the end of the course, every trainee received a \cle{\dots} of participation.
\begin{itemize}
\item[a)] diploma \item[b)] certificate \item[c)] note \item[d)] paper
\end{itemize}
\item Our school organised a training \cle{\dots} democracy and good governance.
\begin{itemize}
\item[a)] in \item[b)] about \item[c)] on \item[d)] of
\end{itemize}
\item The seminar aims to raise \cle{\dots} of citizens' rights among young people.
\begin{itemize}
\item[a)] conscience \item[b)] awareness \item[c)] knowledge \item[d)] sensibility
\end{itemize}
\end{enumerate}
\end{exercice}

\begin{exercice}[Vrai ou faux ?]
Lis le texte suivant, puis indique si chaque affirmation est \textbf{True} ou \textbf{False}. Justifie ta réponse par une courte citation du texte.

\begin{textebox}[A seminar in Douala]
Last Saturday, the Youth Civic Club of Douala organised a one-day seminar on good governance at the Lycée de New-Bell. About sixty students enrolled for the training. In the morning, a guest lecturer from the University of Douala delivered a talk on transparency and the fight against corruption. In the afternoon, the participants were divided into four workshops, each led by a trained facilitator. At the end of the day, the organisers distributed certificates of attendance and invited the students to share what they had learned with their classmates.
\end{textebox}

\begin{enumerate}
\item The seminar lasted a whole week.
\item Only university students attended the training.
\item The morning session focused on transparency and corruption.
\item The participants worked in groups in the afternoon.
\item The students received a diploma at the end of the seminar.
\end{enumerate}
\end{exercice}

\begin{exercice}[Vocabulaire : qui fait quoi ?]
Associe chaque acteur de la formation (colonne A) à sa définition (colonne B).

\begin{center}
\begin{tabularx}{\linewidth}{|l|X|}
\hline
\rowcolor{popblueL} \textbf{A — Actor} & \textbf{B — Definition} \\
\hline
1. facilitator & a) a person who officially signs up for a course \\
\hline
2. trainee & b) a person who guides group discussions and activities \\
\hline
3. lecturer & c) a person who receives training \\
\hline
4. participant & d) a person who gives a formal talk to an audience \\
\hline
5. organiser & e) a person who plans and prepares the event \\
\hline
\end{tabularx}
\end{center}
\end{exercice}

\begin{exercice}[Complète avec le bon mot]
Complète chaque phrase avec la forme correcte du verbe entre parenthèses.

\begin{enumerate}
\item Last year, over two hundred students \cle{\dots} (enrol) for the citizenship course.
\item The trainer \cle{\dots} (deliver) an interesting lecture on voting rights yesterday.
\item Our club \cle{\dots} (organise) a workshop on the rule of law next term.
\item Since January, I \cle{\dots} (attend) three seminars on human rights.
\item While the facilitator was speaking, the participants \cle{\dots} (take) notes.
\end{enumerate}
\end{exercice}

\courssub{Je m'entraîne}

\begin{exercice}[Traduction]
Traduis les phrases suivantes en anglais.

\begin{enumerate}
\item J'aimerais suivre une formation sur la démocratie pendant les vacances.
\item L'animateur nous a divisés en groupes de cinq participants.
\item À la fin du séminaire, chaque stagiaire a reçu une attestation de participation.
\item Cette formation vise à sensibiliser les jeunes aux droits de l'homme.
\item Plus de cent élèves se sont inscrits à l'atelier sur la bonne gouvernance.
\end{enumerate}
\end{exercice}

\begin{exercice}[Attention aux faux amis !]
Les phrases suivantes ont été écrites par des élèves francophones. Chacune contient une erreur de vocabulaire (faux ami ou calque du français). Corrige chaque phrase.

\begin{enumerate}
\item I made a formation on democracy last year.
\item The animator trained us during the seminar.
\item I assisted to a workshop on human rights.
\item She received a diploma after a two-day training course.
\item The conference was very interesting; I learned a lot.
\end{enumerate}
\end{exercice}

\begin{exercice}[Texte à trous : a training camp]
Complète le compte rendu suivant avec les mots de la liste : \emph{facilitator — certificate — awareness — enrolled — rule of law — workshop}.

\begin{textebox}[Report : a civic training camp]
From 10 to 15 August, fifty students from different schools in Bamenda \cle{\dots} for a civic training camp on democracy. The programme included a daily \cle{\dots} on topics such as elections, tolerance and the \cle{\dots}. Each session was led by an experienced \cle{\dots} who encouraged everyone to speak. The main goal of the camp was to raise \cle{\dots} of citizens' rights and duties among young people. On the last day, every participant received a \cle{\dots} of participation from the organisers.
\end{textebox}
\end{exercice}

\begin{exercice}[Appariement : collocations]
Associe chaque verbe de la colonne A au nom qui convient dans la colonne B pour former une collocation correcte, puis utilise trois de ces collocations dans des phrases de ton invention.

\begin{center}
\begin{tabularx}{\linewidth}{|X|X|}
\hline
\rowcolor{popblueL} \textbf{A — Verb} & \textbf{B — Noun} \\
\hline
1. raise & a) a workshop \\
\hline
2. attend & b) corruption \\
\hline
3. deliver & c) democracy \\
\hline
4. fight & d) awareness \\
\hline
5. promote & e) a lecture \\
\hline
\end{tabularx}
\end{center}
\end{exercice}

\courssub{Je me prépare au Bac}

\begin{exercice}[Compréhension et langue — type Bac]
Lis attentivement le texte suivant, puis réponds aux questions.

\begin{textebox}[Training young citizens in Yaoundé]
Every year, during the second term, the bilingual high school of Yaoundé organises a citizenship week for its final-year students. The programme, prepared by teachers in collaboration with a local association, combines lectures, debates and practical workshops. On Monday, a magistrate delivers a lecture on the Constitution and the rule of law. On Tuesday and Wednesday, the students attend workshops led by trained facilitators: they learn how to debate respectfully, how to run a meeting and how to vote in a transparent election. Thursday is devoted to a mock election: the students elect their class representatives under the supervision of their teachers. Finally, on Friday, a closing ceremony is held, during which every participant receives a certificate of attendance.

According to the head teacher, this training is essential. “Young people must understand that democracy is not only about voting,” she explains. “It is about listening to others, accepting differences and taking responsibility.” Many students agree. “Before this week, I thought elections were just about putting a paper in a box,” says André, a Terminale student. “Now I know that being a citizen means much more.” The organisers hope that the participants will share their knowledge with their families and become active citizens in their communities.
\end{textebox}

\textbf{A. Comprehension questions.} Answer in your own words.

\begin{enumerate}
\item Who organises the citizenship week, and for whom?
\item Mention two activities that take place during the week.
\item What happens on Thursday?
\item According to the head teacher, what does democracy mean besides voting?
\item How has André's idea of citizenship changed? Quote from the text to support your answer.
\end{enumerate}

\textbf{B. Language work.}

\begin{enumerate}
\item Find in the text: (a) a word meaning “a person who guides a workshop”; (b) a word meaning “a document given at the end of a course”; (c) an expression meaning “a practice election”.
\item Complete with the right preposition: (a) The students attended a workshop \cle{\dots} the rule of law. (b) Democracy is not only \cle{\dots} voting. (c) The election took place \cle{\dots} the supervision of the teachers.
\item Put the verbs in brackets in the correct tense: “Last year, the school (organise) its first citizenship week, and since then, many students (join) the civic club.”
\end{enumerate}
\end{exercice}

\begin{exercice}[Traduction — type Bac]
Traduis en français le passage suivant, extrait du texte ci-dessus :

\begin{textebox}[Passage to translate]
“Young people must understand that democracy is not only about voting,” she explains. “It is about listening to others, accepting differences and taking responsibility.”
\end{textebox}

Puis traduis en anglais les deux phrases suivantes :

\begin{enumerate}
\item Chaque année, notre lycée organise une semaine de la citoyenneté pour les élèves de Terminale.
\item Les participants ont appris à débattre avec respect et à accepter les différences.
\end{enumerate}
\end{exercice}

\begin{exercice}[Composition — type Bac]
\textbf{Topic :} \emph{Describe a training activity on citizenship organised in your school (or in a school you know) and explain its importance for young people.}

Write a composition of \textbf{80 to 100 words}. Your composition should :

\begin{itemize}
\item say what the training activity was (seminar, workshop, training camp, lecture\ldots) and who organised it;
\item describe at least two activities that took place;
\item explain why such training is important for young people in Cameroon.
\end{itemize}

Use at least \textbf{five} words or expressions from the lesson, such as : \emph{to enrol, facilitator, to raise awareness, certificate of attendance, workshop, rule of law, participant, to deliver a lecture}.
\end{exercice}

\newpage

\coursec{Corrigés des exercices}

\begin{corrige}[Exercice 1 — QCM : le bon mot]
\begin{enumerate}
\item \textbf{b) attend}. \eng{I would like to attend a workshop on human rights next month.} « J'aimerais assister à un atelier sur les droits de l'homme le mois prochain. » \emph{Justification :} \eng{to attend} = assister à (un cours, une réunion). \eng{To assist} signifie « aider » (faux ami) ; \eng{assist to} et \eng{join to} sont des calques du français « assister à » et « adhérer à ».
\item \textbf{b) facilitator}. \eng{The facilitator guided the participants through the group activities.} « L'animateur a guidé les participants à travers les activités de groupe. » \emph{Justification :} \eng{animator} désigne en anglais un dessinateur de films d'animation ; l'animateur d'un atelier est un \eng{facilitator}. \eng{Trainerer} n'existe pas.
\item \textbf{b) certificate}. \eng{At the end of the course, every trainee received a certificate of participation.} « À la fin du cours, chaque stagiaire a reçu une attestation de participation. » \emph{Justification :} un \eng{diploma} sanctionne de longues études ; pour une courte formation, on délivre un \eng{certificate}.
\item \textbf{c) on}. \eng{Our school organised a training on democracy and good governance.} « Notre école a organisé une formation sur la démocratie et la bonne gouvernance. » \emph{Justification :} la préposition qui introduit le thème d'une formation, d'un cours ou d'un livre est \eng{on} (« sur »).
\item \textbf{b) awareness}. \eng{The seminar aims to raise awareness of citizens' rights among young people.} « Le séminaire vise à sensibiliser les jeunes aux droits des citoyens. » \emph{Justification :} la collocation correcte est \eng{to raise awareness of}. \eng{Conscience} est un faux ami : il signifie « conscience » (morale), non « sensibilisation ».
\end{enumerate}
\end{corrige}

\begin{corrige}[Exercice 2 — Vrai ou faux ?]
\begin{enumerate}
\item \textbf{False.} The text says it was \eng{“a one-day seminar”}, not a whole week. « Le séminaire a duré un seul jour, pas une semaine entière. »
\item \textbf{False.} \eng{“About sixty students enrolled for the training”} at the \eng{Lycée de New-Bell} : they were high school students, not university students. (The guest lecturer came from the university, but he was a speaker, not a trainee.)
\item \textbf{True.} \eng{“In the morning, a guest lecturer \dots\ delivered a talk on transparency and the fight against corruption.”}
\item \textbf{True.} \eng{“In the afternoon, the participants were divided into four workshops.”}
\item \textbf{False.} They received \eng{“certificates of attendance”}, not diplomas. « Ils ont reçu des attestations de présence, pas des diplômes. »
\end{enumerate}
\emph{Point de vigilance :} bien distinguer \eng{certificate} (attestation de courte formation) et \eng{diploma} (diplôme d'études longues).
\end{corrige}

\begin{corrige}[Exercice 3 — Vocabulaire : qui fait quoi ?]
\begin{center}
\begin{tabularx}{\linewidth}{|l|l|X|}
\hline
\rowcolor{popblueL} \textbf{Acteur} & \textbf{Réponse} & \textbf{Traduction de la définition} \\
\hline
1. facilitator & b) & une personne qui guide les discussions et activités de groupe \\
\hline
2. trainee & c) & une personne qui reçoit une formation (un stagiaire) \\
\hline
3. lecturer & d) & une personne qui fait un exposé formel devant un public \\
\hline
4. participant & a) & une personne qui s'inscrit officiellement à un cours \\
\hline
5. organiser & e) & une personne qui planifie et prépare l'événement \\
\hline
\end{tabularx}
\end{center}
\emph{Astuce :} le suffixe \eng{-ee} (\eng{trainee}, \eng{employee}) désigne celui qui \emph{subit} l'action ; le suffixe \eng{-er/-or} (\eng{trainer}, \eng{organiser}, \eng{facilitator}) désigne celui qui \emph{fait} l'action.
\end{corrige}

\begin{corrige}[Exercice 4 — Complète avec le bon mot]
\begin{enumerate}
\item \eng{Last year, over two hundred students \textbf{enrolled} for the citizenship course.} « L'année dernière, plus de deux cents élèves se sont inscrits au cours de citoyenneté. » \emph{Justification :} \eng{last year} appelle le prétérit simple ; attention au doublement du \textbf{l} : \eng{enrol} $\rightarrow$ \eng{enrolled}.
\item \eng{The trainer \textbf{delivered} an interesting lecture on voting rights yesterday.} « Le formateur a fait hier une conférence intéressante sur le droit de vote. » \emph{Justification :} \eng{yesterday} $\rightarrow$ prétérit simple.
\item \eng{Our club \textbf{will organise} (ou \textbf{is going to organise}) a workshop on the rule of law next term.} « Notre club organisera un atelier sur l'État de droit le trimestre prochain. » \emph{Justification :} \eng{next term} $\rightarrow$ futur.
\item \eng{Since January, I \textbf{have attended} three seminars on human rights.} « Depuis janvier, j'ai assisté à trois séminaires sur les droits de l'homme. » \emph{Justification :} \eng{since} + point de départ $\rightarrow$ present perfect.
\item \eng{While the facilitator was speaking, the participants \textbf{were taking} notes.} « Pendant que l'animateur parlait, les participants prenaient des notes. » \emph{Justification :} deux actions longues simultanées $\rightarrow$ prétérit continu (\eng{while}).
\end{enumerate}
\end{corrige}

\begin{corrige}[Exercice 5 — Traduction]
\begin{enumerate}
\item \eng{I would like to attend a training course on democracy during the holidays.} \emph{Attention :} « suivre une formation » se dit \eng{to attend a training course}, jamais \eng{to follow a formation}.
\item \eng{The facilitator divided us into groups of five participants.} \emph{Attention :} « diviser en » = \eng{to divide into}.
\item \eng{At the end of the seminar, every trainee received a certificate of participation.} \emph{Attention :} « attestation » = \eng{certificate}, pas \eng{attestation} (mot rare en anglais).
\item \eng{This training course aims to raise young people's awareness of human rights.} Variante : \eng{This course aims to sensitise young people to human rights.}
\item \eng{More than a hundred students enrolled for the workshop on good governance.} \emph{Attention :} « s'inscrire à » = \eng{to enrol for} (ou \eng{to register for}), jamais \eng{to inscribe}.
\end{enumerate}
\end{corrige}

\begin{corrige}[Exercice 6 — Attention aux faux amis !]
\begin{enumerate}
\item « \eng{I made a formation} » $\rightarrow$ \eng{I \textbf{attended a training course} on democracy last year.} « Faire une formation » ne se traduit ni par \eng{make} ni par \eng{formation} (qui signifie « formation » au sens géologique ou militaire).
\item « \eng{The animator} » $\rightarrow$ \eng{The \textbf{facilitator} trained us during the seminar.} Un \eng{animator} est un créateur de dessins animés.
\item « \eng{I assisted to} » $\rightarrow$ \eng{I \textbf{attended} a workshop on human rights.} \eng{To assist} = aider ; « assister à » = \eng{to attend} (sans préposition).
\item « \eng{She received a diploma} » $\rightarrow$ \eng{She received a \textbf{certificate} after a two-day training course.} Un \eng{diploma} sanctionne des études longues et officielles.
\item « \eng{The conference} » $\rightarrow$ \eng{The \textbf{lecture} was very interesting; I learned a lot.} Une \eng{conference} est un grand colloque de plusieurs jours ; un exposé donné par un intervenant est une \eng{lecture} (ou un \eng{talk}).
\end{enumerate}
\end{corrige}

\begin{corrige}[Exercice 7 — Texte à trous : a training camp]
\begin{enumerate}
\item \textbf{enrolled} : \eng{fifty students \dots\ enrolled for a civic training camp} (« se sont inscrits », prétérit, doublement du \textbf{l}).
\item \textbf{workshop} : \eng{a daily workshop on topics such as elections\dots} (« un atelier quotidien »).
\item \textbf{rule of law} : \eng{topics such as elections, tolerance and the rule of law} (« l'État de droit » ; expression figée avec l'article \eng{the}).
\item \textbf{facilitator} : \eng{each session was led by an experienced facilitator} (« un animateur expérimenté » ; \eng{an} s'accorde avec \eng{experienced}).
\item \textbf{awareness} : \eng{to raise awareness of citizens' rights and duties} (collocation \eng{to raise awareness of} = sensibiliser à).
\item \textbf{certificate} : \eng{every participant received a certificate of participation} (« une attestation de participation »).
\end{enumerate}
\end{corrige}

\begin{corrige}[Exercice 8 — Appariement : collocations]
\begin{center}
\begin{tabularx}{\linewidth}{|l|l|X|}
\hline
\rowcolor{popblueL} \textbf{Verbe} & \textbf{Nom} & \textbf{Collocation et traduction} \\
\hline
1. raise & d) & \eng{raise awareness} : sensibiliser \\
\hline
2. attend & a) & \eng{attend a workshop} : assister à un atelier \\
\hline
3. deliver & e) & \eng{deliver a lecture} : faire une conférence \\
\hline
4. fight & b) & \eng{fight corruption} : lutter contre la corruption \\
\hline
5. promote & c) & \eng{promote democracy} : promouvoir la démocratie \\
\hline
\end{tabularx}
\end{center}
\textbf{Exemples de phrases personnelles (modèles) :}
\begin{itemize}
\item \eng{The association organised a campaign to raise awareness of voting rights in Buea.} « L'association a organisé une campagne pour sensibiliser aux droits de vote à Buea. »
\item \eng{Journalists and judges must work together to fight corruption.} « Les journalistes et les juges doivent travailler ensemble pour lutter contre la corruption. »
\item \eng{School clubs play an important role in promoting democracy among young people.} « Les clubs scolaires jouent un rôle important dans la promotion de la démocratie chez les jeunes. »
\end{itemize}
\end{corrige}

\begin{corrige}[Exercice 9 — Compréhension et langue — type Bac]
\textbf{Barème indicatif (sur 20) :} A. Compréhension : 10 pts (2 pts par question) ; B. Langue : 10 pts (Q1 : 3 pts ; Q2 : 3 pts ; Q3 : 4 pts).

\textbf{A. Comprehension questions — réponses modèles :}
\begin{enumerate}
\item \eng{The bilingual high school of Yaoundé organises the citizenship week for its final-year students.} (Le programme est préparé par les enseignants avec une association locale.)
\item \eng{During the week, students attend lectures and debates, and they take part in practical workshops.} (Deux activités suffisent : \eng{lectures}, \eng{debates}, \eng{workshops}, \eng{a mock election}, \eng{a closing ceremony}.)
\item \eng{On Thursday, the students hold a mock election: they elect their class representatives under the supervision of their teachers.}
\item \eng{According to the head teacher, democracy also means listening to others, accepting differences and taking responsibility.}
\item \eng{Before, André thought elections were “just about putting a paper in a box”; now he knows that “being a citizen means much more.”}
\end{enumerate}

\textbf{B. Language work :}
\begin{enumerate}
\item (a) \eng{facilitator} (« a person who guides a workshop ») ; (b) \eng{certificate (of attendance)} ; (c) \eng{a mock election} (« une élection simulée »).
\item (a) \eng{The students attended a workshop \textbf{on} the rule of law.} (b) \eng{Democracy is not only \textbf{about} voting.} (c) \eng{The election took place \textbf{under} the supervision of the teachers.}
\item \eng{Last year, the school \textbf{organised} its first citizenship week, and since then, many students \textbf{have joined} the civic club.} \emph{Justification :} \eng{last year} $\rightarrow$ prétérit simple ; \eng{since then} $\rightarrow$ present perfect.
\end{enumerate}

\textbf{Points de vigilance :} répondre avec ses propres mots (pas de recopie intégrale) ; ne pas confondre \eng{mock election} (simulation) et \eng{real election} ; soigner l'accord des temps.
\end{corrige}

\begin{corrige}[Exercice 10 — Traduction — type Bac]
\textbf{Barème indicatif (sur 10) :} passage anglais $\rightarrow$ français : 5 pts ; deux phrases français $\rightarrow$ anglais : 2,5 pts chacune.

\textbf{1. Traduction en français :}

« Les jeunes doivent comprendre que la démocratie ne consiste pas seulement à voter, explique-t-elle. Elle consiste à écouter les autres, à accepter les différences et à prendre ses responsabilités. »

\emph{Points de vigilance :} \eng{it is about + V-ing} = « il s'agit de / cela consiste à » ; \eng{taking responsibility} = « prendre ses responsabilités » (et non « prendre la responsabilité » au singulier).

\textbf{2. Traduction en anglais :}
\begin{enumerate}
\item \eng{Every year, our school organises a citizenship week for final-year students.} \emph{Attention :} « chaque année » = \eng{every year} (présent simple, \eng{organises} avec \textbf{s}) ; « les élèves de Terminale » = \eng{final-year students}.
\item \eng{The participants learned to debate respectfully and to accept differences.} \emph{Attention :} « avec respect » se rend par l'adverbe \eng{respectfully} ; \eng{learn} est irrégulier : \eng{learn — learned/learnt — learned/learnt}.
\end{enumerate}
\end{corrige}

\begin{corrige}[Exercice 11 — Composition — type Bac]
\textbf{Barème indicatif (sur 20) :} respect de la consigne et du nombre de mots (80--100 mots) : 4 pts ; contenu (les trois points demandés) : 6 pts ; réemploi d'au moins cinq mots de la leçon : 4 pts ; correction grammaticale et orthographique : 4 pts ; organisation et cohérence : 2 pts.

\textbf{Proposition de rédaction modèle (environ 95 mots) :}

\begin{textebox}[Model composition]
Last term, our school organised a two-day workshop on citizenship and the rule of law. About eighty students enrolled for the training. On the first day, a guest lecturer delivered a talk on human rights, and then a trained facilitator divided us into groups to debate real-life situations. On the second day, we held a mock election to learn how to vote transparently. At the end, every participant received a certificate of attendance. This training is important because it raises young people's awareness of their rights and duties and helps them become responsible citizens in Cameroon.
\end{textebox}

\textbf{Points de vigilance :}
\begin{itemize}
\item Utiliser le prétérit simple pour le récit (\eng{organised}, \eng{enrolled}, \eng{delivered}, \eng{received}) et le présent pour l'importance générale (\eng{it raises}, \eng{helps}).
\item Employer les collocations exactes : \eng{deliver a talk/lecture}, \eng{raise awareness}, \eng{certificate of attendance}, \eng{enrol for}.
\item Éviter les faux amis : pas de \eng{formation}, pas de \eng{animator}, pas de \eng{assist to}.
\item Vérifier le nombre de mots : en dessous de 80 ou au-dessus de 100, des points sont perdus.
\end{itemize}
\end{corrige}

\newpage

\coursec{Fiche bilan}

\begin{popfigure}
\begin{tikzpicture}[mindmap, grow cyclic, every node/.style={concept, font=\small}, concept color=popblue, root concept/.append style={concept color=popblue, font=\small\bfseries}, level 1/.append style={level distance=3.4cm, sibling angle=51}, level 2/.append style={level distance=2.1cm, sibling angle=45, font=\scriptsize}]
\node[root concept]{Training on democracy}
  child[concept color=poporange]{ node {Types of courses} 
    child { node {workshop} }
    child { node {seminar} }
    child { node {refresher course} }
  }
  child[concept color=popgreen]{ node {Actors}
    child { node {trainer / facilitator} }
    child { node {trainee / participant} }
  }
  child[concept color=poppurple]{ node {Democratic themes}
    child { node {elections} }
    child { node {rule of law} }
    child { node {human rights} }
  }
  child[concept color=poppink]{ node {False friends}
    child { node {formation $\neq$ training} }
    child { node {stage $\neq$ stage} }
  }
  child[concept color=popgold]{ node {Word families}
    child { node {train / trainer / trainee} }
    child { node {elect / election / elective} }
  }
  child[concept color=popmuted]{ node {Collocations}
    child { node {attend a course} }
    child { node {raise awareness} }
  }
  child[concept color=popblueL]{ node {In context}
    child { node {dialogue} }
    child { node {guided writing} }
  };
\end{tikzpicture}
\legende{Carte mentale de la leçon : le vocabulaire des formations sur la démocratie}
\end{popfigure}

\begin{bilan}
\textbf{1. Lexique clé (English -- Français)}
\begin{itemize}
  \item \eng{a training course / a training session} : une formation / une session de formation
  \item \eng{a workshop} : un atelier ; \eng{a seminar} : un séminaire ; \eng{a refresher course} : un cours de remise à niveau
  \item \eng{a trainer / a facilitator} : un formateur / un animateur ; \eng{a trainee / a participant} : un stagiaire / un participant
  \item \eng{civic education} : l'éducation civique ; \eng{the rule of law} : l'État de droit
  \item \eng{free and fair elections} : des élections libres et transparentes ; \eng{a polling station} : un bureau de vote
  \item \eng{human rights} : les droits de l'homme ; \eng{good governance} : la bonne gouvernance
  \item \eng{to raise awareness} : sensibiliser ; \eng{to attend a course} : suivre une formation ; \eng{to hand out certificates} : remettre des attestations
\end{itemize}

\textbf{2. Points de langue à connaître par cœur}
\begin{center}
\begin{tabularx}{\linewidth}{|X|X|}\hline
\rowcolor{popblueL} \textbf{Règle} & \textbf{Exemple} \\ \hline
\cle{train} + \cle{-er} (celui qui forme) / \cle{-ee} (celui qui est formé) & \eng{The trainer congratulated the trainees.} \\ \hline
Collocations : \eng{to attend / to follow a course} (jamais \emph{to assist a course}) & \eng{She attended a workshop in Yaoundé.} \\ \hline
\eng{to be interested in + -ing} & \eng{They are interested in learning about elections.} \\ \hline
Faux amis : \eng{training} = formation ; \eng{an internship} = un stage ; \eng{a lecture} = une conférence & \eng{He did an internship at the town hall.} \\ \hline
\end{tabularx}
\end{center}

\textbf{3. Méthode express}
\begin{enumerate}
  \item Repérer le thème du texte (formation, démocratie, citoyenneté).
  \item Souligner les collocations et les mots de la même famille.
  \item Traduire en évitant les calques du français (\emph{formation}, \emph{assister}, \emph{stage}).
  \item Réemployer le lexique dans une phrase personnelle à l'oral ou à l'écrit.
\end{enumerate}
\end{bilan}

\begin{aretenir}[Checklist avant le Bac]
\begin{itemize}
  \item[$\square$] Je connais les types de formations : \eng{workshop, seminar, refresher course, training session}.
  \item[$\square$] Je distingue \eng{trainer} (formateur) et \eng{trainee} (stagiaire).
  \item[$\square$] Je sais dire « suivre une formation » : \eng{to attend a training course}.
  \item[$\square$] Je ne traduis pas « formation » par \emph{formation} mais par \eng{training}.
  \item[$\square$] Je ne traduis pas « stage » par \emph{stage} mais par \eng{internship}.
  \item[$\square$] Je connais les familles de mots : \eng{elect / election / elective} ; \eng{govern / government / governance}.
  \item[$\square$] Je maîtrise les collocations : \eng{to raise awareness, to hold elections, to hand out certificates}.
  \item[$\square$] Je connais le lexique démocratique : \eng{rule of law, polling station, human rights, good governance}.
  \item[$\square$] Je sais utiliser \eng{to be interested in + -ing} sans faute.
  \item[$\square$] Je peux rédiger un court paragraphe sur une formation citoyenne en réemployant au moins cinq mots de la leçon.
\end{itemize}
\end{aretenir}

\findecours

\end{document}
